\documentclass[11pt,a4paper]{article}
\usepackage[T1]{fontenc}
\usepackage[utf8]{inputenc}
\usepackage{lmodern,microtype}
\usepackage[margin=25mm,headheight=14pt]{geometry}
\usepackage{amsmath,amssymb,amsthm,mathtools,mathrsfs}
\usepackage{booktabs,longtable,array,enumitem}
\usepackage{xcolor,hyperref,fancyhdr,fancyvrb,xurl,needspace}
\hypersetup{colorlinks=true,linkcolor=blue!45!black,citecolor=blue!45!black,urlcolor=blue!45!black,pdftitle={Frohlich weak-coupling spectral theorem: mathematical blueprint},pdfauthor={Research formalization working document}}
\pagestyle{fancy}\fancyhf{}\fancyhead[L]{\small Fr\"ohlich weak coupling: proof blueprint}\fancyhead[R]{\small v5}\fancyfoot[C]{\thepage}
\setlength{\parindent}{0pt}\setlength{\parskip}{5pt}
\setlist{nosep,leftmargin=1.8em}
\allowdisplaybreaks[2]
\setlength{\emergencystretch}{2em}
\newtheorem{theorem}{Theorem}[section]
\newtheorem{lemma}[theorem]{Lemma}
\newtheorem{proposition}[theorem]{Proposition}
\newtheorem{definition}[theorem]{Definition}
\newtheorem{corollary}[theorem]{Corollary}
\theoremstyle{remark}\newtheorem{remark}[theorem]{Remark}
\newcommand{\R}{\mathbb R}\newcommand{\C}{\mathbb C}\newcommand{\N}{\mathbb N}
\newcommand{\F}{\mathcal F}\newcommand{\Q}{\mathcal Q}\newcommand{\HH}{\mathcal H}
\newcommand{\Pf}{P_{\mathrm f}}\newcommand{\Id}{I}\newcommand{\Dom}{\operatorname{Dom}}
\newcommand{\Ran}{\operatorname{Ran}}\newcommand{\esssup}{\operatorname*{ess\,sup}}
\newcommand{\ip}[2]{\langle #1,#2\rangle}\newcommand{\norm}[1]{\lVert #1\rVert}
\newcommand{\abs}[1]{\lvert #1\rvert}\newcommand{\rankone}[2]{\lvert #1\rangle\langle #2\rvert}
\newcommand{\dd}{\,\mathrm d}\newcommand{\almosteverywhere}{\text{a.e.}}
\newcommand{\node}[3]{\par\Needspace{6\baselineskip}\smallskip\hypertarget{node:#1}{}\label{nodepage:#1}\noindent\textbf{Node \texttt{#1}.} \textit{#2.}\par\noindent\textbf{Direct dependencies:} #3.\par}
\newcommand{\contract}[1]{\textsf{#1}}
\newcommand{\module}[1]{\texttt{\path{#1}}}
\newcommand{\Ecal}{\mathcal E}
\newcommand{\sigR}{\sigma_{\R}}
\newcommand{\resR}{\rho_{\R}}
\newcommand{\weak}{\mathbin{\rightharpoonup}}
\title{The weak-coupling Fr\"ohlich spectral theorem\\\large A dependency-ordered mathematical blueprint for Lean\\\normalsize with fixed GW objects, explicit adapters, and limit/domain contracts}
\author{Mathematical specification: mathlib plus locally retained GW proofs}
\date{13 November 2026 \quad Version 5}
\begin{document}
\maketitle
\begin{abstract}
This document gives a dependency-ordered natural-language proof of the weak-coupling spectral equality for the three-dimensional Fr\"ohlich fibers with $|P|\leq1/4$. Its local lemmas include the measure-space and graph-domain constructions, all smeared-field commutations used in the argument, the associated self-adjoint operator, energy-space Schur inversion, cutoff pull-through, the integrated kernel, the unnormalized Seiringer cancellation, and the weighted comparison estimates. The conclusion concerns the full spectrum, not only eigenvalues. The proof does not use an essential-spectrum theorem, a global minimum theorem for the dispersion, positivity improvement, a separate simplicity theorem, or the all-momentum extension.

This revision retains the previously formalized GW objects as canonical, supplies explicit adapters for their domains and weighted coordinates, and expands the norm and source contracts at the interfaces of the analytic argument. In particular, the scalar-to-Bochner $L^2$ identification is the explicitly restricted inverse of GW's existing product unitary, not an assumed new identification; the finite-cutoff estimates precede the singular form limit; and pull-through is extracted from a single Bochner $L^2$ residual rather than from pointwise unbounded operators or a parameter-dependent family of null sets. Precise mathematical-library leaves and source-paper locations are recorded separately from the local arguments. Version 5 implements the targeted reuse and dependency audit: it removes an optional circular proof route, makes coercive inversion a direct GW reuse, specifies noncommutative inverse/calculus interfaces, and supplies staged three-agent interface gates without changing the target theorem.
\end{abstract}
\paragraph{Status and reading contract.}
The mathematical proof is presented in full modulo the precisely indexed mathlib and retained GW results. The new mathematical constructions and their proofs are part of this document; they are not delegated to a Lean implementer. Section~\ref{sec:adapters} fixes their actual GW realizations and every domain/norm conversion that is used. The external Lake package is mathlib alone; the supplied GW proof source is retained locally unchanged. The retained source identity is the unchanged ZIP shipped under \path{inputs/GW-Froehlich-GitHub-release.zip}, Lean \texttt{v4.36.0-rc2}, and mathlib commit \texttt{9bac2d0692944a775113f7c130fd5e5b4608a854}. Exact GW names, source ranges, and hashes are recorded in the source index and companion manifest.

This remains a natural-language mathematical manuscript, not a kernel-checked certificate. The source references have been inspected, but neither the retained GW release nor a translation of this manuscript has been built here. The distinction is between a supplied proof and a mechanically verified proof, not permission to invent additional analytic arguments. The mathlib reference index distinguishes inherited GW occurrences, older documentation-only lookup aids, and the selected pinned-source interfaces rechecked for version 5 in Section~\ref{sec:pinnedcontracts}. None of these is an elaboration test of its new fiber specialization. The companion metadata records evidence levels separately from proof status; it supplies no pre-verified new nodes. The proof uses ordinary Hilbert-space and measure-theory statements with the mathematical hypotheses stated at their applications.

\clearpage
\tableofcontents
\newpage

\section{Target, conventions, and the corrected architecture}\label{sec:target}
\paragraph{How to read a dependency node.}
Each node gives fixed objects, a statement, its direct mathematical prerequisites, and a proof. References such as ML-L2 are the explicitly listed library leaves, not unnamed analytic assumptions. A reference to a source paper records provenance; it never replaces a proof of a modified argument. Definitions of measurable kernels always specify exceptional-point conventions, and equalities in a form dual or an $L^p$ space retain that meaning until a later lemma proves a stronger assertion.


The starting papers are Mitrouskas~\cite{mitrouskas} and Seiringer~\cite{seiringer}, on top of the existing Griesemer--W\"unsch formalization~\cite{gw}. Section~\ref{sec:concordance} identifies precisely which arguments are retained and which are replaced.
\begin{theorem}[Target theorem]\label{thm:target}
There is $g_*>0$ such that, for every $g\in\R$ and $P\in\R^3$,
\[
 \abs g<g_*,\qquad \abs P\leq\frac14
 \quad\Longrightarrow\quad
 \sigma(H_g(P))\cap(-\infty,E_g(0)+1)=\{E_g(P)\}.
\]
Here
\[
 H_g(P)=(\Pf-P)^2+N+\Phi(v_g),\qquad v_g(k)=g/\abs k,
 \qquad E_g(P)=\inf\sigR(H_g(P)),
\]
and $H_g(P)$ is the operator associated with the closed interaction form. Moreover, $E_g(P)$ has a nonzero eigenvector.
\end{theorem}
The real interval in the displayed statement is identified with its image in $\C$ if the existing spectrum API is complex-valued. The theorem proved first below uses $\sigR$; its bounded-inverse definition agrees with the complex spectrum restricted to this real interval. No statement about nonreal spectral parameters is needed for that identification. Every coupling-dependent object, including the symbol $e=E_g(0)$, retains the \emph{same} fixed $g$ throughout a proof.

\paragraph{Exact release scope.}
The sole designated main target is Theorem~\ref{thm:target}, assembled at MN03. Its quantifiers include signed and zero coupling, zero momentum, and $|P|=1/4$, with a strictly open spectral interval at $E_g(0)+1$. Ground-state simplicity and the all-momentum corollary are not targets of this release. Nor is $E_g(0)+1$ identified here with the essential-spectrum threshold. These explicit scope choices do not authorize weakening any clause of the displayed theorem.

\paragraph{Conventions.}
All Hilbert spaces of wavefunctions are complex; $\ip{x}{y}$ is conjugate-linear in $x$ and linear in $y$. Real kernels are also used as complex integral kernels. Parameter differentiation is over $\R$. Integrals use Lebesgue measure with no $(2\pi)^{-3}$ factor. The Euclidean norm is denoted $\abs\cdot$. Values of singular measurable functions at their finitely many singular points are set to zero, unless a different explicit value is specified. This never changes an $L^p$ class. The interaction parameter is $g$ in $v_g=g/\abs k$, not the conventional physical parameter proportional to $g^2$.

\paragraph{What changes from the initial dependency sketch.}
We use a \emph{directed acyclic graph}, not a literal tree: the form solver, kernel bounds, and singular-integral estimates are shared. We replace the normalized rank-one reduction by an elementary completion of the square. We also replace the claim ``a positive Schur complement has no zero eigenvector'' by a construction of a bounded inverse for $H_g(P)-\lambda$ on its actual domain. These changes remove both a spectral gap in the argument and an unnecessary perturbative calculation.

The essential interfaces at the top of the graph are
\[
\begin{array}{c}
\text{GW Fock/form foundation + mathlib Hilbert and measure theory}\\
\downarrow\\
\text{uniform form bounds and an energy-space inverse lemma}\\
\swarrow\hspace{4cm}\searrow\\[-2mm]
\text{vacuum root / ground vector}\qquad\text{one-phonon Schur complement}\\
\downarrow\hspace{4cm}\downarrow\\[-2mm]
\text{weighted resolvents and threshold kernel}\\
\swarrow\hspace{4cm}\searrow\\[-2mm]
\text{unnormalized Seiringer identity}\qquad\text{weighted comparison estimates}\\
\searrow\hspace{4cm}\swarrow\\[-2mm]
\text{threshold positivity at all }\abs P\leq1/4\\
\downarrow\\
\text{resolvent exclusion + vacuum root = target theorem.}
\end{array}
\]
This display is a dependency schematic, not a new analytic assertion. A precise module/lemma registry appears in Section~\ref{sec:registry}.

\section{Local GW reuse and the fixed mathematical boundary}\label{sec:gw}
The existing GW development is valuable primarily for its Fock-space and form foundations, not for the final dressed-domain theorem. The required paper-level source is Griesemer--W\"unsch, Appendix B, especially the sector formulas (38)--(39) and Lemma B.1; its abstract form argument in Theorem A.1 is also relevant. The full-space statements of its Lemma 2.1 and Theorem 2.2 are not silently applied to fibers. Uniform fiber bounds and the stronger energy-dual convergence are proved locally below.\cite{gw}

\subsection{Results to retain, with fixed canonical implementations}\label{sec:localreuse}
\begin{longtable}{>{\raggedright\arraybackslash}p{.24\textwidth}>{\raggedright\arraybackslash}p{.32\textwidth}>{\raggedright\arraybackslash}p{.35\textwidth}}
\toprule
GW reuse target & Exact mathematical content & Local construction/proof in this document\\\midrule\endhead
Fock realization & Symmetric $L^2$ sectors, Hilbert sum, vacuum, sector maps and normalization. & FD01, MS01--MS03.\\
Multiplication and graph domains & Number, total momentum, free form and operator domains, completeness, dense core. & FD02, FC01.\\
Smeared fields & Sector formulas, adjointness, number estimates, all commutations used here. & FC02--FC04, FC06--FC07.\\
Form representation & Dense self-adjoint associated operator with exact domain characterization. & BA01--BA03, FA00--FA02.\\
Interaction form limit & Cutoff inequalities and mixed tail convergence, with no domain guess. & FB01, FB02A, FB02.\\
Operator/spectrum compatibility & Equality of domain-and-action data; bounded two-sided inverse predicate. & FA05, RE01.\\
\bottomrule
\end{longtable}
These are not theorem assumptions in the final result. After the source audit in Section~\ref{sec:gwverify}, the GW Fock carrier, smeared fields, product-unitary, number operator, and abstract form-representation objects are \emph{mandatory reuse points}. Their corresponding local proof paragraphs below explain the mathematics, but \emph{must not} be reimplemented where the mapped GW declarations apply. Fiber momentum, its graph core, the singular fiber form bounds, and the unsmeared annihilation field remain new. RE01 supplies a transport criterion only for any genuine mismatch; it does not justify creating a second Fock space.

\paragraph{Source-extraction record.}
The source audit is performed at the declaration-text and local-import-graph level; see Section~\ref{sec:gwverify} and the expanded companion machine-readable source manifest. The uploaded snapshot, rather than a remotely fetched GitHub revision, is the audited source identity. The exact public Git commit associated with this ZIP has not been established. Copy the transitive GW proof bodies, not just their statements, and initially preserve all 115 local modules unchanged. They import only mathlib and local GW files. Retain attribution and check the permission status before public redistribution. An actual Lean compilation of this archive remains an independent, pending check. The fiber-specific mathematics, and not the full-space Gross-transform/Hamiltonian theorem, is to be newly formalized.

\subsection{Complex conjugation, duals, and the load of a singular source}\label{sec:dualconventions}
Use inner products conjugate-linear in the first variable. Write $V^{\times}$ for continuous conjugate-linear functionals on a complex Hilbert energy space $V$. A load $F\in V^{\times}$ has its unique Riesz vector $z_F\in V$ characterized by
\[
 F(w)=\ip w{z_F}_V.
\]
The map $F\mapsto z_F$ is complex-linear and isometric. This notation need not become a new general duality library: the implementation may store the Riesz vector directly. The ordinary continuous linear dual instead uses $w\mapsto\ip{z}{w}$ and the conjugate-linear Riesz equivalence; the conversion is complex conjugation of the functional. These two conventions must not be interchanged silently.

For the real kernel $v_g$, the scalar functional
$\ell_g(f)=\int v_g(k)f(k)\dd k$ is linear. The load on the right side of a weak Hamiltonian equation is
\[
 \mathfrak v_g(w)=\overline{\ell_g(w)}
           =\int\overline{w(k)}v_g(k)\dd k,
\]
not $\ell_g(w)$. Thus the notation $h=Rv_g$ means
\[
 a(w,h)=\mathfrak v_g(w)\quad(w\in V).
\]
The scalar $\ip{v_g}{Rv_g}$ means $\ell_g(h)$, which is real and nonnegative for a positive inverse. In Riesz coordinates it is $\ip{z_{\mathfrak v_g}}{T^{-1}z_{\mathfrak v_g}}_V$. This convention also fixes the sign and conjugation of the vacuum block.

For a bounded form $a(w,u)$, its representative $T$ is specified by
$\ip w{Tu}_V=a(w,u)$. If $j:V\to\HH$ is the embedding, the physical load $f\in\HH$ has Riesz vector $j^*f$. If $W:V_1\to V_2^{\times}$ is a coupling, its bounded representative $B:V_1\to V_2$ is characterized by
\[
 \ip{w_2}{Bu_1}_{V_2}=W(u_1)(w_2).
\]
The reverse block is $B^*$, and a Schur correction is $B^*D^{-1}B$ on energy spaces. In a notation such as $A^*XA$, each occurrence of $X$ on a form-dual load means this energy inverse, not merely the bounded physical inverse $jD^{-1}j^*$.

\paragraph{Changing Hilbert structures.}
The form-domain norm is fixed locally before invoking Riesz representation. A different equivalent graph norm changes its Riesz representatives, but not the weak solution or the associated operator. Indeed both solve the same equalities $a(w,u)=F(w)$, and uniqueness follows from coercivity. Every result stated with an energy inverse is therefore invariant under the norm changes used to compare momenta. This observation is a proof, not a type-level identification of different inner products.

\subsection{Concrete form domains and projections}
For $P\in\R^3$ put
\begin{align}
 h^{\mathrm{fr}}_P[u]&=\sum_{n=0}^{\infty}\int\bigl(\abs{K_n-P}^2+n\bigr)\abs{u_n}^2,\\
 \norm u_{\Q_P}^2&=\norm u^2+h^{\mathrm{fr}}_P[u].\label{eq:qnorm}
\end{align}
For $n=0$, $K_0=0$ and the integral is the vacuum scalar. The set underlying $\Q_P$ is independent of $P$, since
\[
 \abs{K-P}^2\leq2\abs K^2+2\abs P^2,
 \qquad \abs K^2\leq2\abs{K-P}^2+2\abs P^2.
\]
We denote it by $\Q$. For $\abs P\leq1/4$ these norms are uniformly equivalent. Sector restrictions are closed subspaces for the graph norm because the projections commute with the weights. Write $\Q_{\geq n}=\Pi^{\geq n}\Q$.

For the isolated one-particle energy space we use
$\|f\|_{\Q_{1,P}}^2=\int(1+|k-P|^2)|f(k)|^2dk$.
Its underlying space is the one-particle restriction of $\Q_P$; the latter has weight $2+|k-P|^2$. The two squared norms lie between one and two times each other. This equivalence is used only to transfer bounds and completeness, never to identify their Riesz representatives. In particular $\Q_1$ without a momentum subscript denotes the isolated norm at $P=0$.

The source $v_g$ is a continuous functional on $\Q_1$, not a vector of $L^2$ when $g\ne0$. Write
\begin{equation}
 \ell_g(f)=\int_{\R^3}v_g(k)f(k)\dd k,
 \qquad
 \abs{\ell_g(f)}\leq
 \left(\int\frac{g^2}{\abs k^2(1+\abs{k-P}^2)}\dd k\right)^{1/2}
 \norm f_{\Q_{1,P}}.\label{eq:sourcebound}
\end{equation}
The integral is uniformly bounded by $C g^2$ by Lemma~\ref{lem:shiftedsource}. Thus the weak load represented by $v_g$ is $\mathfrak v_g(w)=\overline{\ell_g(w)}$, as specified in Section~\ref{sec:dualconventions}; it is not the linear functional $\ell_g$ itself. The notation $Rv_g$ always means the form-domain solution to that equation. The pairing $\ip{v_g}{Rv_g}$ means $\ell_g(Rv_g)$.

\paragraph{Core and symmetry conventions.}
The core $\mathscr C$ consists of finite sequences of symmetric $L^2$ functions with bounded support in each momentum coordinate, and includes $\Omega$; FC01 proves its density and cutoff-field invariance. Smoothness is not required. Complex conjugation $J$ acts componentwise. Momentum reflection $\mathcal R$ sends $u_n(k_1,\ldots,k_n)$ to $u_n(-k_1,\ldots,-k_n)$. It is enough to prove invariance under this single reflection; a formal $O(3)$ representation is unnecessary. Both operations preserve the core; conjugation preserves each graph norm, while reflection identifies the graph norm at $P$ with that at $-P$. Reflection changes $\Pf$ to $-\Pf$, fixes $N$, and fixes the even kernel $v_g$.

\section{The exact GW objects and their mathematical adapters}\label{sec:adapters}
This section is normative: the symbols in the later analytic proof denote the objects specified here. A proof later in the manuscript explaining a Hilbert-space construction is not a request to create a competing definition. The retained source is part of the input, with its original module names. The source index in Section~\ref{sec:gwverify} and its JSON companion give the exact declaration and file for every GW name used here.

\node{GW01}{Canonical sectors, scalar vacuum, one-particle inclusion and projections}{GW-Sectors, GW-Coordinates, GW-Generalized}
\begin{lemma}\label{lem:gwsectors}
Take $X=\texttt{GW.Space}$, $L^2(X)=\texttt{GW.OnePhonon}$, $H_n=\texttt{GW.Fock.Sector}\ n$, and $\F=\texttt{GW.Fock.Carrier}$. For $n\geq1$ the physical representative of an element of $H_n$ is its inclusion into \texttt{AmbientSector} $n$; for $n=0$ it is its image under \texttt{vacuumAdapter}. These maps are isometries. Let
\begin{align*}
 j_1&=(\texttt{ordinarySectorKernel}\ 1)^{-1}\circ
       \texttt{kernelOneBoson},\\
 \iota_1 f&=\operatorname{single}_1(j_1f),\qquad
 \iota_0 z=\operatorname{single}_0z.
\end{align*}
Here \texttt{kernelOneBoson} is specialized to coefficient Hilbert space $\C$. Then $j_1:L^2(X)\simeq H_1$ is unitary, $\iota_0 1$ is the existing GW vacuum, and $\iota_1$ is an isometry onto the one-particle sector. Coordinate extraction followed by $j_1^{-1}$ is its adjoint. The orthogonal projections $\Pi^n$, $\Pi^{\leq m}$, and $\Pi^{\geq m}$ are bounded contractions on $\F$.
\end{lemma}
\begin{proof}
GW defines $H_0=\C$ and $H_{n+1}$ as the closed permutation-fixed subspace of $L^2(X^{n+1},\mu_{n+1})$, with $\mu_n=\texttt{sectorMeasure}\ n$. It defines $\F$ as the existing $\ell^2$ Hilbert sum of these sectors. The empty product measure is a unit point mass, and \texttt{vacuumAdapter} is the isometry sending $z$ to the constant function $z$ on that singleton. For positive sectors \texttt{sectorInclusion} is the isometric subtype inclusion. Thus the descriptions do not replace the scalar vacuum by an inequivalent measure-space convention.

Both factors in $j_1$ are existing linear isometric equivalences. The representative of $\texttt{kernelOneBoson}\ f$ at a one-coordinate configuration $k$ is $f(k_0)$: this follows from \path{kernelOneBoson_coe}, the definition of the zero-sector split, and \path{oneBosonCoordinate_preserving}. The formula therefore identifies exactly the one-variable functions used throughout the proof. There is no normalization factor.

The GW sum norm identity is $\|u\|^2=\sum_n\|u_n\|^2$. Consequently the map which keeps just coordinate $n$ is contractive, its square is itself, and the sum inner-product formula shows it is self-adjoint. The finite sum of these maps for $n\leq m$ is GW's \texttt{P} $m$. Define $\Pi^{\geq0}=I$ and $\Pi^{\geq m}=I-\texttt{P}(m-1)$ for $m\geq1$. Coordinatewise these retain precisely the indicated coordinates, so the same norm and inner-product arguments prove their contraction and orthogonality assertions. Pairing $\iota_1f$ with $u$ leaves only $\langle j_1f,u_1\rangle=\langle f,j_1^{-1}u_1\rangle$, proving the adjoint formula. Similarly $\iota_0^*u=u_0$. These are equalities on the existing GW carrier, not a transport to another Fock space.
\end{proof}

\node{GW02}{GW currying in the correct direction, including symmetric sections}{GW01, GW-Product, GW-Generalized, GW-Coordinates, ML-LP, ML-PROD}
\begin{lemma}\label{lem:gwcurrying}
There is a specified linear isometry
\[
 S_n:H_{n+1}\longrightarrow L^2(X;H_n)
\]
whose representative is the slice $S_nu(k)=u(k,\cdot)$ almost everywhere. It is the restriction of the following composition, not a newly constructed product unitary:
\begin{align*}
 H_{n+1}&\hookrightarrow\operatorname{AmbientSector}(n+1)
 \xrightarrow{\ \texttt{splitUnitary}\ n\ 0\ }L^2(X\times X^n)\\
 &\xrightarrow{\ (\texttt{productUnitary}\ \mathrm{vol}\ \mu_n)^{-1}\ }
 L^2(X;\operatorname{AmbientSector}n),
\end{align*}
followed by the pointwise symmetric-subspace identification and, at $n=0$, the inverse vacuum adapter. Thus \emph{the inverse} of GW's product unitary curries a scalar product-space function.
\end{lemma}
\begin{proof}
\texttt{splitCoordinates\_preserving} and \texttt{splitUnitary\_ae} give exactly the coordinate change which distinguishes the first variable. No integration factor appears. GW's \texttt{productUnitary} maps a Bochner $L^2$ space to the product $L^2$ space, as its declaration explicitly states. Its inverse maps rectangle tensors to their curried fields, by \texttt{productUnitary\_rectangle} and injectivity. The almost-everywhere section statement for this inverse follows from the following norm argument, which also specifies it for arbitrary representatives.

Choose rectangle-simple $u_j$ with $\|u_j-u\|_2\leq2^{-2j}$. Fubini gives
$\int\|u_j(k,\cdot)-u(k,\cdot)\|_2^2\,dk=\|u_j-u\|_2^2$.
Tonelli applied to the sum shows the sum of these section errors is finite for almost every $k$, hence the errors tend to zero there. The same estimate for the images under the isometry, whose errors have equal $L^2$ norms, shows their pointwise Hilbert-norm errors tend to zero outside another null set. On rectangle-simple functions the sections and the images agree. Taking the two limits identifies the image with the section for almost every $k$. Fubini also ensures that $u(k,\cdot)$ is square integrable there. Changing a representative changes sections only outside a null set of $k$, by the product almost-everywhere theorem, and therefore does not change the Bochner class. Rectangle density can be reused from \texttt{rectangle\_dense\_criterion}; MS01--MS02 below give the full mathematical density argument for reference.

For every permutation $\pi$ of the remaining $n$ coordinates, the input $u$ is fixed by its extension that leaves the distinguished coordinate fixed. On rectangle tensors, currying intertwines this operation with pointwise application of \texttt{kernelPermutation}\ $n\ \pi$. Equality for all inputs follows by continuity and density. The resulting Bochner vector is therefore fixed by each such pointwise permutation. There are finitely many of them; discard the finite union of their exceptional null sets. The resulting slices lie in the closed fixed subspace almost everywhere.

For a representative-free construction of this last step, apply the existing bounded map \texttt{kernelFixedProjection} pointwise, then $(\texttt{ordinarySectorKernel}\ n)^{-1}$. This defines a strongly measurable sector-valued field by \texttt{amplify}. The projection changes no slice almost everywhere, by \texttt{kernelSymmetrizer\_of\_fixed}; the inverse is isometric. Consequently the resulting map is an isometry, not merely a contraction. At $n=0$ this same construction is exactly the inverse scalar-vacuum identification. The displayed slice formula and norm identity follow. No representative is evaluated at a prescribed momentum.
\end{proof}

\node{GW03}{Number graph coordinates and the assembled annihilation field}{GW02, GW-Number, GW-Bochner, ML-SERIES}
\begin{lemma}\label{lem:gwnumberfield}
The complete number-form space $V_N$ may be implemented as a named copy of $\F$ with the physical embedding
\[
 (j_Nx)_n=(1+n)^{-1/2}x_n.
\]
Its range is precisely the domain of the existing \texttt{numberPower Sector $(1/2)$}, and
\[
 \|x\|^2=\|j_Nx\|^2+\|N^{1/2}j_Nx\|^2.
\]
On this space the annihilation field is the already-specified composition
\[
 x\longmapsto\bigl(\sqrt{n+1}\,S_n(j_Nx)_{n+1}\bigr)_{n\geq0}
 \xrightarrow{\ \texttt{bochnerUnitary}\ \mathrm{vol}\ \texttt{Sector}\ }
 L^2(X;\F).
\]
In physical notation it satisfies all the assertions of FC02, including
$\|\mathcal Au\|^2=\|N^{1/2}u\|^2$.
\end{lemma}
\begin{proof}
GW's \texttt{numberPower\_apply}, \texttt{rootNumberCoefficient}, and \texttt{numberPower\_norm\_sq} identify $N^{1/2}$ coordinatewise with $\sqrt n$. The domain condition is $\sum_n n\|u_n\|^2<\infty$. Since $u\in\F$, adding the already finite sum $\sum_n\|u_n\|^2$ gives the equivalent condition $\sum_n(1+n)\|u_n\|^2<\infty$. Multiplication by $(1+n)^{-1/2}$ is bounded by one, hence defines $j_N$ coordinatewise. It is injective. For $x\in\F$ its image has weighted sum exactly $\sum_n\|x_n\|^2$. Conversely, given a vector in the domain, the sequence $x_n=\sqrt{1+n}\,u_n$ belongs to $\F$ and maps to it. This proves the range and norm assertions. Completeness is inherited from the named coordinate copy of $\F$, not incorrectly from the ambient subtype norm of an unbounded-operator domain.

GW02 gives
$\|\sqrt{n+1}S_nu_{n+1}\|_{L^2(X;H_n)}^2=(n+1)\|u_{n+1}\|^2$.
Their sum is finite by the domain condition. Thus they form an element of the existing \texttt{BochnerSectorSum}. Apply \texttt{bochnerUnitary}; its isometry and coordinate theorem give the asserted norm and every component formula. The countable common-exceptional-set assertion is already provided by \texttt{bochnerReconstruct\_ae}. Linearity follows from linearity of every constituent map. The bound $\|\mathcal A j_Nx\|\leq\|x\|$ makes it a continuous linear map without a second completion argument. The coordinate $n$ depends only on input coordinate $n+1$, proving the projection and vacuum formulas. This explicitly constructs the new momentum field from GW's existing maps; it does not reuse the unrelated electron-position-indexed field.
\end{proof}

\node{GW04}{Restrictions of the maximal GW smeared operators, not replacement domains}{GW03, GW-Fields, GW-Bounds, GW-Adjoints, GW-CCR, GW-Product, ML-BOCHNER}
\begin{lemma}\label{lem:gwrestrictions}
For $f\in L^2(X)$ and physical $u\in V_N$, let
\begin{align*}
 a(f)u&=\texttt{ordinaryAnnihilation}\ f\,
       (\texttt{ordinaryAnnihilationRootInput}\ f\ u),\\
 a^*(f)u&=\texttt{ordinaryCreation}\ f\,
       (\texttt{ordinaryCreationRootInput}\ f\ u).
\end{align*}
The domain witnesses use the physical image $j_Nu$ when $V_N$ is stored in coordinates. These restrictions are continuous from $V_N$ to $\F$, with
\[
 \|a(f)u\|\leq\|f\|\|N^{1/2}u\|,\qquad
 \|a^*(f)u\|\leq\|f\|\|(N+1)^{1/2}u\|.
\]
Moreover
\[
 a(f)u=\int_X\overline{f(k)}\mathcal Au(k)\,dk,
 \qquad \Phi(f)u=a(f)u+a^*(f)u
       =\texttt{ordinaryField}\ f\ u.
\]
The displayed domains do not assert that the maximal domain of either individual smeared operator is exactly $V_N$.
\end{lemma}
\begin{proof}
The input witnesses are existing GW definitions. They use \texttt{annihilation\_mem\_root}, \texttt{creation\_mem\_shiftedRoot}, and equality of the number and shifted-number square-root domains. The norm bounds are \texttt{annihilation\_root\_norm} and \texttt{creation\_root\_norm} with $G_n=H_n$, $R_n=\texttt{ordinarySectorKernel}\ n$, and $F=\texttt{ordinaryKernel}\ f$. The last map sends $z\in\C$ to $zf$, so its norm equals $\|f\|$: the upper bound is $\|zf\|=|z|\|f\|$, and evaluation at $z=1$ gives the reverse bound. These facts supply all generic arguments to the quoted bounds. The graph-coordinate norm in GW03 is exactly the shifted-number norm, so the claimed continuity follows.

The field integral is a Bochner integral because
$\int|f(k)|\|\mathcal Au(k)\|\,dk\leq\|f\|_2\|\mathcal Au\|_2$.
For a finite-sector $u$, pair its $n$th component with arbitrary $w_n\in H_n$. The creation-sector formula \texttt{ordinaryCreationSector\_ae} and sector pairing \texttt{sectorCreation\_pairing} identify the pairing with
\[
 \sqrt{n+1}\int\overline{f(k)}\,
      \langle w_n,S_nu_{n+1}(k)\rangle\,dk.
\]
To justify interchanging the $k$ integral and the sector inner product, the integral of the absolute value of the product integrand is at most
$\sqrt{n+1}\|f\|\|w_n\|\|u_{n+1}\|$, by product-space Cauchy--Schwarz. Thus this is precisely the $n$th component of the asserted Bochner integral. A vector having the same pairing with every $w_n$ is the same sector vector; equality of all coordinates gives equality in $\F$.

For a general $u\in V_N$, use the existing finite particle cutoff $\texttt{P}\ M$. Its convergence in the number graph norm follows by deleting a tail of the convergent sum $\sum_n(1+n)\|u_n\|^2$. The GW annihilation bound controls the difference of the operator values. The Bochner-integral difference is bounded by $\|f\|\|\mathcal A(u-\texttt{P}\ M\,u)\|_2$, which also tends to zero. Therefore the equality extends to $u$. Finally \texttt{ordinaryField\_domain} and \texttt{ordinaryField\_action} give the field sum on exactly this root-number domain. For adjointness and the creation norm identity use the existing \texttt{creation\_annihilation\_pairing} and \texttt{ordinaryCreation\_norm\_identity}; FC06 records their sector calculation, not a second definition. The value of a partial linear map is independent of the proof used to witness the same domain membership, since it is a function on a subtype.
\end{proof}

\node{GW05}{Signed coupling and closed ultraviolet balls}{GW-Model, GW01, GW04, ML-REAL}
\begin{lemma}\label{lem:gwkerneladapter}
Put $\alpha=g^2$ and $\epsilon_g=1$ for $g\geq0$, $\epsilon_g=-1$ otherwise. For every $k$, including $k=0$,
\[
 v_g(k)=\epsilon_g\,\texttt{GW.Model.v}(g^2,k).
\]
Use the closed-ball cutoff $v_{g,\Lambda}=\mathbf1_{\{|k|\leq\Lambda\}}v_g$ throughout this manuscript. For $\Lambda>0$ its $L^2$ class is exactly $\epsilon_g$ times
\texttt{GW.Model.cutoffL2} $(g^2,\Lambda)$, with the proofs $g^2\geq0$, $\Lambda>0$ supplied. In particular, with $v_3=\texttt{GW.Model.unitBallVolume}$,
\[
 \|v_{g,\Lambda}\|_2^2\leq7v_3g^2\Lambda.
\]
\end{lemma}
\begin{proof}
$\sqrt{g^2}=|g|$ and $\epsilon_g|g|=g$ by the two cases in the definition. The kernel equality follows by division by $|k|$, using the shared convention that division at zero is zero. Both cutoff definitions test $|k|\leq\Lambda$, so their functions are equal everywhere, not merely outside a boundary sphere. The existing \texttt{cutoffL2\_ae} identifies their $L^2$ classes, and \texttt{cutoffL2\_norm\_sq} gives the stated bound since $|\epsilon_g|=1$. If a calculation uses an open ball instead, \texttt{spaceSphereVolume} proves that the two classes agree; no such change is needed for the chosen implementation. Real linearity of $a(f)+a^*(f)$ in $f$ follows from the conjugate linearity of $a(f)$ and linearity of $a^*(f)$, so multiplication by the real sign changes the field by the same sign. No Fourier constant or $1/\sqrt2$ is introduced.
\end{proof}

\node{EN01}{The complete fiber energy carrier and its embedding into GW Fock space}{GW01, GW-Weighted, GW-Number, ML-SERIES, ML-BOCHNER, GW-Multiplier}
\begin{lemma}\label{lem:energycoordinates}
For any $P\in X$ let
\[
 w_{n,P}(K_n)=\sqrt{1+|K_n-P|^2+n},\qquad K_n=\sum_{j<n}k_j,
\]
with $K_0=0$. Store $\Q_P$ as a named copy of $\F$ with its existing Hilbert norm, and specify its physical embedding by
$(j_Px)_n=w_{n,P}^{-1}x_n$.
Then $j_P$ is a continuous injective linear map with dense range. Its range is precisely the physical form domain in \eqref{eq:qnorm}. The weighted-coordinate map $u\mapsto(w_{n,P}u_n)_n$ is the inverse from that physical range to $\Q_P$, and
\[
 \|x\|_{\Q_P}^2=\|j_Px\|^2+h_P^{\rm fr}[j_Px].
\]
The analogous construction with weight $\sqrt{1+F_P^2}$ gives the free operator graph space, where $F_P=|K_n-P|^2+n$.
\end{lemma}
\begin{proof}
Each weight is measurable, finite, at least one, and invariant under coordinate permutations. GW's bounded reciprocal multiplier therefore preserves every symmetric sector. On the vacuum it is ordinary multiplication by $(1+|P|^2)^{-1/2}$. The squared output norm is bounded coordinatewise by the squared input norm, so the sum is finite and defines a contraction on the existing Hilbert sum. Linearity and continuity follow coordinatewise. If $j_Px=0$, multiply each almost-everywhere sector equality by the finite positive weight to obtain $x_n=0$. Thus $x=0$.

For $u=j_Px$, the weighted sector sequence is exactly $x$; this equality follows from $w_{n,P}w_{n,P}^{-1}=1$ and the multiplier almost-everywhere formulas. Conversely, if the weighted square sum of $u$ is finite, each weighted component is in the same symmetric sector and the weighted sequence defines $x\in\F$ with $j_Px=u$. The norm identity follows from $w_{n,P}^2=1+F_P$ and addition of nonnegative integrals and series. This is the sectorwise use of the already proved \texttt{WeightedL2.embed\_range} and \texttt{weighted\_norm\_embed}, with GW's scalar vacuum handled separately; no new weighted completion is needed.

For density, truncate a physical $u\in\F$ to $n\leq M$ and all $|k_j|\leq R$. Dominated convergence in each retained sector followed by the $\ell^2$ tail estimate gives convergence to $u$ as $M,R\to\infty$. On this support $w_{n,P}$ is bounded by $\sqrt{1+(MR+|P|)^2+M}$, so each truncated vector belongs to the range of $j_P$. This proves dense range. Replacing $w$ by $\sqrt{1+F_P^2}$ gives the same complete coordinate-space construction and identifies exactly the maximal domain of the free multiplier. EN01 establishes the complete carrier, dense embedding, norm identity, and maximal-domain description; it does not use FA00. Self-adjointness, when required later, is obtained from the retained GW multiplier interface with its hypotheses discharged, or from the subsequent associated-form construction after EN01 has been proved. It is not a premise of this node. In particular we have specified the graph norm, rather than inheriting the incomplete ambient norm on a domain subtype.
\end{proof}

\node{EN02}{Energy norm changes, restrictions, and the dual change-of-coordinates rule}{EN01, GW03, GW01, ML-HILBERT}
\begin{lemma}\label{lem:energytransport}
All physical $\Q_P$ have the same underlying set. For $P,P'\in X$, the map preserving the physical vector is
\[
 I_{P\leftarrow P'}x=(w_{n,P}w_{n,P'}^{-1}x_n)_n,
 \qquad j_PI_{P\leftarrow P'}=j_{P'}.
\]
It and its inverse are bounded, with
$\|I_{P\leftarrow P'}\|^2\leq2(1+|P-P'|^2)$.
Every particle projection acts on these coordinates by the same coordinate selection as on $\F$. The embedding $\Q_P\to V_N$ is bounded with norm at most one in physical graph norms. On the one-particle space the isolated norm $\int(1+|k-P|^2)|f|^2$ and the inherited norm $\int(2+|k-P|^2)|f|^2$ have squared ratio between one and two.

More generally, if $I:V'\to V$ is an invertible bounded coordinate change preserving the physical embedding, then a form representative $T$ on $V$ and a source Riesz vector $z$ transform as
\[
 T'=I^*TI,\qquad z'=I^*z,\qquad
 (T')^{-1}=(I^{-1})T^{-1}(I^*)^{-1}.
\]
Their physical weak solution is unchanged.
\end{lemma}
\begin{proof}
Using $|K-P|^2\leq2|K-P'|^2+2|P-P'|^2$ gives
$1+n+|K-P|^2\leq2(1+|P-P'|^2)(1+n+|K-P'|^2)$.
Apply the resulting pointwise bound to the multiplier ratio and sum its squared norm. Interchanging $P,P'$ gives the inverse. Multiplication of the ratios is one, proving the inverse identity. Coordinate projections commute with all weights and therefore remain orthogonal contractions. Since $1+n\leq w_{n,P}^2$, the embedding into $V_N$ follows by comparing sums. The one-particle comparison follows from $1+|k-P|^2\leq2+|k-P|^2\leq2(1+|k-P|^2)$.

For the last assertions, take arbitrary $v',u'$. The transformed form is
$\langle Iv',TIu'\rangle_V=\langle v',I^*TIu'\rangle_{V'}$,
and the transformed load is $\langle Iv',z\rangle_V=\langle v',I^*z\rangle_{V'}$. Multiplying the proposed inverse on each side verifies both identities. If $u'= (T')^{-1}z'$, then $Iu'=T^{-1}z$, and the two physical embeddings consequently give the same vector. This proof also covers the sectorwise norm changes and the shifts of momentum used later. GW's auxiliary \texttt{EnergySpace d} has the form inner product; it is not silently substituted for the fixed coordinate space just constructed.
\end{proof}

\section{Explicit foundational constructions and legitimate leaves}\label{sec:foundations}
This section closes the construction steps previously concealed in the words ``by density'' or ``use the Fock realization''. A foundational leaf in this document is either one of the precise library results in Section~\ref{sec:mathlib}, elementary finite-dimensional real/complex algebra, or an earlier local lemma. In particular, no theorem asserting a pull-through identity, an associated Hamiltonian, an $L^2$ currying map, or a weighted integral operator is a new assumption. GW01--GW05 and EN01--EN02 specify which existing constructions implement the interfaces. The explanatory foundational proofs below are retained for mathematical readability; they are not alternative implementation branches.

\subsection{Measure-theoretic extension and scalar sections}
\node{MS01}{Rectangle density with its generating-set argument}{ML-LP, ML-SIMPLE, ML-PROD, ML-VOLUME, ML-BOCHNER, GW-Product}
\begin{lemma}[Rectangle approximation]\label{lem:rectdense}
Let $X=\R^a$ and $Y=\R^b$, with Lebesgue measure and their Borel sigma algebras. The complex linear span of
$\mathbf1_A(x)\mathbf1_B(y)$, for measurable finite-measure bounded sets $A,B$, is dense in $L^2(X\times Y)$. The same assertion holds with each factor restricted to a bounded box.
\end{lemma}
\begin{proof}
First work in a fixed product box $D=D_X\times D_Y$ of finite measure, with its restricted product measure. Let $W$ be the closed linear span of its rectangle indicators. Put
$\mathscr D=\{E\subset D:E\text{ measurable and }\mathbf1_E\in W\}$.
The whole box is in $\mathscr D$. If $E\in\mathscr D$, then $\mathbf1_{D\setminus E}=\mathbf1_D-\mathbf1_E\in W$. If $E_j$ are pairwise disjoint members, the indicator of their union is the $L^2$ limit of $\sum_{j\leq m}\mathbf1_{E_j}$: the squared error is the measure of the tail union and tends to zero by countable additivity and finiteness of the measure of $D$. Thus $\mathscr D$ is a Dynkin system containing the measurable rectangles, a family closed under finite intersections.

Here is the generating-set step, rather than an unspecified appeal to it. For a family $\mathscr P$ closed under intersections, let $\mathscr D_0$ be the intersection of all Dynkin systems containing $\mathscr P$ and the ambient set. A Dynkin system is closed under $B\setminus A$ when $A\subset B$ are its members: write this as the complement of the disjoint union $A\cup(D\setminus B)$. For fixed $A\in\mathscr P$, the family of $B\in\mathscr D_0$ such that $A\cap B\in\mathscr D_0$ is a Dynkin system: use that difference observation for complements and distribute intersections over disjoint unions. It contains $\mathscr P$, so equals $\mathscr D_0$. Now fix $B\in\mathscr D_0$ and repeat with the variable $A$; the family again contains $\mathscr P$ by the preceding conclusion. It follows that $\mathscr D_0$ is closed under finite intersections. Together with complements this gives finite unions; disjointifying a countable union using successive differences gives closure under all countable unions. Hence $\mathscr D_0$ is a sigma algebra and contains the sigma algebra generated by $\mathscr P$.

Apply this argument to rectangles. They generate the product sigma algebra by its definition. Therefore every measurable-set indicator on $D$ belongs to $W$. Finite linear combinations give all simple functions on $D$. The density of simple functions in $L^2$ is the precise leaf ML-SIMPLE; its hypotheses hold because the exponent is $2<\infty$ and $\C$ is a complete normed space. Thus $W=L^2(D)$.

Finally take increasing integer-radius product boxes $D_j$. For $u\in L^2(X\times Y)$, $\mathbf1_{D_j}u\to u$ in $L^2$: the squared error is $\int\mathbf1_{D_j^c}|u|^2$, which tends to zero by dominated convergence with majorant $|u|^2$. Approximate the truncated vector inside the box by the preceding argument, and extend those rectangles by zero. All $L^2$ classes here are those of the fixed Borel measure, so this construction requires no choice of a representative for a completed sigma algebra. No evaluation at a prescribed point is part of the conclusion.
\end{proof}

\node{MS02}{The scalar-product to Bochner-$L^2$ identification}{MS01, ML-LP, ML-L2, ML-PROD, ML-COMPLETE, ML-MEAS, GW02}
\begin{lemma}[Currying as an isometry, with its section statement]\label{lem:L2curry}
There is a linear isometry
\[
 \mathscr U:L^2(X\times Y;\C)\longrightarrow L^2(X;L^2(Y;\C))
\]
which sends a rectangle tensor $f(x)g(y)$ to $x\mapsto f(x)[g]$ whenever these expressions are square integrable. If $u$ is a measurable representative of the source, then for almost every $x$, its section $u(x,\cdot)$ is in $L^2(Y)$ and represents $(\mathscr U[u])(x)$. The result is independent of both source and target representatives. If a finite group acts by measure-preserving transformations in the $Y$ variables and fixes the source class, the target can be regarded as a strongly measurable field in its closed fixed subspace, with the same norm.
\end{lemma}
\begin{proof}
On a finite sum of rectangle tensors, the proposed field is strongly measurable: each scalar measurable function times a fixed vector is strongly measurable, and finite sums preserve this property. Fubini, applied to the integrable squared norm, gives
\[
 \int_X\Big\|\sum_i f_i(x)g_i\Big\|_{L^2(Y)}^2\,dx
 =\int_{X\times Y}\Big|\sum_i f_i(x)g_i(y)\Big|^2\,d(x,y).
\]
For rectangle indicators integrability follows from finite measure and a finite sum; for general square-integrable factors it follows from Cauchy--Schwarz on the finite cross terms. The norm identity shows that two expressions defining the same scalar $L^2$ vector define the same Bochner $L^2$ vector. Thus the map is well-defined and linear on the dense subspace in MS01.

The extension can be constructed without an unspecified completion step. Given $u$, choose $u_j$ in that subspace with $\|u_j-u\|_2\leq2^{-j}$. The images are Cauchy because their differences have exactly the same norm as $u_j-u_l$. Bochner $L^2$ is complete by ML-COMPLETE, so define $\mathscr Uu$ as their limit. Any other approximating sequence has the same limit by the difference estimate. Addition and scalar multiplication commute with limits, proving linearity. Taking limits of the norm identity gives isometry and uniqueness.

Tonelli applied to $|u|^2$ first gives square-integrable sections outside a null set. To identify them with the constructed field, choose the above approximants so that $\sum_j\|u_j-u\|_2^2<\infty$. Tonelli yields
\[
 \int_X\sum_j\|u_j(x,\cdot)-u(x,\cdot)\|_{L^2(Y)}^2\,dx<\infty.
\]
Consequently the summands tend to zero for almost every $x$. In the target, the isometry gives $\sum_j\|\mathscr Uu_j-\mathscr Uu\|_{L^2(X;L^2(Y))}^2<\infty$; a second use of Tonelli gives pointwise convergence in $L^2(Y)$ outside another null set. Each rectangle approximation already agrees with its section; taking the countable union of these exceptional sets proves the claim. If two representatives of $u$ agree almost everywhere on the product, Tonelli applied to their zero discrepancy shows their sections agree almost everywhere in $Y$ for almost every $x$. The target assertion follows from the same equivalence-class rule.

A measure-preserving transformation $S$ of $Y$ induces an isometry $U_S$ of $L^2(Y)$, by change of variables. If $u(x,Sy)=u(x,y)$ as product classes, Fubini gives $U_S[u(x,\cdot)]=[u(x,\cdot)]$ for almost every $x$. A finite intersection of these full-measure sets suffices for a finite group $G$. The common fixed subspace is closed because it is the intersection of the kernels of the continuous maps $U_S-I$.

Here is the needed strong measurability in the subspace, rather than only in its ambient space. Define the bounded averaging map $\mathsf S=|G|^{-1}\sum_{S\in G}U_S$. Left multiplication permutes the elements of $G$, so $U_T\mathsf S=\mathsf S$; also $\mathsf S$ is the identity on fixed vectors and has norm at most one. Thus its range is exactly the fixed subspace. Choose ambient-space simple functions $F_j$ converging almost everywhere to a representative $F$ of $\mathscr Uu$, as in the definition of strong measurability. The functions $\mathsf S F_j$ are simple functions valued in the fixed subspace, and $\|\mathsf SF_j-F\|\leq\|F_j-F\|$ at every point where $F$ is fixed. Their almost-everywhere convergence proves strong measurability as a subspace-valued field. The inclusion is isometric, so its squared norm and integral are unchanged. This applies in particular to the finite permutation groups defining symmetric Fock sectors.
\end{proof}

\node{MS03}{Countable Hilbert sums and fields assembled from components}{ML-L2, ML-LP, ML-COMPLETE, ML-HSUM, ML-PROD, ML-MEAS, GW-Bochner}
\begin{lemma}[Component assembly]\label{lem:fieldassembly}
Let $(H_n)$ be a sequence of complex Hilbert spaces, and let $\mathscr H=\bigoplus_{n\geq0}H_n$ be their square-summable Hilbert sum. If $F_n\in L^2(X;H_n)$ and $\sum_n\|F_n\|_2^2<\infty$, there is a unique $F\in L^2(X;\mathscr H)$ whose $n$th coordinate is $F_n$, and
\[
 \|F\|_2^2=\sum_n\|F_n\|_2^2.
\]
Its representatives can be chosen so that $\sum_n\|F_n(x)\|^2<\infty$ and $F(x)=(F_n(x))_n$ for almost every $x$.
\end{lemma}
\begin{proof}
For the finite partial sums $F^{(m)}=\sum_{n\leq m}\iota_nF_n$, strong measurability follows from finite addition and the continuous coordinate inclusions. Orthogonality gives
$\|F^{(m)}-F^{(l)}\|_2^2=\sum_{l<n\leq m}\|F_n\|_2^2$ for $m\geq l$. Thus the partial sums are Cauchy in the complete Bochner space. Their limit $F$ has the stated coordinates because each projection is continuous of norm at most one. Conversely zero coordinates imply squared norm zero, proving uniqueness. The norm identity follows by taking limits of the finite identities.

For the representative assertion, Tonelli gives
$\int\sum_n\|F_n(x)\|^2=\sum_n\|F_n\|_2^2<\infty$.
Thus the pointwise component sequence is square summable outside a null set. Its finite sums converge there in $\mathscr H$. They are strongly measurable, hence so is their almost-everywhere limit (ML-MEAS). The squared norm is integrable by that same equality. It defines a Bochner $L^2$ vector with the same coordinates as $F$, and uniqueness identifies it with $F$.
\end{proof}

\subsection{A fixed Fock realization and all graph-space completions}
\node{FD01}{Canonical symmetric sectors and their projections}{MS02, MS03, ML-LP, ML-HSUM, ML-VOLUME, GW01}
\begin{definition}[Description of the fixed GW realization]\label{def:canonicalFock}
Write $X_n=(\R^3)^n$, $X_0=\{\ast\}$ with mass one. The existing GW sectors are $H_0=\C$, identified isometrically with $L^2(X_0)$ by GW01, and for $n\geq1$ the subspaces of $L^2(X_n;\C)$ invariant under every permutation. The already fixed carrier is $\F=\bigoplus_{n\geq0}H_n$ and its existing vacuum is $\Omega=(1,0,\ldots)$. All sector norms use the ordinary product measure, without factorial weights. Define $K_n=k_1+\cdots+k_n$. Sector inclusions and projections are the coordinate maps of this Hilbert sum.
\end{definition}
\begin{proof}[Verification of the construction]
A coordinate permutation preserves product Lebesgue measure: on finite-measure rectangles the measures are products of the coordinate measures and permutation does not change that product. Equality of the two sigma-finite measures follows by the generating-set argument in MS01, first on boxes and then by exhaustion. Composition therefore induces a linear isometry of $L^2(X_n)$. Its fixed subspace is closed as the kernel of its difference from the identity. There are finitely many permutations, so $H_n$ is a closed Hilbert subspace. The finite average
\[
 S_nf=\frac1{n!}\sum_{\pi\in\mathfrak S_n}f\circ\pi
\]
is a contraction, maps into $H_n$ by reindexing the finite sum, and restricts to the identity on $H_n$. In particular it is a projection. It is self-adjoint: changing variables by $\pi$ changes the adjoint summand to $\pi^{-1}$, and inversion permutes the finite group. The Hilbert sum exists by ML-HSUM; its inner product is $\sum_n\langle u_n,w_n\rangle$, absolutely convergent by Cauchy--Schwarz for square-summable numerical sequences. Coordinate inclusion is isometric and coordinate projection is contractive, with adjoint inclusion. Finite sector sums converge in norm because the sum of squared tail norms tends to zero. The vacuum has norm one. These facts fix all normalizations used later.
\end{proof}

\node{FD02}{Weighted domains, dense embeddings and multiplication operators}{FD01, ML-LP, ML-L2, ML-COMPLETE, EN01, EN02, GW-Multiplier}
\begin{lemma}[Weighted Hilbert spaces]\label{lem:weightedspace}
Let $w_n:X_n\to[1,\infty)$ be finite measurable permutation-invariant functions. The space
\[
 V_w=\{u\in\F:\sum_n\int w_n^2|u_n|^2<\infty\},\qquad
 \langle u,z\rangle_{V_w}=\sum_n\int w_n^2\overline{u_n}z_n
\]
is a Hilbert space. Multiplication $u\mapsto(w_nu_n)_n$ is an isometric isomorphism $V_w\to\F$, with inverse multiplication by $w_n^{-1}$. The inclusion $j:V_w\to\F$ is continuous, injective, and has dense range.
\end{lemma}
\begin{proof}
Multiplication by a finite measurable function is well-defined on measurable equivalence classes: equal representatives have equal products off the same null set. The weighted sum condition is exactly membership of the product sequence in $\F$. Conversely, given $z\in\F$, the vector $u_n=z_n/w_n$ has unweighted norm at most $\|z\|$, belongs to the symmetric sector, and satisfies $wu=z$ almost everywhere. This gives mutually inverse linear maps and proves all inner-product and completeness properties by transport from $\F$. The inclusion has norm at most one. To prove dense range, truncate any $z\in\F$ in sector number and in the weight: $z^{(m)}_n=\mathbf1_{n\leq m}\mathbf1_{w_n\leq m}z_n$. This vector lies in $V_w$ and converges to $z$ in $\F$, because the squared discrepancy is dominated by the summable sector norms and the indicators tend to one. The same argument with the weighted integrable majorant proves graph-norm approximation for vectors already in $V_w$.
\end{proof}

For each sector let $a_n:X_n\to\R$ be finite and measurable, and suppose that $a_n\circ\pi=a_n$ almost everywhere for every permutation $\pi\in\mathfrak S_n$. Multiplication then preserves bosonic symmetry and, on the maximal domain
$\{u\in\F:\sum_n\int |a_n|^2|u_n|^2<\infty\}$, defines a densely defined closed operator on $\F$. The almost-everywhere symmetry assumption is needed for this bosonic domain and its density. Without it, the multiplication graph restricted to symmetric inputs and symmetric outputs is still closed, but need not be densely defined. Here is the closedness check used below. If $u_j\to u$ and $au_j\to z$ in $\F$, choose a subsequence with summable squared errors in each of the two norms. Tonelli, on the disjoint union of sectors, gives almost-everywhere convergence of both representatives. Since each $a_n$ is finite almost everywhere, $z_n=a_nu_n$ there. Thus $u$ is in the multiplication domain and $au=z$. Density follows by the same weight truncations, which remain symmetric as $L^2$ classes by the assumed almost-everywhere permutation invariance. For the nonnegative free multiplier
$F_P(K_n)=|K_n-P|^2+n$, the form space is $V_{\sqrt{1+F_P}}$ and the operator domain is $V_{\sqrt{1+F_P^2}}$. The form-domain associated operator is precisely multiplication by $F_P$: if $F_Pu\in\F$, the weak identity follows by Cauchy--Schwarz; conversely test the weak identity against bounded-weight, finite-sector truncations of $F_Pu-f$ to prove $F_Pu=f$ almost everywhere and in $\F$. This last argument uses tests $z_m=\mathbf1_{n\leq m,F_P\leq m}(F_Pu-f)$; they are in the form domain, and the identity gives $\|z_m\|^2=0$.

Sector projections preserve every $V_w$ and are contractions in its weighted norm. Since the weights are real, componentwise complex conjugation is a conjugate-linear isometry of $V_w$. Momentum reflection $u_n(K_n)\mapsto u_n(-K_n)$ is an isometry from $V_w$ onto $V_{\widetilde w}$ when $\widetilde w_n(K_n)=w_n(-K_n)$ almost everywhere; it preserves $V_w$ when the weights are reflection-invariant. In particular, reflection identifies the free form and operator graph norms at $P$ with those at $-P$.

Unbounded number and momentum multipliers are used only on their actual multiplication domains. Commutation with a weight identifies products on the joint domains where the required products belong to $\F$; it does not by itself imply preservation of $V_w$. Any stronger mapping property or extension must be justified by a separate estimate. FC01 supplies the bounded-coordinate, finite-sector core invariance used for commutators. For two weights with $0<c\leq C$ and $c w\leq \widetilde w\leq Cw$ almost everywhere, the underlying spaces agree and their norms are equivalent with exactly these constants. This proves, rather than assumes, the common form-domain and momentum-norm equivalence assertions in Section~\ref{sec:target}.

\subsection{Bounded forms, inverses, and a fully defined associated operator}
\node{BA01}{Riesz coordinates for conjugate-dual sources and forms}{ML-RIESZ}
\begin{lemma}\label{lem:rieszcoordinates}
For a continuous conjugate-linear functional $F$ on a complex Hilbert space $V$, there is a unique $z_F\in V$ with $F(w)=\langle w,z_F\rangle$ and $\|z_F\|=\|F\|$. The assignment is complex-linear. If $a(w,u)$ is conjugate-linear in $w$, linear in $u$, and $|a(w,u)|\leq C\|w\|\|u\|$, there is a unique bounded linear $T$ with $a(w,u)=\langle w,Tu\rangle$ and $\|T\|\leq C$. If $a$ is Hermitian, $T$ is self-adjoint.
\end{lemma}
\begin{proof}
Apply the ordinary linear-dual Riesz correspondence ML-RIESZ to the linear functional $w\mapsto\overline{F(w)}$. It represents this functional as $\langle z_F,w\rangle$; conjugating gives the required equality. Norm preservation is part of that correspondence. Uniqueness also follows directly: test the difference against its own representing vector. For a fixed $u$, apply this construction to $F_u(w)=a(w,u)$. Uniqueness implies $z_{F_{u+v}}=z_{F_u}+z_{F_v}$ and $z_{F_{cu}}=cz_{F_u}$. The norm bound gives boundedness of $Tu=z_{F_u}$. Hermitian symmetry gives $\langle w,Tu\rangle=\langle Tw,u\rangle$, which is exactly the defining equality for the bounded adjoint. This also constructs $j^*$ for any bounded Hilbert-space map $j$: apply the same argument to $(w,f)\mapsto\langle jw,f\rangle$.
\end{proof}

\node{BA02}{Geometric inverses, continuity, and order on dual sources}{BA01, ML-COMPLETE, ML-SERIES}
\begin{lemma}\label{lem:inversecalculus}
If $\|K\|\leq r<1$ on a Banach space, $I-K$ has inverse $\sum_{j\geq0}K^j$, of norm at most $(1-r)^{-1}$. For bounded invertible $S,T$,
\[
 S^{-1}-T^{-1}=S^{-1}(T-S)T^{-1}.
\]
If $T$ is positive and invertible on a Hilbert space, its inverse is positive and
\[
 \langle z,T^{-1}z\rangle
 =\max_{u\in V}\{2\Re\langle z,u\rangle-\langle u,Tu\rangle\}.
\]
Consequently $T_2\geq T_1>0$ implies
$\langle z,T_2^{-1}z\rangle\leq\langle z,T_1^{-1}z\rangle$ for every $z$.
\end{lemma}
\paragraph{Implementation contract: reuse the Banach and ring interfaces.}
Use the existing complete space of continuous linear maps (ML-COMPLETE), not a second pointwise-completion construction. In the normed ring of bounded endomorphisms, ML-SERIES supplies \path{Units.oneSub}, \path{geom_series_mul_neg}, \path{mul_neg_geom_series}, and the geometric norm estimate. Match multiplication with composition and prove both inverse identities for the chosen inverse. The scalar formula \path{hasSum_geometric_of_norm_lt_one} has a division-ring hypothesis and is not the operator inverse interface. The inverse-order and form-dual specializations below remain local lemmas if no exact retained result already covers them. The explanatory proof does not instruct the implementer to reprove completeness or geometric summability.
\begin{proof}
The finite partial sum has norm at most $\sum_{j\leq m}r^j$, and its tails have norm at most $r^{l+1}/(1-r)$. The required completeness of bounded maps follows directly from completeness of the target Banach space: for an operator-norm Cauchy sequence $A_j$, define $Ax=\lim_j A_jx$ for each $x$. Pointwise limits preserve addition and scalar multiplication; a uniform bound $\|A_j\|\leq M$ gives $\|Ax\|\leq M\|x\|$. Given $\epsilon>0$, choose $j_0$ with $\|A_j-A_l\|\leq\epsilon$ for $j,l\geq j_0$ and pass $l\to\infty$ for each $x$, obtaining $\|(A_j-A)x\|\leq\epsilon\|x\|$. This is operator-norm convergence. Applying it to the partial sums gives the asserted limit. Multiplication by $I-K$ on either side of the finite sum gives $I-K^{m+1}$; continuity of composition and $r^{m+1}\to0$ prove both inverse identities. The norm bound passes to the limit. The two-inverse identity follows by multiplying out its right side. In particular uniformly bounded inverses are norm continuous wherever their underlying operators are norm continuous.

The inverse of a self-adjoint invertible operator is self-adjoint: take adjoints in its two inverse identities and use uniqueness of a two-sided inverse. For positivity put $y=T^{-1}z$; then $\langle z,T^{-1}z\rangle=\langle Ty,y\rangle=\langle y,Ty\rangle\geq0$. Expanding with $u=T^{-1}z+w$ gives
\[
 2\Re\langle z,u\rangle-\langle u,Tu\rangle
 =\langle z,T^{-1}z\rangle-\langle w,Tw\rangle.
\]
The maximum is attained at $w=0$. If $T_2\geq T_1$, each maximized expression for $T_2$ is at most the corresponding one for $T_1$; taking the maxima proves inverse order. This applies to energy-space representatives and therefore also to form-dual loads, not only physical vectors.
\end{proof}

\node{BA03}{Mixed form bounds from diagonal bounds}{BA01, BA02, GW-Mixed}
\begin{lemma}\label{lem:polarizationbound}
If a Hermitian form $a$ on a complex Hilbert space obeys $|a(u,u)|\leq M\|u\|^2$, then
$|a(w,u)|\leq 2M\|w\|\|u\|$. A sharper constant is unnecessary here. If a sequence of bounded forms is Cauchy in this mixed norm, it has a bounded-form limit. Hermitian symmetry, sector-selection equalities, and uniform diagonal inequalities pass to the limit.
\end{lemma}
\begin{proof}
For unit vectors $w,u$, the complex polarization identity expresses $4a(w,u)$ as a signed sum of the four diagonal values at $w+i^ju$, $j=0,1,2,3$, with unit complex coefficients. The sum of their squared norms is $8$, by expansion and cancellation of the four mixed terms. Hence $|a(w,u)|\leq2M$. For nonzero vectors apply this to their normalizations and multiply by their norms; zero vectors give zero by sesquilinearity. The bounded-form space is complete: BA01 identifies it isometrically with the complete space of bounded linear operators (or, equivalently, take the pointwise scalar limits and pass the uniform mixed estimate to the limit). Each asserted equality or inequality passes to the scalar limit for fixed tests. This proves the last claims, including preservation of vanishing matrix entries between particle sectors.
\end{proof}

\node{DF01}{Domain and spectrum definitions on the fixed partial-operator interface}{BA01, GW-Resolvent}
The following are mathematical definitions of the interfaces, not a second operator library. The implemented operators are the retained GW partial linear maps; GW06 and GW07 establish these descriptions for the actual Hamiltonian and inverse predicate.

\begin{definition}[Associated operator]\label{def:associated}
Let $V,H$ be Hilbert spaces and $j:V\to H$ a continuous injective linear map with dense range. Let $a$ be a bounded Hermitian sesquilinear form on $V$. Its associated operator $A$ is characterized by
\[
 ju\in\Dom A,\quad Aju=f
 \quad\Longleftrightarrow\quad
 a(w,u)=\ip{jw}{f}_H\quad\text{for all }w\in V.
\]
The vector $f$ is unique by density. For $a_\lambda(w,u)=a(w,u)-\lambda\ip{jw}{ju}$, let $T_\lambda\in\mathcal B(V)$ satisfy
\[
 \ip{w}{T_\lambda u}_V=a_\lambda(w,u).
\]
Riesz representation constructs $T_\lambda$; it is self-adjoint on $V$.
\end{definition}

\begin{definition}[Real resolvent without a spectral theorem]\label{def:resolvent}
For a densely defined operator $A$ on $H$, a real number $\lambda$ lies in $\resR(A)$ if there is a bounded complex-linear $R:H\to H$ such that
\[
 Rf\in\Dom A,\quad (A-\lambda)Rf=f\quad(f\in H),
 \qquad R(A-\lambda)u=u\quad(u\in\Dom A).
\]
Set $\sigR(A)=\R\setminus\resR(A)$. This is the usual real slice of the resolvent/spectrum definition, not the point spectrum.
\end{definition}

\node{FA00}{Construction and uniqueness of the associated self-adjoint operator}{BA01, BA02, FD02, ML-HILBERT, GW-Form, DF01}
\begin{theorem}[Associated operator: mathematical explanation of the reused GW theorem]\label{thm:associatedconstruction}
Let $j:V\to H$ be bounded, injective, and have dense range. Let $a$ be a bounded Hermitian form on $V$. Suppose for a real $c$ and $\alpha>0$,
\[
 a(u,u)+c\|ju\|^2\geq\alpha\|u\|_V^2.
\]
There is a unique densely defined self-adjoint operator $A$ on $H$ with the associated-form domain characterization in Definition~\ref{def:associated}. Its domain is the range of the bounded injective map
$R=jT^{-1}j^*$, where $\langle w,Tu\rangle=a(w,u)+c\langle jw,ju\rangle$, and
$A(Rf)=f-cRf$.
\end{theorem}
\begin{proof}
This proof explains the existing GW representation theorem; the implementation uses GW06 and the indexed source. The coercivity inverse argument is written here to avoid a forward logical dependency in the mathematical explanation. Choose $M\geq\max(\|T\|,\alpha)$, put $\eta=\alpha/M^2$, and expand
$\|u-\eta Tu\|^2\leq(1-\alpha^2/M^2)\|u\|^2$. Thus $\|I-\eta T\|<1$ unless it is zero, in which case inversion is immediate. BA02 gives the inverse in either case and $\|T^{-1}\|\leq\alpha^{-1}$ follows from $\alpha\|u\|\leq\|Tu\|$. The map $R$ is bounded and self-adjoint. It is positive by BA02. If $Rf=0$, then with $z=j^*f$ and $y=T^{-1}z$,
\[
 0=\langle f,Rf\rangle=\langle z,T^{-1}z\rangle
   =\langle y,Ty\rangle\geq\alpha\|y\|^2.
\]
Thus $z=0$, and $\langle jw,f\rangle=0$ for every $w$. Density gives $f=0$. The range of $R$ is dense: a vector orthogonal to it satisfies $Ry=0$ by self-adjointness, hence is zero; the orthogonal-complement characterization of closure for a linear subspace is the standard Hilbert-space projection fact ML-HILBERT.

Define $B=R^{-1}$ on $\operatorname{Ran}R$ and $A=B-cI$. Injectivity makes $B$ well-defined and linear. If $u_j\to u$ and $Bu_j\to v$, the identity $u_j=RBu_j$ yields $u=Rv$, proving closedness. To prove self-adjointness without a spectral theorem, a vector $y$ belongs to the adjoint domain of $B$ precisely when some $z$ satisfies $\langle Bx,y\rangle=\langle x,z\rangle$ for all $x\in\operatorname{Ran}R$. Substituting $x=Rf$ makes this
$\langle f,y\rangle=\langle Rf,z\rangle=\langle f,Rz\rangle$ for all $f$. It is equivalent to $y=Rz$, and then $By=z$. This identifies both the adjoint domain and the adjoint value with those of $B$. A real scalar shift preserves this equality, proving self-adjointness of $A$.

For $u=T^{-1}j^*f$, the definition gives
$a(w,u)+c\langle jw,ju\rangle=\langle jw,f\rangle$ and $ju=Rf$. Thus $Aju=f-cju$, establishing the associated identity on its domain. Conversely if $u\in V$ satisfies $a(w,u)=\langle jw,h\rangle$ for all $w$, BA01 gives $Tu=j^*(h+cju)$. Hence $ju=R(h+cju)\in\operatorname{Ran}R$ and $Aju=h$. This proves the exact domain characterization. Any other operator with that characterization has the same domain and values, so uniqueness is literal equality, not just unitary equivalence.
\end{proof}


\section{Concrete Fock calculus and measurable annihilation fields}\label{sec:fockdetail}
This section fixes every sector calculation used below on GW01's carrier, using GW02--GW04 for the actual maps and EN01--EN02 for their graph coordinates. No operation evaluates an arbitrary $L^2$ class at a prescribed momentum.

\node{FC01}{A bounded-momentum algebraic core}{FD01, FD02, ML-BOCHNER, EN01, GW04}
For integers $M,R\geq1$ let $J_{M,R}$ keep sectors $n\leq M$ and multiply the $n$th component by
\[
 \prod_{i=1}^{n}\mathbf1_{\{\abs{k_i}\leq R\}}.
\]
For $n=0$ the product is $1$. This multiplier is invariant under permutations, so it preserves the symmetric sector. Let $\mathscr C$ be the union of the ranges of these projections. Its vectors need not be smooth. No differentiation in a particle momentum occurs in a commutator calculation in this proof.

\begin{lemma}\label{lem:boundedcore}
The space $\mathscr C$ is a linear, dense subspace of $\F$, of $\Dom(N^{1/2})$ in its graph norm, of $\Q$ in its form norm, and of the free operator domain in its graph norm. Every sector projection preserves it. If $f\in L^2(\R^3)$ has bounded momentum support, then $a(f),a^*(f)$ map $\mathscr C$ into $\mathscr C$. The free momentum and number multipliers also preserve $\mathscr C$.
\end{lemma}
\begin{proof}
Write $F_P(K_n)=\abs{K_n-P}^2+n$. On the $n$th sector the four squared norms are integrals with weights $1$, $1+n$, $1+F_P$, and $1+F_P^2$, respectively. For $M=R=j\to\infty$, the multipliers of $J_{j,j}$ increase to $1$ pointwise on every finite configuration. Applying dominated convergence on each sector and then to the summable nonnegative sequence of weighted sector norms proves convergence in each of the four norms. Equivalently use monotone convergence on the disjoint union of the configuration spaces. The bound on each coordinate, not merely on their sum, ensures boundedness of every polynomial multiplier used on the core.

The creation formula is
\[
 (a^*(f)u)_n(k_1,\ldots,k_n)
 =\frac1{\sqrt n}\sum_{i=1}^{n}f(k_i)
           u_{n-1}(k_1,\ldots,\widehat{k_i},\ldots,k_n),\quad n\geq1,
\]
with zero vacuum component. Annihilation is the integral formula stated below. These formulas increase or decrease the largest occupied sector by one and preserve bounded coordinate support, with radius enlarged to include the support of $f$. All products and integrals exist by Cauchy--Schwarz. This proves invariance. The union is linear because two radii and two sector bounds can be replaced by their maxima.
\end{proof}
\paragraph{Implementation point.}
Use $\mathscr C$ for cutoff commutators. Smooth-core density may be reused from GW, but is unnecessary here. In particular a sharp ultraviolet cutoff causes no smoothness problem. On the $n$-particle component use one fixed coordinate ordering; permutation symmetry supplies independence of the chosen removed coordinate and the factors $\sqrt n$.

\node{FC02}{The annihilation map into a Bochner $L^2$ space}{FC01, MS02, MS03, FD01, FD02, GW03}
Let $V_N=\Dom(N^{1/2})$, with norm
$\norm u_{V_N}^2=\norm u^2+\norm{N^{1/2}u}^2$.
For $u\in\mathscr C$ define the field $\mathcal A u$ by
\begin{equation}
 ((\mathcal A u)(k))_n(k_1,\ldots,k_n)
   =\sqrt{n+1}\,u_{n+1}(k,k_1,\ldots,k_n).
 \label{eq:fielddefinition}
\end{equation}
This is an identity of almost-everywhere equivalence classes, not a choice of values of an $L^2$ function at every point.

\begin{lemma}\label{lem:annihilationmap}
There is a unique continuous complex-linear map
\[
 \mathcal A:V_N\longrightarrow L^2(\R^3;\F)
\]
with \eqref{eq:fielddefinition} on the core, and
\begin{equation}
 \norm{\mathcal A u}_{L^2(\R^3;\F)}^2
   =\sum_{m\geq1}m\norm{u_m}_2^2
   =\norm{N^{1/2}u}^2.\label{eq:annihilationisometry}
\end{equation}
For $n\geq1$ it satisfies
$\mathcal A\Pi^{\geq n}=\Pi^{\geq n-1}\mathcal A$;
also $\mathcal A\Omega=0$ and $(\mathcal A\iota_1f)(k)=f(k)\Omega$ a.e.
Here $\iota_1$ is the given one-sector isometry.
\end{lemma}
\begin{proof}
For each input sector $u_{n+1}$, apply MS02 to distinguish its first momentum variable from the other $n$. Symmetry under permutations of the remaining variables places the resulting field in $H_n$ almost everywhere. Multiply that field by $\sqrt{n+1}$. MS02's isometry gives its squared $L^2$ norm as $(n+1)\|u_{n+1}\|^2$. If $u\in V_N$, the sum of these norms is finite. MS03 assembles the coordinates into one strongly measurable Fock-valued $L^2$ field with exactly the norm in \eqref{eq:annihilationisometry}. This constructs the map directly for every $u\in V_N$ and proves linearity, independence of representatives and the norm bound. On the core the section statement in MS02 proves \eqref{eq:fielddefinition}. FC01's density and continuity give uniqueness of a map with these core values. The projection identities hold in each coordinate: the input component $n+1$ survives $\Pi^{\geq m}$ precisely when $n\geq m-1$. For the vacuum all output components vanish; for a one-particle vector only output sector zero survives and its coefficient is $\sqrt1=1$. Coordinate uniqueness in MS03 proves the asserted field identities.
\end{proof}
\paragraph{No hidden evaluation operator.}
The notation $a_ku$ will, when used at all, abbreviate a representative of $(\mathcal A u)(k)$. It never denotes a linear operator on all of $\F$ for a fixed $k$. The expression $a_k^*\Omega$ is not a Hilbert vector; this document's proofs do not construct it.

\subsection{The sector calculations behind every field commutator}
\node{FC06}{Smeared adjointness and the creation norm, with all indices}{FC01, FC02, MS02, MS03, ML-PROD, GW04}
\begin{lemma}\label{lem:sectorcalculation}
For $f\in L^2(\R^3)$, the creation and annihilation formulas on the finite-sector core satisfy
\[
 (a(f)u)_n(K_n)=\sqrt{n+1}\int\overline{f(l)}u_{n+1}(l,K_n)\,dl,
\]
\[
 \|a(f)u\|^2\leq\|f\|_2^2\sum_m m\|u_m\|_2^2,
 \qquad \|a^*(f)u\|^2=\|a(f)u\|^2+\|f\|_2^2\|u\|^2.
\]
They obey $\langle a^*(f)w,u\rangle=\langle w,a(f)u\rangle$ whenever $w,u$ have finitely many sectors. Bounded momentum support of $f$ is needed only for invariance of the bounded-coordinate core, not for these identities or for extension to $V_N$.
\end{lemma}
\begin{proof}
For each sector and almost every $K_n$, Cauchy--Schwarz gives
\[
 \Big|\int\overline{f(l)}u_{n+1}(l,K_n)\,dl\Big|^2
 \leq\|f\|_2^2\int|u_{n+1}(l,K_n)|^2\,dl.
\]
Fubini/Tonelli then gives its $L^2(K_n)$ norm and, after multiplication by $n+1$ and summation, the first bound. Section measurability is MS02. The Bochner integral formulation in FC03 agrees with this formula: apply each continuous coordinate projection and use the scalar Fubini identity, first for rectangle tensors and then by the displayed norm bound.

For adjointness, pair the $(m+1)$-particle creation component with $u_{m+1}$. Its $m+1$ summands have identical integrals after a permutation of coordinates, because both $w_m$ and $u_{m+1}$ are symmetric. Each cross product is integrable by Cauchy--Schwarz on the product space, with norms $\|f\|\|w_m\|$ and $\|u_{m+1}\|$. The factor $(m+1)/\sqrt{m+1}$ is $\sqrt{m+1}$, leaving precisely the pairing with $(a(f)u)_m$. Summing the finitely many sectors proves the identity.

For the creation norm, fix $u_m$ and write
$F_i(k_1,\ldots,k_{m+1})=f(k_i)u_m(K_{\widehat i})$.
The squared norm of the output is
$(m+1)^{-1}\sum_{i,j}\int\overline{F_i}F_j$.
Each of the $m+1$ diagonal terms equals $\|f\|^2\|u_m\|^2$. For $i\ne j$, permutation symmetry makes the integral the same as for $i=1,j=2$. Write the other $m-1$ variables as $L$. That integral is
\[
 \int dL\int dk_1\,dk_2\,
 \overline{f(k_1)}f(k_2)
 \overline{u_m(k_2,L)}u_m(k_1,L)
 =\int dL\,\Big|\int\overline{f(k)}u_m(k,L)\,dk\Big|^2.
\]
The inner square is integrable by the first Cauchy--Schwarz estimate. All rearrangements also follow by absolute integrability of $\overline{F_i}F_j$, which is at most $\|F_i\|_2\|F_j\|_2$ in $L^1$. There are $(m+1)m$ ordered unequal pairs. Dividing by $m+1$ leaves $m$ times the last integral, exactly the squared norm of $a(f)$ on $u_m$. For $m=0$ the unequal-pair sum is empty and the same conclusion holds with zero annihilation term. Orthogonality of distinct output sectors permits summing over $m$, proving the claimed norm identity.

These estimates give bounded maps $a(f),a^*(f):V_N\to\F$, of norms at most $\|f\|$. To define the extension, take $u_j$ from FC01 converging in $V_N$; the images are Cauchy by the norm bounds, their limits are independent of the sequence, and linearity passes to the limits. Both sides of the adjoint relation and the norm equality are continuous in this convergence, so they hold on $V_N$. The integral of $\|\overline f\mathcal Au\|$ is bounded by $\|f\|\|\mathcal Au\|_2$, proving the Bochner integral exists for the extended annihilation operator as well.
\end{proof}


\node{FC03}{Smeared operators and their adjoints}{FC02, FC06, GW04}
\begin{lemma}\label{lem:smearedfield}
For $f\in L^2(\R^3)$ and $u\in V_N$,
\[
 a(f)u=\int\overline{f(k)}(\mathcal A u)(k)\dd k,
 \qquad \norm{a(f)u}\leq\norm f_2\norm{N^{1/2}u}.
\]
The Bochner integral exists. For $u,w\in V_N$,
\[
 \ip{a^*(f)w}{u}=\ip w{a(f)u},\qquad
 \norm{a^*(f)u}^2=\norm{a(f)u}^2+\norm f_2^2\norm u^2.
\]
The latter equality uses the usual creation/annihilation normalization. In particular
$\norm{a^*(f)u}\leq\norm f_2\norm{(N+1)^{1/2}u}$.
\end{lemma}
\begin{proof}
FC06 proves the sector formula, the full finite-sum adjoint calculation, the equal-index and unequal-index contributions to the creation norm, and their graph-norm extensions. For the integral description use FC02 and Cauchy--Schwarz:
$\int\|\overline{f(k)}\mathcal Au(k)\|dk\leq\|f\|_2\|\mathcal Au\|_2$.
Thus the Bochner integral exists. Its coordinates are the annihilation sector integrals by MS02 and the continuity of coordinate projection; MS03 identifies the vector with $a(f)u$. The norm estimates and adjoint identity are precisely those already proved in FC06, so no additional commutator identity is inferred from an unstated CCR theorem.
\end{proof}
The inspected GW normalization is exactly this one. GW04 fixes the restrictions of its maximal operators and proves their agreement with these integral formulas; no factorial-weight adapter is needed.

\paragraph{A safe interpretation of a varying creation field.}
For a field $G\in L^2(\R^3;\F)$, define the continuous conjugate-linear functional
\[
 \Gamma_G(w)=\int\ip{(\mathcal A w)(k)}{G(k)}\dd k,
 \qquad \abs{\Gamma_G(w)}\leq\norm{N^{1/2}w}\norm G_2.
\]
It is a form-dual load, not in general a physical vector. If notation such as
$\int a_k^*G(k)\dd k$ is ever desired, it means precisely this functional. A changing vector $G(k)$ cannot be treated as a single ordinary smeared creation operator. The kernel proof below does not even require this extra definition: it uses scalar pairings and ordinary $a^*(f)$ at finite cutoff.

\node{FC04}{Energy control for cutoff annihilation fields}{FC01, FC02, FD02, MS02, MS03, ML-PROD, EN02}
\begin{lemma}\label{lem:fieldenergy}
Let $u\in\Dom H_0(P)$ and put $F_P=H_0(P)=(\Pf-P)^2+N$. For almost every $k$, $(\mathcal A u)(k)\in\Q_{P-k}$, and
\begin{equation}
 \int\norm{(\mathcal A u)(k)}_{\Q_{P-k}}^2\dd k
 =\sum_{m\geq1}\int m\bigl(\abs{K_m-P}^2+m\bigr)\abs{u_m}^2
 \leq\norm{F_Pu}^2.\label{eq:fieldenergyexact}
\end{equation}
On every bounded set $K$ of external momenta this implies integrability with values in the fixed graph space $\Q_0$. The corresponding field map is continuous from the free operator graph norm to $L^2(K;\Q_0)$.
\end{lemma}
\begin{proof}
On the output sector $m-1$, the weight of $\Q_{P-k}$ is
$1+\abs{(K_m-k)-(P-k)}^2+(m-1)=\abs{K_m-P}^2+m$.
Insert the factor $m$ from \eqref{eq:fielddefinition}, integrate, and use Tonelli. The inequality is the pointwise numerical bound
$m(\abs{K_m-P}^2+m)\leq(\abs{K_m-P}^2+m)^2$.
For fixed bounded $K$, the shifted graph norms are uniformly equivalent to $\Q_0$ by the inequalities in Section~\ref{sec:gw}. To make the energy-valued measurability explicit, prove the result first for $J_{j,j}u$ and then use the displayed estimate on differences. The resulting limit in $L^2(K;\Q_0)$ has the same physical representative as $\mathcal A u$ by the continuous embedding. Exhaust $\R^3$ by balls.
\end{proof}
\paragraph{A domain implication that must not be used.}
Membership $u\in\Q$ alone does \emph{not} imply that $a_ku\in\Q$ a.e. For example take a normalized even compactly supported orbital $\chi$ with finite nonzero second moment, and set
$u_n=n^{-3/2}\chi^{\otimes n}$ for $n\geq1$.
The $N$-form expectation is $\sum n^{-2}<\infty$; the kinetic form expectation is a constant times the same sum, because the mean momentum of $\chi$ is zero. Hence $u\in\Q$. But for every $k$ with $\chi(k)\ne0$, the number-form expectation of its annihilation slice contains
\[
 \abs{\chi(k)}^2\sum_{n\geq2}(n-1)n\,n^{-3}=\infty.
\]
Thus the stronger field-energy fact is used only for \emph{cutoff} resolvent images, after their free operator-domain membership has been proved. The singular limit is taken in $L^2$ fields, not in this stronger graph norm.

\node{FC07}{Field commutations and the compressed free shift}{FC01, FC02, FC06, FD02}
\begin{lemma}[All commutations used later]\label{lem:fieldCCR}
For $u$ with finitely many sectors and $f\in L^2$, the following are equalities of $L^2(\R^3;\F)$ fields:
\begin{align*}
 \mathcal A(a^*(f)u)(k)&=a^*(f)(\mathcal Au(k))+f(k)u,\\
 \mathcal A(a(f)u)(k)&=a(f)(\mathcal Au(k)).
\end{align*}
For bounded-coordinate core vectors, as equalities of locally square-integrable fields with the indicated polynomial multipliers,
\[
 \mathcal A N u=(N+1)\mathcal Au,\qquad
 \mathcal A(P_{\rm f})_j u=((P_{\rm f})_j+k_j)\mathcal Au.
\]
If $f$ has bounded support, the smeared core commutators are
\[
 [(P_{\rm f})_j,a^*(f)]=a^*(k_jf),\qquad
 [(P_{\rm f})_j,a(f)]=-a(k_jf).
\]
For $n=1,2$, put
$L_{n,p,\Lambda}=\Pi^{\geq n}(H_{g,\Lambda}(-p)-e+1-n)\Pi^{\geq n}$.
On the core in the $n$-sector complement,
\begin{equation}
 \mathcal A L_{n,p,\Lambda}u(k)
 =L_{n-1,p+k,\Lambda}\mathcal Au(k)+v_{g,\Lambda}(k)u.
 \label{eq:projectedcorePT}
\end{equation}
The field on the right is extended by zero below sector $n-1$.
\end{lemma}
\begin{proof}
For the first formula examine the output sector $r$. The left side is $\sqrt{r+1}$ times the $(r+1)$-sector creation formula evaluated at $(k,K_r)$. The normalization factors cancel. The summand using the distinguished coordinate $k$ is $f(k)u_r(K_r)$. Each of the remaining $r$ summands is $f(k_i)u_r(k,K_{\widehat i})$. The latter sum is exactly the $r$-sector creation formula applied to the field: its factor $1/\sqrt r$ cancels the factor $\sqrt r$ in $(\mathcal Au(k))_{r-1}$. For $r=0$ this sum is empty. For the second identity, both sides in sector $r$ are
\[
 \sqrt{(r+1)(r+2)}\int\overline{f(l)}u_{r+2}(k,l,K_r)\,dl;
\]
interchanging $k,l$ uses the input symmetry. These integral expressions are defined for almost every $(k,K_r)$ by Fubini. Finite-sector vectors lie in every number-power domain. The bounds of FC06, integrated in $k$, show that the field terms just calculated have finite $L^2$ norms; for example, on an input with at most $M$ particles, $a^*(f)$ on each slice has norm at most $\sqrt M\|f\|$ times its slice norm. Thus the equality is in $L^2$, not just a formal coordinate statement.

Removing one momentum changes total momentum from $k+K_r$ to $K_r$ and number from $r+1$ to $r$, proving the free-field identities directly. For the smeared creation commutator the difference between total output momentum and the momentum of the omitted-coordinate input is $k_i$ in the $i$th summand. This is exactly creation with $k_jf$. For annihilation, total input momentum is $l+K_r$, so subtracting the two compositions leaves $-l_j\overline{f(l)}$, which is annihilation with $k_jf$ and a minus sign. Bounded support of $f$ ensures $k_jf\in L^2$; all vectors stay in FC01's core, so both compositions are defined.

To include projections, check coordinates: $(\mathcal A\Pi^{\geq n}u)_r$ is zero for $r<n-1$ and otherwise is $(\mathcal Au)_r$. Thus $\mathcal A\Pi^{\geq n}=\Pi^{\geq n-1}\mathcal A$. Applying this to the preceding field identities, the interaction part is
$\Pi^{\geq n-1}\Phi(v_\Lambda)\mathcal Au+v_\Lambda(k)u$; the source term survives because $u$ already has at least $n$ particles. The free part becomes
$(P_{\rm f}+p+k)^2+(N+1)+1-n-e$, which equals
$(P_{\rm f}+p+k)^2+N+1-(n-1)-e$.
This proves the exact shift in \eqref{eq:projectedcorePT}, including its constant and projection. There is no omitted boundary sector: the projection after annihilation is precisely the one just computed.
\end{proof}

\node{FC08}{Measurable bounded operator families acting on fields}{MS02, MS03, ML-MEAS, ML-SIMPLE}
\begin{lemma}\label{lem:operatorfield}
Let $K\subset\R^3$ be bounded and measurable. Let $T(k):V\to W$ be operator-norm continuous on a set containing $K$, with $\sup_K\|T(k)\|\leq M$. For every $F\in L^2(K;V)$ the field $k\mapsto T(k)F(k)$ is strongly measurable, defines a class in $L^2(K;W)$ independent of the representative of $F$, and has norm at most $M\|F\|_2$. The same statements hold for a measurable family which is a pointwise operator-norm limit of such continuous families with a common bound.
\end{lemma}
\begin{proof}
For fixed $v$, the map $k\mapsto T(k)v$ is continuous on its parameter set. To see strong measurability without assuming a separable operator space, partition $[-j,j]^3$ into finitely many half-open cubes of diameter at most $2^{-j}$. In every cube meeting the parameter set choose one point of that intersection, and use the value $T(k_{j,i})v$ on the intersection with that cube; use zero outside the box. These are finite simple functions. At each fixed parameter point, the chosen point in its cube tends to it as $j\to\infty$, so continuity gives pointwise convergence. Thus the vector field is strongly measurable by the defining simple approximation criterion. Only the intersections with the measurable set $K$ enter the integrals. A simple field $F=\sum_j\mathbf1_{E_j}v_j$ gives a finite sum of these fields. A strongly measurable representative is an almost-everywhere pointwise limit of simple fields (ML-MEAS). Boundedness of $T(k)$ transfers that pointwise limit to the images, establishing strong measurability. The estimate $\|T(k)F(k)\|^2\leq M^2\|F(k)\|^2$ proves square integrability and the norm bound. Equal representatives have equal images outside the same exceptional set. For the last assertion take the pointwise limit of the already strongly measurable vector-valued maps and use the same bound. In particular continuous energy-space representatives and physical resolvents may be composed with $L^2$ fields without choosing pointwise values for equivalence classes globally.
\end{proof}

\node{FC09}{A common weak equation from an $L^2$ residual}{FC08, ML-SIMPLE, ML-L2, ML-LP}
\begin{lemma}\label{lem:residualzero}
Let $V$ be a complex Hilbert space, $K$ have finite measure, and $r\in L^2(K;V)$. Let $\mathscr D\subset V$ be a dense linear subspace. Suppose
\[
 \int_K\overline{\chi(k)}\langle w,r(k)\rangle\,dk=0
 \quad(w\in\mathscr D,\ \chi\in L^2(K)).
\]
Then $r=0$ as a Bochner $L^2$ vector, hence $r(k)=0$ outside one null set. No separability or parameter-independent choice of a countable test core is required.
\end{lemma}
\begin{proof}
For fixed $\chi$, the absolute value of the integral with test $w-w'$ is at most $\|\chi\|_2\|w-w'\|\|r\|_2$. Approximation of any $w\in V$ by $\mathscr D$ extends the equality to all constant vector tests. A finite simple field $s=\sum_i\mathbf1_{E_i}w_i$ has
$\langle s,r\rangle_{L^2}=\sum_i\int\mathbf1_{E_i}\langle w_i,r\rangle=0$; the indicators are in $L^2$ because $K$ has finite measure. Simple fields are dense by ML-SIMPLE. Choose $s_j\to r$ in $L^2$. Cauchy--Schwarz gives
$|\langle r,r\rangle-\langle s_j,r\rangle|\leq\|r-s_j\|\|r\|\to0$.
Thus $\|r\|^2=0$. The equivalence between zero $L^2$ norm and almost-everywhere zero is ML-LP. This is a vector equation outside one exceptional set, so at every remaining point it can be paired with all $w\in V$ simultaneously.
\end{proof}

\node{FC05}{The weak-equation extraction interface used in pull-through}{FC08, FC09, BA01}
Let $V$ be a fixed energy Hilbert space, $j:V\to H$ its embedding, $T(k)$ a bounded continuous family on a bounded parameter set, $U\in L^2(K;V)$ and $G\in L^2(K;H)$. The residual $r(k)=T(k)U(k)-j^*G(k)$ belongs to $L^2(K;V)$ by FC08. If its integrals against all scalar $L^2$ tests and a dense linear space of constant vector tests vanish, FC09 proves $T(k)U(k)=j^*G(k)$ as a vector identity outside one null set. Pairing this identity with any $w\in V$ then yields the weak equation for all tests at once. This is the precise common-null-set interface used in PT01; it does not require a separability argument.\label{lem:localweakextraction}


\paragraph{Quantifiers for parameterized fields.}
A statement ``for every $p$, for almost every $k$'' must not be strengthened to a single exceptional set for all $p$. In particular it does not authorize the substitution $p=k$. The scalar kernel proof first integrates out its field variable for each fixed parameter and obtains an identity for \emph{every} parameter; only that scalar identity is subsequently composed with a momentum translation. This avoids a diagonal-null-set problem entirely.

\section{The bounded energy-space bridge to the full spectrum}\label{sec:abstract}
This section fixes the mathematical functional-analysis interface and gives its proofs. GW06 and GW07 prove its exact realization in the retained source. All operator matrices in this section are bounded operators on \emph{energy spaces}; a singular coupling is not silently promoted to a bounded operator on the physical Hilbert space.

\node{FA01}{Energy-space inversion}{FA00, BA01}
\begin{lemma}\label{lem:forminverse}
If $T_\lambda$ is boundedly invertible on $V$, then
\[
 R_\lambda=jT_\lambda^{-1}j^*
\]
is a bounded two-sided inverse of $A-\lambda$. If $T_\lambda u=0$ for a nonzero $u\in V$, then $ju$ is a nonzero $\lambda$-eigenvector of $A$.
\end{lemma}
\begin{proof}
For $f\in H$ set $u=T_\lambda^{-1}j^*f$. Then, for every $w\in V$,
\[
 a_\lambda(w,u)=\ip{w}{j^*f}_V=\ip{jw}{f}_H.
\]
The representation characterization gives $ju\in\Dom A$ and $(A-\lambda)ju=f$. Conversely, for $ju\in\Dom A$, the same characterization gives $T_\lambda u=j^*(A-\lambda)ju$. Apply $T_\lambda^{-1}$ and $j$. The operator norm is at most $\norm j^2\norm{T_\lambda^{-1}}$. The kernel statement follows from the same characterization with $f=0$ and the injectivity of $j$.
\end{proof}

\node{FA02}{Coercivity gives an energy-space inverse}{BA01, BA02, GW-Coercive}
\begin{lemma}\label{lem:coercive}
If a bounded self-adjoint $T$ on a Hilbert space satisfies $\ip u{Tu}\geq\alpha\norm u^2$ for $\alpha>0$, then $T$ has a bounded inverse and $\norm{T^{-1}}\leq\alpha^{-1}$.
\end{lemma}
\paragraph{Implementation contract: direct GW reuse.}
Use the complete complex Hilbert energy carrier already fixed by the predecessors. The real inequality in the statement means
$\alpha\|u\|^2\leq\Re\langle u,Tu\rangle$, and bounded self-adjointness means $T^*=T$. Do not reconstruct coercivity, range density, a Neumann series, or a complex Lax--Milgram theorem for this node.
\begin{proof}
Instantiate \path{GW.Operators.coerciveInverse} from GW-Coercive with $B=T$, $c=\alpha$, the displayed positivity hypothesis, and $T^*=T$. It gives a bounded complex-linear $R$. The retained lemmas \path{coerciveInverse_left} and \path{coerciveInverse_right} give $RT=TR=I$, and \path{coerciveInverse_norm_le} gives $\|R\|\leq\alpha^{-1}$. Thus $R=T^{-1}$ with exactly the claimed bound. If BA02 already supplies a different named inverse $S$, the identity $S=S(TR)=(ST)R=R$ identifies them; no second inverse theory is needed. Completeness is the energy-space completeness, not a norm inherited from an unclosed physical domain.
\end{proof}

\node{FA03}{Two estimates imply graph-norm coercivity}{ML-REAL}
\begin{lemma}\label{lem:interpolation}
Suppose a Hermitian form $s$ on $\Q_1$ satisfies
\[
 s[f]\geq\varepsilon\norm f^2,
 \qquad s[f]\geq a\norm{\abs k f}^2-b\norm f^2
\]
with $\varepsilon,a>0$ and $b\geq0$. Then
\[
 s[f]\geq\theta a\norm{\abs k f}^2+\frac\varepsilon2\norm f^2,
 \qquad \theta=\frac\varepsilon{2(b+\varepsilon)}.
\]
\end{lemma}
\begin{proof}
Take $\theta$ times the second inequality plus $1-\theta$ times the first. The coefficient of $\norm f^2$ is $\varepsilon-\theta(b+\varepsilon)=\varepsilon/2$.
\end{proof}
This is the missing domain step in a proof that merely observes $s\geq\varepsilon I$. For a singular form, Hilbert-norm positivity alone is not an energy-space inverse statement.

\node{FA04}{Schur elimination on energy spaces}{FA01, BA02}
\begin{lemma}\label{lem:schur}
Let $V=V_0\oplus V_1$ and let the bounded Hermitian representative of a form be
\[
 T=\begin{pmatrix}A&B^*\\B&D\end{pmatrix},\qquad D^{-1}\in\mathcal B(V_1).
\]
Set $S=A-B^*D^{-1}B$. Then
\begin{equation}
 T=\begin{pmatrix}I&B^*D^{-1}\\0&I\end{pmatrix}
 \begin{pmatrix}S&0\\0&D\end{pmatrix}
 \begin{pmatrix}I&0\\D^{-1}B&I\end{pmatrix}.\label{eq:schurfactor}
\end{equation}
If $S$ is boundedly invertible, so is $T$. Its inverse is the bounded map obtained by solving
\[
 x=S^{-1}(f_0-B^*D^{-1}f_1),\qquad y=D^{-1}(f_1-Bx).
\]
Its kernel is $\{(x,-D^{-1}Bx):Sx=0\}$. If $D\geq0$ and $S\geq0$, then $T\geq0$.
\end{lemma}
\begin{proof}
Multiply the three matrices, or expand
\[
 \ip{(x,y)}{T(x,y)}=\ip x{Sx}+\ip{y+D^{-1}Bx}{D(y+D^{-1}Bx)}.
\]
The triangular factors have bounded triangular inverses, obtained by changing the sign of the off-diagonal entry. All claims follow; use Lemma~\ref{lem:forminverse} only after this bounded-energy calculation is complete.
\end{proof}

The associated operator is not an external assumption: FA00 constructs it and proves its self-adjointness and uniqueness. FA01 and FA04 now provide its physical inverse from an invertible bounded energy representative. In particular the complement block need not be positive in a Schur calculation; invertibility is enough. Positivity is used only for those separate claims that explicitly require it.


\node{FA05}{Kernel obstruction, uniqueness of inverses, and the real spectrum}{FA00, FA01, FD02, GW-Resolvent, DF01}
\begin{lemma}\label{lem:spectraldefinitions}
For the domain-and-action operator constructed in FA00, the following facts require no spectral theorem.
\begin{enumerate}
\item Two bounded two-sided inverses at the same parameter agree.
\item A nonzero eigenvector at a real parameter forces that parameter to lie in $\sigR(A)$.
\item If $e\in\sigR(A)$ and every real $\lambda<e$ has a bounded two-sided inverse, then $e=\inf\sigR(A)$.
\item Defining the complex resolvent with the same two-sided identities using $A-z$, its intersection with the real line is exactly the real resolvent. Consequently a real-spectrum equality below a real threshold is the desired complex-spectrum intersection with that real interval.
\end{enumerate}
\end{lemma}
\begin{proof}
If $R,S$ are inverses, then $Sf$ is in the operator domain and
$Rf=R(A-\lambda)Sf=Sf$ for all $f$. If $(A-\lambda)u=0$ and an inverse existed, its left identity would give $u=R0=0$. For the infimum assertion, $e$ is a lower bound because no smaller real number is in the spectrum, and every lower bound is at most $e$ because $e$ is a member. Thus $e$ is the greatest lower bound; the set is nonempty and bounded below, so its real infimum is $e$. For the last assertion substitute the scalar embedding $\R\hookrightarrow\C$ in the two definitions: the same bounded complex-linear map satisfies the same domain and algebraic identities. Complements within the real line and intersections with its intervals preserve this equality. No bounded-operator spectrum API is applied directly to an unbounded Hamiltonian.
\end{proof}

\paragraph{The domain object which this proof specifies.}
Use the existing GW \texttt{LinearPMap} representation with its specified domain submodule. No new partial-operator or adjoint library is to be defined. The domain predicate is the one supplied by \path{formOperator_domain_iff}, and the inverse predicate is the existing \texttt{IsResolvent}. GW06 supplies the exact form record; GW07 proves the spectrum correspondence. FA00 is a mathematical explanation of the retained representation proof, not a parallel implementation.


\subsection{The exact partial-operator and resolvent interface}
\node{GW07}{Interpreting the spectrum by GW's actual inverse predicate}{FA01, FA05, GW-Form, GW-Resolvent}
\begin{lemma}\label{lem:gwresolvent}
For an associated operator $H=\texttt{GW.Operators.formOperator}\ d$ define
\[
 \rho_{\C}(H)=\{z\in\C:\exists R\in\mathcal B(\F),\
       \texttt{GW.Operators.IsResolvent}\ H\ z\ R\},\qquad
 \sigma_{\C}(H)=\C\setminus\rho_{\C}(H).
\]
Then the real resolvent and spectrum used in FA05 are exactly the preimages of these sets under $\R\hookrightarrow\C$. Every inverse supplied by FA01 or FA04 satisfies the three fields of this GW predicate on the \emph{actual} partial-operator domain. A weak kernel vector gives an eigenvector on that same domain.
\end{lemma}
\begin{proof}
GW's predicate has the following exact three requirements for a bounded map $R$: (i) $Rf\in\Dom H$ for every physical $f$; (ii) $(H-z)Rf=f$ for every $f$; (iii) $R(H-z)u=u$ for every domain vector $u$. In the source the fields are called respectively \texttt{mem\_domain}, \texttt{left\_inv}, and \texttt{right\_inv}. These names are recorded to avoid reversing the two identities during translation.

For real $\lambda$, let $T_\lambda$ represent the form $h-\lambda\langle j\cdot,j\cdot\rangle$ on the fixed coordinate space and suppose it is boundedly invertible. Put $u=T_\lambda^{-1}j^*f$ and $Rf=ju$. The weak equation is
$h(v,u)=\langle jv,f+\lambda ju\rangle$ for every $v$.
Apply the existing \path{formOperator_domain_of_equation} to exactly this right-hand-side vector. It produces the first field and the equality $H(Rf)=f+\lambda Rf$, which is the second field after subtracting $\lambda Rf$.

For any $x\in\Dom H$, \path{operatorFormCoordinate} gives $u_x$ with $ju_x=x$, and \path{formOperator_form_identity} gives
$T_\lambda u_x=j^*(Hx-\lambda x)$.
Applying $T_\lambda^{-1}$ and then $j$ proves the third field. The equalities are independent of membership witnesses for $x$ by subtype equality. If $T_\lambda u=0$, the same domain theorem with source $\lambda ju$ gives $H(ju)=\lambda ju$. Injectivity of $j$ preserves nonzeroness. These are precisely the arguments of FA01, now with their existing GW interface fixed.

Two inverse witnesses coincide by \texttt{isResolvent\_unique}, so no choice in the construction alters the resolvent set. The real-set identification follows by unfolding the displayed definition at $z=\lambda+0i$; it does not require a bounded-operator spectrum API. For completeness, self-adjointness of \texttt{formOperator} and the existing \texttt{resolvent\_isResolvent} give every $z$ with nonzero imaginary part in the resolvent. Thus this is the usual full spectrum of the self-adjoint partial operator. The spectral theorem, a spectral measure, and essential spectrum are not needed for any of these assertions.
\end{proof}

\section{Three-dimensional measure estimates with no exact constants}\label{sec:integrals}
The analysis needs finite uniform constants, not the exact convolution value $\pi^3/\abs P$. This section specifies sufficient elementary estimates. It can be implemented using dyadic shells and Lebesgue scaling, avoiding a new spherical-coordinate integration framework.

\node{AN00}{Translations, scalar dilations and positive finite ball volume}{MS01, ML-VOLUME, ML-PROD, ML-BOCHNER, GW-Model}
\begin{lemma}\label{lem:affinemeasure}
On $\R^3$ with product Lebesgue measure, translations and coordinate reflections preserve measure. For $a\ne0$, dilation by $a$ scales the measure of a measurable set by $|a|^3$. Thus for any nonnegative measurable $f$, and also for any integrable complex $f$,
\[
 \int f(ax+b)\,dx=|a|^{-3}\int f(y)\,dy.
\]
The unit Euclidean ball has finite positive volume $v_3$, every radius-$R$ ball has volume $v_3R^3$, and every singleton has zero measure.
\end{lemma}
\paragraph{Implementation contract: reuse the geometric clauses.}
On the prescribed \path{GW.Space} with its specified volume, use \path{GW.Model.unitBallVolume_pos}, \path{GW.Model.spaceBallVolume}, \path{GW.Model.spaceSphereVolume}, and \path{GW.Model.spaceClosedBallVolume} for the applicable volume clauses, including their radius conditions. The proof below explains these facts but is not a request to rebuild them. For translations, reflections, and dilations, first instantiate the pinned measure-preserving/Haar change-of-variables interfaces on this same carrier; if a direct match is not established, the supplied rectangle argument remains the proof route. Do not assert a new measure identification from notation alone.
\begin{proof}
On a product of three bounded intervals, volume is the product of their lengths (ML-VOLUME). Translation leaves these lengths unchanged; inverse dilation divides each by $|a|$. Reversing an interval swaps its endpoints but leaves its length. These calculations prove the pushforward identity on half-open bounded rectangles. Such rectangles form an intersection-stable generating family. To extend to all Borel sets, restrict both measures to a bounded rectangle and apply the Dynkin-system argument in MS01 to the class of subsets on which the two finite measures agree: complements use subtraction from the common finite total measure and disjoint unions use countable additivity. Exhaustion by rectangles then gives the equality of the sigma-finite measures on the full space. This also proves translation and reflection invariance. For indicator functions the integral formula is the just-proved measure identity. Positive simple functions follow by finite addition, and nonnegative measurable functions by monotone convergence of their simple approximants. Integrable real and imaginary parts, or an $L^1$ simple approximation, give the complex formula. Thus every later substitution comes with a measure identity, not an assumed coordinate Jacobian formula.

The cube $(-1/\sqrt3,1/\sqrt3)^3$ lies inside the unit Euclidean ball and has positive volume; the unit ball lies inside $[-1,1]^3$ of finite volume. This proves $0<v_3<\infty$. Apply dilation to obtain the radius formula, and translation for other centers. A singleton is contained in cubes with side length $2/j$, whose measures are $8/j^3$; monotonicity of measure and passage to the infimum give zero. The arguments apply to the fixed identification with $\mathrm{EuclideanSpace}\,\R\,(\mathrm{Fin}\,3)$; its norm is the Euclidean norm used in the cube inclusions.
\end{proof}

\node{AN01}{Power integrals and infrared nonintegrability}{AN00, ML-SERIES, ML-PROD, GW-Model}
\begin{lemma}\label{lem:powers}
For $0<a<3$ and $R>0$, $\int_{\abs x<R}\abs x^{-a}\dd x\leq C_aR^{3-a}$. For $a>3$, $\int_{\abs x>R}\abs x^{-a}\dd x\leq C_aR^{3-a}$. Moreover $\abs x^{-4}$ is not integrable on any neighborhood of zero.
\end{lemma}
\begin{proof}
Let $v_3$ be the volume of the unit ball. It is finite and strictly positive, and volume scales as $R^3$. On the shells $2^{-n-1}R<\abs x\leq2^{-n}R$, dominate the integral by
\[
 v_3(2^{-n}R)^3(2^{-n-1}R)^{-a}
 =2^av_3R^{3-a}2^{(a-3)n}.
\]
Sum the geometric series. For the exterior use shells $2^nR<\abs x\leq2^{n+1}R$ instead. For nonintegrability use the shell $B(0,2^{-n})\setminus B(0,2^{-n-1})$, whose volume is the difference of the two ball volumes; its integral is at least
\[
 (2^{-n})^{-4}\,v_3(2^{-3n}-2^{-3n-3})
 =\tfrac78v_3\,2^n.
\]
Finite unions already have arbitrarily large integrals. Scaling gives every neighborhood. Boundaries can be assigned to neighboring shells for the upper bounds. The lower bound uses a difference of nested open balls, so it also needs no sphere-measure calculation.
\end{proof}

\node{AN02}{Uniform shifted source norm}{AN00, AN01}
\begin{lemma}\label{lem:shiftedsource}
There is $C_v$ such that
\[
 \sup_{p\in\R^3}\int\frac{\dd k}{\abs k^2(1+\abs{k+p}^2)}\leq C_v.
\]
For fixed $R$, the same integral restricted to $\abs k>\Lambda$ is at most $C/\Lambda$ when $\abs p\leq R$ and $\Lambda\geq\max\{1,2R\}$.
\end{lemma}
\begin{proof}
For $\abs p\leq2$, split at $\abs k=4$: the inner integral is bounded by $\int_{\abs k\leq4}\abs k^{-2}$, and the outer integrand by $4\abs k^{-4}$. For $r=\abs p>2$, split into $\abs k<r/2$, $\abs{k+p}<r/2$, the remaining part of $\abs k\leq2r$, and $\abs k>2r$. On the first region the bound is $4r^{-2}\abs k^{-2}$, giving $C/r$. On the second, translate by $p$ and use
\[
 4r^{-2}\int_{\abs y<r/2}(1+\abs y^2)^{-1}\dd y\leq C/r.
\]
Here split at $\abs y=1$ and use Lemma~\ref{lem:powers}. On the third region both distances are at least $r/2$, giving $C r^{-4}$ on a set of volume at most $C r^3$. On the last use $4\abs k^{-4}$. The tail statement follows from $\abs{k+p}\geq\abs k/2$.
\end{proof}

\node{LD01}{Weighted source functionals and their exact convergence norm}{AN02, EN01, EN02, BA01, GW01, ML-BOCHNER, GW-Domination}
\begin{lemma}\label{lem:weightedloads}
Let $V$ be a one-particle energy space with norm $\|wf\|_2$, where $w\geq1$ is a finite positive measurable weight. For measurable $s$ with $s/w\in L^2$, define the conjugate-linear load
\[
 F_s(f)=\int\overline{f(k)}s(k)\,dk.
\]
This integral is absolutely convergent, its dual norm is exactly $\|s/w\|_2$, and its Riesz coordinate on the weighted-coordinate carrier is $s/w$. If $s_j/w\to s/w$ in $L^2$, then $F_{s_j}\to F_s$ in dual norm. It suffices that $s_j/w\to s/w$ pointwise a.e. and
$|(s_j-s)/w|^2\leq M$ with $M$ integrable, in which case dominated convergence proves this norm limit.
\end{lemma}
\begin{proof}
Write $q=wf\in L^2$. Then $\overline f s=\overline q(s/w)$ since $w$ is positive real. Cauchy--Schwarz gives the $L^1$ bound $\|q\|_2\|s/w\|_2$ and the pairing $F_s(f)=\langle q,s/w\rangle$. This proves the asserted Riesz vector and upper norm bound. If $s/w\ne0$, choose $q=(s/w)/\|s/w\|_2$ to obtain equality of norms; if it is zero both sides vanish. The difference of two such functionals has source $s_j-s$, so the exact norm statement proves the convergence assertion. For the final criterion the nonnegative integrands $|(s_j-s)/w|^2$ are measurable, tend to zero a.e., and are dominated by $M$; dominated convergence gives convergence of their integrals to zero, hence of the square-root norms. Changes on null sets do not alter any functional.

In particular $s=v_g$, $w(k)=\sqrt{1+|k+p|^2}$ satisfies the hypothesis uniformly in $p$ by AN02. The cutoff error is the square root of the weighted tail integral there. On a Fock energy space, the source load uses only the one-particle projection and the map $j_1$ in GW01. This projection is bounded by EN02. Thus no occurrence of $Y_pv_g$ treats $v_g$ as a physical Hilbert vector. A bounded measurable factor $b(k)$ may multiply the source since $|b s/w|\leq\|b\|_\infty|s/w|$; this is the precise source operation used in SE02.
\end{proof}

\node{TR01}{Momentum translation as a specific form-domain isomorphism}{AN00, AN02, EN01, EN02, LD01, ML-LP}
\begin{lemma}\label{lem:translationadapter}
For $P\in X$, $\tau_Pf(q)=f(q+P)$ is a unitary on $L^2(X)$. It maps the isolated form domain with norm $\int(1+|k-P|^2)|f(k)|^2\,dk$ isometrically onto the domain with norm $\int(1+|q|^2)|f(q)|^2\,dq$. The inverse is $\tau_{-P}$. On the common physical domain,
\[
 \ell_g(\tau_{-P}f)=\int\frac{g}{|q+P|}f(q)\,dq.
\]
This integral is absolutely convergent and its functional norm in the unshifted isolated graph norm is bounded uniformly for $|P|\leq1/4$ (indeed for all $P$ by AN02 and a change of variables).
\end{lemma}
\begin{proof}
Translation preserves Lebesgue measure by AN00. Composition with it therefore preserves $L^2$ equivalence classes and their inner products. Substituting $k=q+P$ in the nonnegative weighted norm integral gives the stated isometry, with $|k-P|=|q|$; composing with the negative translation gives identity. The source equation is the change-of-variables formula for an absolutely integrable function. Its absolute integral is bounded by
\[
 |g|\left(\int\frac{dq}{|q+P|^2(1+|q|^2)}\right)^{1/2}
       \left(\int(1+|q|^2)|f(q)|^2dq\right)^{1/2}.
\]
The first integral is the shifted-source integral of AN02 after $k=q+P$. This proves integrability before making the substitution, and proves the uniform bound. On energy-coordinate copies the isometry is the map obtained by conjugating physical translation with the weights; the norm formula just proved specifies it uniquely. Conjugating a sesquilinear form means applying its inverse to both arguments, not conjugating a nonexistent everywhere-defined operator. In particular $O\geq0$ if and only if $\tau_PO\tau_P^{-1}\geq0$ on the translated domain.
\end{proof}

For $\omega(q)=\min\{\abs q,1\}$ and $P\in\R^3$ define
\[
 \mu_P(q)=\frac1{\abs q^2\abs{q+P}^2},\qquad
 w_P(q)=\frac1{\abs q\abs{q+P}}.
\]
Set their singular-point values to zero.

\node{AN03}{Moving singularities}{AN00, AN01}
\begin{lemma}\label{lem:singular}
Uniformly for $\abs P\leq1/4$,
\begin{align}
 \int\omega(q)(\omega(q)+\abs P)\mu_P(q)\dd q&\leq C_I,\label{eq:weightint}\\
 \int\omega(q)^3\abs{\mu_P(q)-\mu_0(q)}\dd q&\leq C_I\abs P.\label{eq:weightdiff}
\end{align}
For $P\ne0$,
\begin{equation}
 \int\mu_P(q)\dd q\leq C_I/\abs P.\label{eq:tint}
\end{equation}
\end{lemma}
\begin{proof}
For $P=0$, the first integrand is $\abs q^{-2}$ inside the unit ball and $\abs q^{-4}$ outside; the second vanishes. Let $\rho=\abs P>0$, $\nu_P=P/\rho$, and split at $\abs q=2\rho$.

On the inner region put $q=\rho x$. Since $2\rho\leq1/2$, its contribution to \eqref{eq:weightint} equals
\[
 \rho\int_{\abs x\leq2}\frac{\abs x(\abs x+1)}{\abs x^2\abs{x+\nu_P}^2}\dd x\leq C\rho.
\]
The last integral is uniformly bounded: partition into $\abs x<1/4$, $\abs{x+\nu_P}<1/4$, and the rest of the ball. In the first region the other denominator is at least $(3/4)^2$; in the second $\abs x\geq3/4$; in the rest both denominators are bounded below. Apply Lemma~\ref{lem:powers} to the remaining powers $1$ or $2$. The same partition shows
\[
 \int_{\abs q\leq2\rho}\omega(q)^3(\mu_P(q)+\mu_0(q))\dd q\leq C\rho^2.
\]
On the outer region,
\[
 \abs{q+P}\geq\abs q/2,\quad
 \mu_P\leq4\abs q^{-4},\quad
 \abs{\mu_P-\mu_0}\leq10\rho\abs q^{-5}.
\]
The last bound follows by subtracting the two fractions and bounding the numerator $\abs{2q\cdot P+\rho^2}\leq(5/2)\rho\abs q$. Also $\omega(q)+\rho\leq(3/2)\omega(q)$. Thus \eqref{eq:weightint} reduces to the already integrable $\omega^2\abs q^{-4}$. For \eqref{eq:weightdiff}, the remaining bound is $C\rho$ times the integral of $\abs q^{-2}$ between $2\rho$ and $1$, plus that of $\abs q^{-5}$ outside $1$. The inner contribution $C\rho^2$ is also at most $C\rho$.

Finally, scaling over all of $\R^3$ gives $\int\mu_P=\rho^{-1}\int\abs x^{-2}\abs{x+\nu_P}^{-2}\dd x$. The same two small balls control the local singularities, and on $\abs x>2$ the integrand is bounded by $4\abs x^{-4}$. This proves \eqref{eq:tint} without evaluating the integral.
\end{proof}

\paragraph{Consequences used later.}
The integrals $\int\omega^2\mu_P$ and $\int\omega^3\mu_P$ are uniformly bounded. For $g\ne0$, $t_0(q)=g/\abs q^2$ is not in $L^2$ by Lemma~\ref{lem:powers}, whereas $t_P=g w_P\in L^2$ for $P\ne0$. If a scalar $a$ satisfies $a t_0\in L^2$, then $a=0$: otherwise divide the $L^2$ function by $a$.

\section{Fiber forms, ultraviolet limits, and sector lower bounds}\label{sec:forms}
\node{FB01}{Finite-cutoff Lieb--Yamazaki inequalities, before any singular form exists}{FC01, FC03, FC07, AN01, BA03, GW05}
For a finite $\Lambda>0$ put $v_{g,\Lambda}(k)=g|k|^{-1}\mathbf1_{|k|\leq\Lambda}$, with value zero at $k=0$. AN01 gives $v_{g,\Lambda}\in L^2$. Define the cutoff Hermitian form on $\Q$ by
\[
 \Phi_\Lambda(w,u)=\langle w,a^*(v_{g,\Lambda})u+a(v_{g,\Lambda})u\rangle.
\]
Both field terms are defined on $\Q\subset V_N$, and FC03 bounds this form continuously on $\Q_p$ for every fixed $\Lambda,p$. At this node the symbol $\Phi(v_g)$ has not yet been defined.
\begin{lemma}\label{lem:lycutoff}
There is a finite constant $C_{\rm LY}$ independent of $\Lambda,g,p,\theta$ such that, for $0<\theta\leq1$,
\[
 |\Phi_\Lambda[u]|\leq\theta h_p^{\rm fr}[u]
       +C_{\rm LY}(g^2/\theta+g^4/\theta^3)\|u\|^2.
\]
For $|g|\leq\theta$ the remainder is at most $C_Lg^2\theta^{-1}\|u\|^2$, with one fixed $C_L$. For $1\leq L\leq\Lambda'$ the form with form factor
$g|k|^{-1}\mathbf1_{L<|k|\leq\Lambda'}$ satisfies, uniformly in $p$,
\begin{equation}
 |\Phi_{(L,\Lambda']}[u]|
 \leq C|g|L^{-1/2}\| (P_{\rm f}-p)u\|
                  \|(N+1)^{1/2}u\|.
 \label{eq:finite-tail-preyoung}
\end{equation}
\end{lemma}
\begin{proof}
Work first on FC01's invariant core. Choose $L\geq1$ and split at $\min(L,\Lambda)$. The low form factor has squared $L^2$ norm at most $Cg^2L$, by AN01 with exponent two. Smeared adjointness gives its diagonal form as $2\Re\langle u,a(v_{\rm low})u\rangle$. Therefore
\[
 |\Phi_{\rm low}[u]|
 \leq2\|v_{\rm low}\|\|u\|\|N^{1/2}u\|
 \leq\frac\theta2\|N^{1/2}u\|^2+Cg^2L\theta^{-1}\|u\|^2.
\]
The last inequality is $2ab\leq\delta a^2+\delta^{-1}b^2$ with $\delta=\theta/2$.

For the high part put
$w_j(k)=gk_j|k|^{-3}\mathbf1_{L<|k|\leq\Lambda}$ and
$C_j=a^*(w_j)-a(w_j)$. The functions $w_j$ and $k_iw_j$ are square integrable at finite cutoff. Thus FC07 proves, as genuine core-operator equalities,
\[
 \sum_j[(P_{\rm f})_j-p_j,C_j]
   =\sum_j\Phi(k_jw_j)=\Phi_{\rm high},
 \qquad \sum_j\|w_j\|_2^2
   =g^2\int_{L<|k|\leq\Lambda}|k|^{-4}dk\leq Cg^2/L.
\]
Here $\sum_jk_jw_j=g|k|^{-1}\mathbf1_{L<|k|\leq\Lambda}$ follows from $\sum_jk_j^2=|k|^2$. On the core, smeared adjointness gives $C_j^*=-C_j$ in pairings. For $D_j=(P_{\rm f})_j-p_j$, which is symmetric on its multiplication domain,
\[
 \langle u,[D_j,C_j]u\rangle
 =\langle D_ju,C_ju\rangle+\langle C_ju,D_ju\rangle
 =2\Re\langle D_ju,C_ju\rangle.
\]
Also $\|C_ju\|\leq2\|w_j\|\|(N+1)^{1/2}u\|$. Cauchy--Schwarz in the three components proves \eqref{eq:finite-tail-preyoung}. This calculation does not use an uncut smeared operator.

Young's inequality now bounds the high term by
\[
 \frac\theta2\|(P_{\rm f}-p)u\|^2
 +\frac{C g^2}{\theta L}(\|N^{1/2}u\|^2+\|u\|^2).
\]
With $C$ the high-term coefficient in the preceding inequality, set $C_0=\max\{1,2C\}$ and $L=C_0\max(1,g^2/\theta^2)$. Then $Cg^2/(\theta L)\leq\theta/2$. This ensures the combined coefficients of $N$ and kinetic energy are at most $\theta$. The low scalar term is bounded by
$CC_0(g^2/\theta+g^4/\theta^3)$; the high scalar term $Cg^2/(\theta L)$ is at most $Cg^2/\theta$. This proves the claimed remainder. If $|g|\leq\theta$, then $g^4/\theta^3\leq g^2/\theta$, giving the second form of the bound. If $L>\Lambda$, the high term is absent and the same upper bounds remain valid.

For a fixed finite cutoff the form and the free quadratic quantities are continuous in $\Q_p$. Approximate $u\in\Q$ by FC01's core in that norm. The squared norms and the cutoff form values converge, so the inequalities pass to all $u\in\Q$. The same argument extends the tail estimate. Every constant in the extension is the one already independent of $p,\Lambda$; the fixed-cutoff continuity is used only to identify the limiting left side.
\end{proof}

\node{FB02A}{Norm construction of the singular interaction and its limiting estimates}{FB01, BA03, FD02, EN02, GW-Mixed}
\begin{lemma}\label{lem:formlimit}\label{lem:ly}
There is a unique bounded Hermitian form $\Phi_g$ on the common form space $\Q$ which is the form-norm limit of $\Phi_\Lambda$ as $\Lambda\to\infty$. On every $\Q_p$, uniformly in $p$, it satisfies
\[
 |(\Phi_g-\Phi_\Lambda)(w,u)|
 \leq C|g|\Lambda^{-1/2}\|w\|_{\Q_p}\|u\|_{\Q_p}
 \quad(\Lambda\geq1).
\]
Both inequalities of FB01 hold with $\Phi_\Lambda$ replaced by $\Phi_g$. In particular
\begin{equation}
 |\Phi_g[u]|\leq\theta h_p^{\rm fr}[u]
                  +C_Lg^2\theta^{-1}\|u\|^2,
 \qquad |g|\leq\theta\leq1.
 \label{eq:ly}
\end{equation}
\end{lemma}
\begin{proof}
For finite $\Lambda'\geq\Lambda\geq1$, their difference is exactly the tail form in \eqref{eq:finite-tail-preyoung}. Since
$\|(P_{\rm f}-p)u\|\|(N+1)^{1/2}u\|\leq\|u\|_{\Q_p}^2/2$,
its diagonal absolute value is at most $C|g|\Lambda^{-1/2}\|u\|_{\Q_p}^2$. BA03 turns this into the displayed mixed estimate with another fixed factor two. First fix $p=0$ and let $\Lambda$ run over positive integers. This is a Cauchy sequence in the complete space of bounded Hermitian forms, so has a unique limit. Comparing an arbitrary real cutoff with integer cutoffs larger than it shows convergence of all real cutoffs with exactly the same tail estimate. The underlying domains agree for every $p$, with equivalent norms (FD02); thus the scalar limit for fixed $(w,u)$ is the same limit, not a new momentum-dependent form. Taking $\Lambda'\to\infty$ in the uniform mixed bound proves the estimate on $\Q_p$.

For fixed $u,p,\theta$, the diagonal cutoff inequalities have left sides converging to $|\Phi_g[u]|$ and right sides independent of the cutoff. Passing to the real limit proves both versions. This occurs strictly after the form has been constructed; no limiting-form estimate has been used in the Cauchy argument. We henceforth write $\Phi(v_g)$ for this form, never for an everywhere-defined physical operator.
\end{proof}

\node{GW06}{The precise coercive-form record defining the fiber Hamiltonian}{FB02A, EN01, EN02, BA03, GW-Form, GW04, DF01}
\begin{lemma}\label{lem:gwformrecord}
For every real $g$ and $P\in X$, define on the energy-coordinate space from EN01
\[
 a_{g,P}(x,y)=h_P^{\rm fr}(j_Px,j_Py)+\Phi(v_g)(j_Px,j_Py).
\]
There is a specified \texttt{GW.Operators.CoerciveFormData} record with this form and embedding $j_P$. Define $H_g(P)$ to be \texttt{GW.Operators.formOperator} of this record. This is a densely defined self-adjoint partial linear map on the existing \texttt{GW.Fock.Carrier}, and its domain and action are exactly the weak form-domain characterization used throughout the manuscript.
\end{lemma}
\begin{proof}
FB02A with $\theta=1/2$ provides a finite nonnegative number $b_g$ such that
$|\Phi(v_g)[u]|\leq\tfrac12h_P^{\rm fr}[u]+b_g\|u\|^2$, uniformly in $P$.
Fix $b_g=C_{\rm LY}(2g^2+8g^4)\geq0$, using the constant of FB01. Polarization BA03 gives continuity of the associated sesquilinear form on $\Q_P$. The free form is continuous because, in coordinates,
$h_P^{\rm fr}(j_Px,j_Py)=\langle x,y\rangle-\langle j_Px,j_Py\rangle$.
Both are Hermitian; their sum is an \texttt{AlgebraicForm} with the first variable conjugate-linear. The \texttt{continuous\_form} field follows from the mixed bound and continuity of bounded bilinear (conjugate-linear/linear) maps.

Set \texttt{embedding} $=j_P$, \texttt{shift} $=b_g+1$, \texttt{lower} $=1/2$, and \texttt{upper} $=2+2b_g$. EN01 supplies the required injectivity, dense range, and complete Hilbert carrier. For physical $u=j_Px$,
\begin{align*}
 h_{g,P}[u]+(b_g+1)\|u\|^2
 &\geq\tfrac12 h_P^{\rm fr}[u]+\|u\|^2\geq\tfrac12\|x\|^2,\\
 h_{g,P}[u]+(b_g+1)\|u\|^2
 &\leq\tfrac32 h_P^{\rm fr}[u]+(2b_g+1)\|u\|^2
 \leq(2+2b_g)\|x\|^2.
\end{align*}
The numerical fields \texttt{lower\_pos} and \texttt{lower\_le\_upper} follow from $b_g\geq0$. These calculations supply every field of the GW record, with no appeal to an operator associated to a form before that operator has been defined.

Apply \path{formOperator_selfadjoint}. Its definition of self-adjointness includes density of the domain. The exact source theorems \path{formOperator_domain_iff}, \path{formOperator_domain_of_equation}, and \path{formOperator_form_identity} say respectively that a physical $u$ belongs to this domain exactly when $u=j_Px$ for an energy coordinate $x$ and there is $f\in\F$ with $a_{g,P}(y,x)=\langle j_Py,f\rangle$ for every $y$, and that its value is this $f$. This is Definition~\ref{def:associated} literally on the fixed carrier. The value is unique because $j_P$ has dense range.

For sector compressions use the same construction on the closed coordinate subspace with the indicated particle coordinates zero. The embedding has dense range in the corresponding physical sector because the bounded-coordinate cutoffs preserve it. The inequalities and Hermiticity restrict. The record therefore constructs the actual compressed operator as well; it is not an assertion that an arbitrary compression of an unbounded operator is self-adjoint. For a finite cutoff replace $\Phi(v_g)$ by its bounded form on $V_N$; GW04 identifies it with the existing \texttt{ordinaryField}. The same construction and argument apply.
\end{proof}

\node{FB02}{Closed fiber forms, associated operators and selection rules}{FB02A, BA03, FA00, FD02, FC03, AN02, GW06, LD01}
\begin{lemma}\label{lem:cutoff}
For every real $g$, the form
$h_{g,p}=h_p^{\rm fr}+\Phi_g$ is closed and semibounded on $\Q$, and has a uniquely determined self-adjoint associated operator $H_g(p)$. The tail estimate is
\begin{equation}
 |(\Phi_g-\Phi_\Lambda)[u]|
 \leq Cs\Lambda^{-1/2}\|u\|_{\Q_p}^2.
 \label{eq:tailform}
\end{equation}
For $s\leq1$, $|\Phi_g(w,u)|\leq Cs\|w\|_{\Q_p}\|u\|_{\Q_p}$. Matrix entries vanish unless the two particle numbers differ by one; the vacuum couples only to the one-particle source $v_g$. Complex conjugation fixes the real form, and momentum reflection changes $p$ to $-p$.
\end{lemma}
\begin{proof}
For $0<s\leq1$, set $\theta=s$ in \eqref{eq:ly}. The right side is at most $C s\|u\|_{\Q_p}^2$. Apply BA03 to get the mixed estimate. For $s=0$ every cutoff form and its limit are zero. For arbitrary $g$, the full two-term remainder of FB01 at $\theta=1/2$ gives $h_{g,p}[u]\geq\tfrac12h_p^{\rm fr}[u]-b_g\|u\|^2$ with finite $b_g\geq0$. Adding $(1+b_g)\|u\|^2$ gives a positive form bounded below by $\tfrac12\|u\|_{\Q_p}^2$ and above by a finite constant times that norm. A Cauchy sequence in this shifted form norm is Cauchy in the complete graph norm; convergence in the graph norm implies convergence in the shifted norm by the upper bound. This proves completeness of the form norm, which is closedness. The operator is the specific GW form operator of GW06. Its exact domain and action are those of the supplied GW representation theorem; FA00 explains that theorem mathematically. No second associated-operator type is defined.

At finite cutoff, the sector formulas show the selection rule and give the vacuum--one-particle entry $\int v_\Lambda(k)f(k)dk$. Their norm limits preserve all vanishing entries by BA03. For the vacuum source, Cauchy--Schwarz with the one-particle graph weight and AN02 identifies the limit as $\ell_g(f)$. Conjugating the finite-cutoff formulas, which have real coefficients, proves $h_{g,p}(Jw,Ju)=\overline{h_{g,p}(w,u)}$. Changing every coordinate $k_i$ to $-k_i$ preserves product measure, fixes the even real form factor, and changes the free multiplier from $|K_n-p|^2$ to $|K_n+p|^2$. Thus reflection transports the cutoff forms and their limits. The associated-domain characterization transports these form symmetries to the operators. Coercive weak inverses inherit them by uniqueness of their solutions. This proves all the symmetries used later without a positivity-improving or energy-minimum theorem.
\end{proof}


\paragraph{A fixed mixed constant for all later estimates.}
From $|\Phi_g[u]|\leq s\,h_p^{\rm fr}[u]+C_Ls\|u\|^2$ for $0<s\leq1$ and BA03, fix once and for all
\[
 C_\Phi=2\max\{1,C_L\},\qquad
 |\Phi_g(w,u)|\leq C_\Phi s\|w\|_{\Q_p}\|u\|_{\Q_p}.
\]
At $s=0$ the same inequality follows from the vanishing form. Define $C_R=3C_\Phi$ and $C_D=16C_\Phi^2$. The proofs of RS01 and SP03 below establish their advertised bounds with precisely these constants. Thus these constants do not depend on an as-yet-constructed ground energy or on a particular momentum.

\node{SY01}{Conjugation and reflection on forms, sources and weak inverses}{GW01, GW04, GW05, EN01, EN02, FB02A, GW06, BA01, LD01, AN00}
\begin{lemma}\label{lem:symmetrytransport}
Let $J$ be coordinatewise complex conjugation and $\mathcal R$ coordinatewise momentum reflection:
$(Ju)_n=\overline{u_n}$ and $(\mathcal Ru)_n(k_1,\ldots,k_n)=u_n(-k_1,\ldots,-k_n)$, with scalar vacuum unchanged by $\mathcal R$. Then $J$ is an antilinear isometric involution, $\mathcal R$ is a linear unitary involution, and
\[
 h_{g,P}(Ju,Jw)=\overline{h_{g,P}(u,w)},\qquad
 h_{g,-P}(\mathcal Ru,\mathcal Rw)=h_{g,P}(u,w).
\]
They preserve the relevant sector domains. The same statements hold with any finite cutoff and after real scalar shifts. For every coercive compressed form used later, these identities pass to its weak inverse, both on physical sources and on the real, even source $v_g$.
\end{lemma}
\begin{proof}
On each ambient sector conjugation is an antilinear isometry since $|\overline z|=|z|$, and commutes with coordinate permutations; it therefore restricts to the existing symmetric sector. The scalar vacuum is conjugated in the usual way. Reflection preserves product Lebesgue measure and commutes with permutations, so its pullback is unitary on each sector. Their direct sums are well-defined on $\F$ by the squared norm sum. They satisfy $J^2=\mathcal R^2=I$. Multiplication by real $K_n$ commutes with $J$, while $K_n(-k_1,\ldots,-k_n)=-K_n(k_1,\ldots,k_n)$ proves the reflected kinetic equality. The weights show the asserted energy-domain preservation and isometries between the appropriate norms.

For a real even cutoff $v_{g,\Lambda}$, the sector creation formula and the integrated annihilation formula of GW04 imply
$J\Phi(v_{g,\Lambda})u=\Phi(v_{g,\Lambda})Ju$ and
$\mathcal R\Phi(v_{g,\Lambda})u=\Phi(v_{g,\Lambda})\mathcal Ru$ on $V_N$.
For annihilation, conjugate the absolutely integrable expression or substitute $l\mapsto-l$; for creation, apply the same operations to its finite sum. The norm bounds permit both operations for arbitrary $V_N$ inputs, or equivalently extend from finite-sector vectors. Pairing gives the form identities at cutoff. The mixed norm convergence FB02A, together with the boundedness of $J$ and $\mathcal R$ on the energy spaces, passes these equalities to the singular form. Sector projections commute with both maps. Real scalar shifts have the same identities.

Here is the inverse argument, including the antilinear case. Suppose $a(w,u)=\langle w,f\rangle$ for all physical energy tests. Then
$a(w,Ju)=\overline{a(Jw,u)}=\overline{\langle Jw,f\rangle}=\langle w,Jf\rangle$.
Uniqueness of the coercive weak solution gives $J Rf=R Jf$. For reflection the identical calculation with the linear unitary gives $\mathcal R R_pf=R_{-p}\mathcal R f$ whenever $p$ is the momentum parameter in $H(-p)$. For the singular source the right side is a load rather than a vector. Its defining integral has the same conjugation and reflection identities because $v_g$ is real and even and LD01 guarantees absolute integrability; substituting these in the same weak equation proves $Jh_p=h_p$ and $\mathcal Rh_p=h_{-p}$. At this point the conclusion is conditional only on having a coercive inverse; RS01 supplies those inverses. No statement about uniqueness of ground states is used.
\end{proof}

\subsection{Transporting locally retained GW definitions without an axiom}
\node{RE01}{Unitary transport of the entire form construction}{FD01, FD02, FC01, FC03, BA01, FA00, FB02A}
\begin{lemma}\label{lem:GWtransport}
Suppose an existing Fock realization has sectorwise unitary maps $U_n:H_n\to H_n^{\rm old}$ intertwining coordinate permutation symmetry, $K_n$, and multiplication by $n$. Suppose its smeared creation formula, on the bounded-coordinate finite-sector core, is the transport of the formula in FC01. Then $U=\bigoplus_nU_n$ transports all graph spaces, cutoff interaction forms, their form-norm limits, associated operators, eigenvectors, and real bounded-inverse resolvent sets in this document.
\end{lemma}
\begin{proof}
For finite sequences, $\|Uu\|^2=\sum_n\|U_nu_n\|^2=\sum_n\|u_n\|^2$; extend by norm limits and use the inverse sector maps to obtain a unitary. Intertwining of the real multipliers gives equality of weighted graph norms on the core. Approximate in the relevant graph norm by FC01's truncations (or their transported versions) to extend this equality and obtain the full graph-space isomorphism. Smeared annihilation is determined by the smeared adjoint relation on the dense core, so transports together with creation. The cutoff form therefore agrees on the core and, by its graph continuity, on the full form domain. The difference of the two transported limiting forms is the norm limit of zero cutoff differences; it is zero.

For associated operators, $a^{\rm old}(Uw,Uu)=a(w,u)$ and $\langle Uw,Uh\rangle=\langle w,h\rangle$ turn the domain characterization into the same characterization on the old space. FA00's uniqueness gives $A^{\rm old}=UAU^{-1}$ with domains transported exactly. If $R$ is a bounded two-sided inverse of $A-\lambda$, then $URU^{-1}$ has both inverse identities and maps into the transported domain. Apply the same argument to $U^{-1}$ for the converse. A nonzero eigenvector is transported to a nonzero eigenvector because $U$ is injective. Thus every spectral conclusion is invariant under the choice of these realizations.
\end{proof}
The audited GW source has a canonical Fock carrier already; no general transport should be introduced unless an actual type or convention mismatch is exhibited. The lemma states and proves the precise compatibility criterion for that fallback. There is a complete local canonical construction above; importing compatible existing proofs is a source-reuse choice, not a mathematical assumption in the main theorem. If GW uses factorial sector weights, the proposed $U_n$ is multiplication by the square root of the corresponding constant weight; substituting this in the two displayed creation formulas checks the normalization explicitly before any field identity is used.


\node{FB03}{Uniform fiber and sector bounds}{FB02A, FB02, FA02, GW06}
\begin{lemma}\label{lem:sectorbounds}
For $0\leq s=\abs g\leq1/4$, take $C_E=2C_L$, $C_1=1+C_L$, and $C_2=4C_L$, where $C_L$ is fixed by FB02A. Uniformly in $p$,
\begin{align}
 h_{g,p}[u]&\geq-C_Es^2\norm u^2,\label{eq:globalformlower}\\
 h_{g,p}[u]&\geq(1-s)\bigl(\norm{(\Pf-p)u}^2+\norm{N^{1/2}u}^2\bigr)
                  -C_Ls\norm u^2.\label{eq:linearbound}
\end{align}
Set $C_1=1+C_L$ and $L_s=1-C_1s$. On $\Q_{\geq1}$, for $\lambda<L_s$,
\begin{align}
 (h_{g,p}-\lambda)[u]
 \geq{}&(1-s)\norm{(\Pf-p)u}^2
 +(1-s)\ip u{(N-1)u}
 +(L_s-\lambda)\norm u^2.\label{eq:nonvac}
\end{align}
Consequently this form is coercive in the $\Q_p$ norm.

If $-C_Es^2\leq e\leq0$, then for $u\in\Q_{\geq2}$ and $\varepsilon\geq0$,
\begin{equation}
 (h_{g,p}-e-1+\varepsilon)[u]\geq\tfrac18\norm u_{\Q_p}^2,
 \qquad (h_{g,p}-e-1+\varepsilon)[u]\geq\tfrac14\norm u^2.\label{eq:twosector}
\end{equation}
The last two bounds require $C_2s^2\leq1/8$, independently of $p$.
\end{lemma}
\begin{proof}
For $s=0$, the interaction vanishes, so all assertions follow directly from the nonnegative free multipliers. Hence assume $s>0$ for the following choices of the relative-bound parameter. Use $\theta=1/2$ for \eqref{eq:globalformlower}, and $\theta=s$ for \eqref{eq:linearbound}. On the nonvacuum sector rewrite $N=(N-1)+1$, which gives \eqref{eq:nonvac}. Its right side dominates
\[
 \min\{1-s,(L_s-\lambda)/2\}\norm u_{\Q_p}^2,
\]
since $N+1=(N-1)+2$.

For the last assertion take $\theta=1/4$; its admissibility follows from $s\leq1/4$. With $C_2=4C_L$ the left side is at least
\[
 \tfrac34\norm{(\Pf-p)u}^2+\tfrac34\ip u{Nu}-(1+C_2s^2)\norm u^2.
\]
Subtract $\tfrac18\norm u_{\Q_p}^2$ and use $N\geq2$. The remaining coefficient of $\norm u^2$ is at least $1/8-C_2s^2$, hence nonnegative under the stated condition $C_2s^2\leq1/8$. Independently, the same expression is at least $(1/2-C_2s^2)\norm u^2$, which is at least $\norm u^2/4$. All estimates are uniform in $p$.
\end{proof}

\section{The low-energy root and the ground vector}\label{sec:vacuum}
This section constructs the ground-state energy before using it in any resolvent. It avoids circular reliance on a pre-existing eigenvalue or on the essential spectrum.

\node{SP00}{Monotonicity of a scalar source resolvent}{BA01, BA02, FA02}
\begin{lemma}\label{lem:sourcefunction}
Let $V,H$ be Hilbert spaces, $j:V\to H$ bounded, $z\in V$, and
$T_\lambda=T_0-\lambda j^*j$ coercive for $\lambda$ in an open real interval $I$. Set $a(\lambda)=\langle z,T_\lambda^{-1}z\rangle$. Then $a$ is real, nonnegative, continuous, and nondecreasing. Consequently $F(\lambda)=c-\lambda-a(\lambda)$ is continuous and
$F(\lambda_2)\leq F(\lambda_1)-(\lambda_2-\lambda_1)$ for $\lambda_1<\lambda_2$ in $I$.
\end{lemma}
\begin{proof}
Reality and nonnegativity follow from positivity of the inverse in BA02. At a fixed $\lambda_0$, write
$T_\lambda=T_{\lambda_0}(I-(\lambda-\lambda_0)T_{\lambda_0}^{-1}j^*j)$.
For $|\lambda-\lambda_0|\|T_{\lambda_0}^{-1}\|\|j\|^2<1/2$, the geometric inverse gives a local inverse norm bound. BA02's two-inverse identity gives
\[
 \|T_\lambda^{-1}-T_{\lambda_0}^{-1}\|
 \leq |\lambda-\lambda_0|\|T_\lambda^{-1}\|\|j\|^2\|T_{\lambda_0}^{-1}\|.
\]
Pair with $z$ to obtain continuity. If $\lambda_1<\lambda_2$, then $T_{\lambda_1}\geq T_{\lambda_2}>0$, and inverse order gives $a(\lambda_1)\leq a(\lambda_2)$. Subtraction gives the displayed strict-decrease bound for $F$. This argument avoids inverse differentiation in the ground-state construction.
\end{proof}

\node{SP01}{Vacuum Schur function}{FA01, FA02, FA04, FA05, FB03, AN02, SP00, ML-IVT, GW07, LD01}
Fix $\abs P\leq1/4$ and $\lambda<L_s$. Let $K_{P,\lambda}$ be the weak-form map on $\Q_{\geq1}$ for the compression of $h_{g,P}-\lambda$. It is invertible between the energy space and its conjugate dual by \eqref{eq:nonvac}. Define
\[
 h_{P,\lambda}=K_{P,\lambda}^{-1}v_g,\quad
 a_P(\lambda)=\ell_g(h_{P,\lambda}),\quad
 F_P(\lambda)=\abs P^2-\lambda-a_P(\lambda).
\]
Here and below applying $\ell_g$ to a nonvacuum Fock vector means applying it to that vector's one-particle component.

\begin{lemma}\label{lem:vacuumroot}
For sufficiently small $s$, there is a unique $e_P\in(-1,\abs P^2]$ with $F_P(e_P)=0$. Moreover,
\begin{equation}
 \sigR(H_g(P))\cap(-\infty,L_s)=\{e_P\},\qquad
 \Psi_P=\Omega-h_{P,e_P}\ne0,
 \quad H_g(P)\Psi_P=e_P\Psi_P.\label{eq:lowresult}
\end{equation}
The form $h_{g,P}-e_P$ is nonnegative. In particular $e_P=E_g(P)$ and
\[
 -C_Es^2\leq e_P\leq\abs P^2.
\]
\end{lemma}
\begin{proof}
Use the decomposition $\C\Omega\oplus\Q_{\geq1}$, with the usual scalar inner product on the vacuum factor and the graph inner product on the second factor. This direct-sum energy norm is equivalent to the original form norm. The energy-space Schur complement is exactly the scalar $F_P(\lambda)$.

The map $K_{P,\lambda}$ has bounded energy representative $T_0-\lambda j^*j$ for the fixed graph inner product. The Riesz source vector $z$ is independent of $\lambda$. Thus $a_P(\lambda)=\langle z,K_{P,\lambda}^{-1}z\rangle$. SP00 gives continuity and monotonicity of $a_P$, and hence the explicit strict-decrease inequality
\begin{equation}
 F_P(\lambda_2)\leq F_P(\lambda_1)-(\lambda_2-\lambda_1)
 \quad(\lambda_1<\lambda_2<L_s).
 \label{eq:vacderivative}
\end{equation}
No differentiation of a source-resolvent expression is needed in this step.

Choose $s$ so that $L_s>1/2$ and $C_Es^2<1/4$. The full form $h_{g,P}+1$ is at least $(3/4)\norm\cdot^2$ by \eqref{eq:globalformlower}. Evaluate it on $\Omega-h_{P,-1}$; Schur completion gives $F_P(-1)>0$. At $\lambda=\abs P^2\leq1/16<L_s$, positivity of $K_{P,\lambda}^{-1}$ gives $F_P(\abs P^2)=-a_P(\abs P^2)\leq0$. Apply ML-IVT to the continuous real function on the closed interval $[-1,|P|^2]$: its endpoint values straddle zero, so there is a root. The left endpoint is excluded by strict positivity there, and \eqref{eq:vacderivative} proves uniqueness even on the whole interval $(-\infty,L_s)$.

For every other $\lambda<L_s$, the scalar Schur complement and the nonvacuum block are invertible on their energy spaces. Lemmas~\ref{lem:schur} and~\ref{lem:forminverse} give $\lambda\in\resR(H_g(P))$. At the root the Schur kernel vector is $\Omega-h_{P,e_P}$, whose vacuum component is $1$; Lemma~\ref{lem:forminverse} makes it an operator-domain eigenvector. Schur completion at $e_P$ also proves $h_{g,P}-e_P\geq0$.

The set $\sigR(H_g(P))$ contains $e_P$ and contains no smaller number, since $e_P<L_s$. Thus its infimum is $e_P$, without invoking an independent theorem about the bottom of a spectrum. The lower bound follows either from \eqref{eq:globalformlower} evaluated on the eigenvector or from the nonnegativity just proved; the upper bound is the root's construction.
\end{proof}

\paragraph{Definition used from now on.}
Put
\[
 e=e_0=E_g(0),\quad \tau=1+e,\quad h_0=h_{0,e},\quad\Psi_0=\Omega-h_0.
\]
The vacuum Schur equation states
\begin{equation}
 e=-\ell_g(h_0),\qquad (H_g(0)-e)\Psi_0=0.\label{eq:groundidentities}
\end{equation}
In particular $-C_Es^2\leq e\leq0$. Note that $h_{g,P}[\Omega]=\abs P^2$ is a \emph{form} identity. Neither this argument nor any later one asserts $\Omega\in\Dom(H_g(P))$ for nonzero $g$.

\section{The near-threshold Schur complement at the resolvent level}\label{sec:onephonon}
\node{SP03}{One-phonon elimination and monotonicity}{FA01, FA02, FA03, FA04, FB02, FB03, SP01, BA02, GW07, EN02}
Let $0\leq\varepsilon\leq C_1s$ and $\abs P\leq1/4$. Take $s$ small enough that
\begin{equation}
 d_{P,\varepsilon}=1+e-\abs P^2-\varepsilon\geq\tfrac12.\label{eq:vacdenom}
\end{equation}
Let $K^{(\varepsilon)}_P$ be the positive form compression of $H_g(P)-e-1+\varepsilon$ to $\Q_{\geq2}$ and $X^{(\varepsilon)}_P$ its weak inverse. On the physical high sector it has norm at most $4$, and on the energy space it has the inverse bounds in \eqref{eq:twosector}.

Let $W_P:\Q_1\to\Q_{\geq2}^{\times}$ denote the off-diagonal interaction coupling, defined by its form pairing. Define the nonnegative form
\[
 D_P^{(\varepsilon)}(f,h)=\ip{W_P f}{X_P^{(\varepsilon)}W_P h}
\]
by the energy/dual pairing, and set
\begin{equation}
 O_P^{(\varepsilon)}[f]=\varepsilon\norm f^2+\norm{(k-P)f}^2-e\norm f^2
 -D_P^{(\varepsilon)}[f]+\frac{\abs{\ell_g(f)}^2}{d_{P,\varepsilon}}.
 \label{eq:onephonon}
\end{equation}
All these forms have domain $\Q_1$.

\begin{lemma}\label{lem:nearresolvent}
Under $C_Ds^2\leq1/2$, there are fixed $a=b=1/2$, independent of the displayed parameters, such that
\begin{equation}
 O_P^{(\varepsilon)}[f]\geq a\norm{\abs{k-P}f}^2-b\norm f^2.\label{eq:schurgraphlower}
\end{equation}
Also
\begin{equation}
 O_P^{(\varepsilon)}\geq\varepsilon I+O_P^{(0)}.\label{eq:monotone}
\end{equation}
If $O_P^{(0)}\geq0$ and $\varepsilon>0$, then
\[
 e+1-\varepsilon\in\resR(H_g(P)).
\]
\end{lemma}
\begin{proof}
The mixed constant just fixed and the isolated-to-inherited norm comparison in EN02 give
$\norm{W_Pf}_{\Q_{\geq2}^{\times}}\leq\sqrt2 C_\Phi s\norm f_{\Q_{1,P}}$.
The high-sector energy inverse has norm at most $8$ by \eqref{eq:twosector}; hence
\[
 0\leq D_P^{(\varepsilon)}[f]
 \leq8(\sqrt2 C_\Phi s)^2\norm f_{\Q_{1,P}}^2
 = C_Ds^2\norm f_{\Q_{1,P}}^2,
 \qquad C_D=16C_\Phi^2.
\]
Impose explicitly $C_Ds^2\leq1/2$. Dropping the nonnegative vacuum correction, $-e\|f\|^2$ and $\varepsilon\|f\|^2$ gives
\[
 O_P^{(\varepsilon)}[f]\geq
 (1-C_Ds^2)\|(k-P)f\|^2-C_Ds^2\|f\|^2
 \geq\tfrac12\|(k-P)f\|^2-\tfrac12\|f\|^2.
\]
Thus $a=b=1/2$ are valid fixed constants in \eqref{eq:schurgraphlower}. The smallness ledger includes this specific guard. The same mixed bounds and LD01 prove continuity on $\Q_1$.

The inverse of $K_P^{(0)}+\varepsilon j^*j$ decreases with $\varepsilon$ as a quadratic form on dual sources. One direct proof is the variational identity
\[
 \ip z{T^{-1}z}=\sup_{u\in V}\bigl(2\Re\ip z u-\ip u{Tu}\bigr),
\]
proved by completing the square with $u-T^{-1}z$. Increasing $T$ can only decrease this supremum. Thus $D_P^{(\varepsilon)}\leq D_P^{(0)}$. In addition, $d_{P,\varepsilon}^{-1}\geq d_{P,0}^{-1}$. Subtract the two expressions \eqref{eq:onephonon} to prove \eqref{eq:monotone}.

If the threshold form is nonnegative, \eqref{eq:monotone} gives $O_P^{(\varepsilon)}\geq\varepsilon I$. Combine this with \eqref{eq:schurgraphlower} and Lemma~\ref{lem:interpolation}, using the shifted one-particle graph norm, to obtain \emph{energy-space} coercivity. Its bounded energy representative is invertible by Lemma~\ref{lem:coercive}.

Now decompose the full form domain as
\[
 \Q=\C\Omega\oplus\Q_1\oplus\Q_{\geq2}.
\]
There is no vacuum--high-sector coupling by the particle-number selection rule. The complementary block to $\Q_1$ is the invertible energy block $\operatorname{diag}(-d_{P,\varepsilon},K_P^{(\varepsilon)})$. Its Schur complement is exactly \eqref{eq:onephonon}. Lemmas~\ref{lem:schur} and~\ref{lem:forminverse} therefore construct a bounded inverse of $H_g(P)-e-1+\varepsilon$ on the physical Hilbert space. This proves absence of \emph{all} spectrum at that energy.
\end{proof}

\paragraph{Explicit inverse equations for the implementation.}
For a physical right-hand side $(a,f,z)$, regarded as loads in the energy duals, solve
\begin{align*}
 O_P^{(\varepsilon)}u_1&=f+d_{P,\varepsilon}^{-1}\ell_g^*a
                         -W_P^*X_P^{(\varepsilon)}z,\\
 u_0&=d_{P,\varepsilon}^{-1}(\ell_g(u_1)-a),\\
 u_2&=X_P^{(\varepsilon)}(z-W_Pu_1).
\end{align*}
All right-hand-side maps are continuous into the indicated energy duals; the solutions are continuous into the energy spaces. These formulas, followed by the embedding into $\F$, are the bounded physical inverse. They are also a useful sign check on the positive vacuum term in \eqref{eq:onephonon}.

\section{Weighted resolvents and parameter differentiation}\label{sec:resolvents}
\node{RS01}{Uniform auxiliary resolvents}{FB02A, FB02, FB03, SP01, AN02, FD02, BA01, BA02, LD01, EN02}
For $n=0,1,2$, working on the Hilbert subspace $\Pi^{\geq n}\F$, define the positive multiplication operator
\[
 W_{n,p}^2=(\Pf+p)^2+N+1-n-e,
 \qquad G_{n,p}=W_{n,p}^{-1}.
\]
Here $W_{n,p}^2\geq1$ and $\norm{G_{n,p}}\leq1$. Also
\begin{equation}
 (\Pf+p)^2+N+1\leq3W_{n,p}^2.\label{eq:weightcompare}
\end{equation}
Let $V_{n,p}$ be the weighted energy space with norm $\norm{W_{n,p}u}$. As a map from the physical sector to this energy space, $G_{n,p}$ is an isometric isomorphism. Sandwich the interaction as a form and use Riesz representation to define
\[
 T_{n,p}=G_{n,p}\Phi(v_g)G_{n,p},\qquad M_{n,p}=I+T_{n,p}.
\]
These are bounded physical-sector operators, not products of unbounded operators on an unspecified domain.

\begin{lemma}\label{lem:resolvents}
Uniformly in $p\in\R^3$ and $n\in\{0,1,2\}$,
\[
 \norm{T_{n,p}}\leq C_Rs.
\]
For $C_Rs\leq1/2$, define
\begin{equation}
 R_p^{(n)}=G_{n,p}M_{n,p}^{-1}G_{n,p}.\label{eq:Rformula}
\end{equation}
It is the inverse of the compressed form $H_g(-p)-e+1-n$. It is positive and has physical operator norm at most $2$.
\end{lemma}
\begin{proof}
For physical $u,w$ in the compressed sector, \eqref{eq:weightcompare} gives $\|G_{n,p}u\|_{\Q_{-p}}\leq\sqrt3\|u\|$. The fixed mixed estimate therefore gives
$|\Phi_g(G_{n,p}w,G_{n,p}u)|\leq3C_\Phi s\|w\|\|u\|$.
BA01 identifies its representative with $T_{n,p}$, proving the bound with $C_R=3C_\Phi$. The geometric series in $-T_{n,p}$ gives $M_{n,p}^{-1}$, whose norm is at most $2$. Positivity follows from $M_{n,p}\geq I/2$ and the positive-form inverse identity. The weak form equation verifies \eqref{eq:Rformula}. No assertion $E_g(-p)\geq e$ is used.
\end{proof}

\paragraph{Uniform parameter range for all subsequent constants.}
The construction and estimates in this node are valid, with the same constants, for every reference number $e\in[-1/4,0]$, not only for the eventually selected ground energy. Indeed the only facts about $e$ used in the weights and their comparisons are $-e\geq0$ and $-e\leq1/4$. The constants $C_R$ and all weight-derivative constants may therefore be chosen before $g$, $P$, or $e=E_g(0)$ are chosen. SP01 and the preliminary cutoff in SM01 put that actual ground energy in this fixed interval. This order of quantifiers is used in the smallness ledger; no constant is defined using the particular coupling to which it is later applied.

Set, as in the source papers,
\begin{equation}
 Z_p=R_p^{(0)},\qquad Y_p=R_p^{(1)},\qquad X_P=R_{-P}^{(2)},\qquad
 h_p=Y_pv_g,\qquad\Ecal(p)=-\ell_g(h_p).
 \label{eq:auxdefs}
\end{equation}
The high resolvents are extended by zero to the other Fock sectors when paired with full Fock vectors. In particular $X_P=X_P^{(0)}$ from Section~\ref{sec:onephonon}. The vacuum root gives $\Ecal(0)=e$.

\subsection{Elementary parameter calculus in the norms actually used}
\node{CA01}{Norm differentiation, product rules and integrated remainders}{BA02, ML-FTC, ML-DIFF}
\begin{lemma}\label{lem:normcalculus}
The following statements hold for Banach-space-valued functions on an open subset of $\R^d$.
\begin{enumerate}
\item If a candidate derivative $D_iF$ is continuous and
$F(p+te_i)-F(p)=\int_0^t D_iF(p+re_i)\,dr$, then $F$ has that partial derivative in norm. Continuous partial derivatives give a Fr\'echet derivative $DF(p)h=\sum_i h_iD_iF(p)$.
\item For a continuous bilinear map $B$, differentiation of $B(F,G)$ has the two-term product rule, and repeated multi-index differentiation has coefficient $\binom\alpha\beta$ for $B(\partial^\beta F,\partial^{\alpha-\beta}G)$.
\item If $M(p)$ is norm differentiable and boundedly invertible, with its inverse locally bounded, then
$D_i(M^{-1})=-M^{-1}(D_iM)M^{-1}$. If $M$ is $C^r$, then its inverse is $C^r$ and
\[
 \partial^\alpha M^{-1}
 =-M^{-1}\sum_{0<\beta\leq\alpha}\binom\alpha\beta
       (\partial^\beta M)(\partial^{\alpha-\beta}M^{-1}).
\]
\item If $F$ has second derivative bounded by $L$ on a segment, then
$\|F(p+h)-F(p)-DF(p)h\|\leq L\|h\|^2/2$.
\end{enumerate}
\end{lemma}
\paragraph{Implementation contract: bounded calculus before new proofs.}
Reuse ML-DIFF for the bilinear product rule and first inverse derivative. For inversion use the real normed algebra of bounded complex-linear endomorphisms of the fixed Banach space. Package $M$ and its actual bounded two-sided inverse as a unit. The unit inverse agrees with that chosen inverse by BA02's uniqueness argument, so \path{hasFDerivAt_ringInverse} and the chain rule yield $A\mapsto-M^{-1}AM^{-1}$. This is not differentiation of an unbounded Hamiltonian. Preserve the displayed higher-derivative recurrence, the actual carrier/real-scalar instances, and the $L\|h\|^2/2$ remainder: those quantitative specializations still have to be supplied using the proof below. No equivalent general first-derivative theorem should be reimplemented.
\begin{proof}
For the first partial derivative, divide the integral formula by $t\ne0$ and subtract $D_iF(p)$. The resulting norm is at most the supremum of $\|D_iF(p+re_i)-D_iF(p)\|$ on the intervening segment, which tends to zero by continuity. To pass to a total derivative, telescope along the $d$ coordinate segments from $p$ to $p+h$. On segment $i$, subtract $h_iD_iF(p)$. Each error has norm at most $|h_i|$ times the supremum of the difference of that derivative from its value at $p$ in the radius-$\|h\|_1$ neighborhood. Summing gives at most $\|h\|_1$ times that modulus of continuity. Since $\|h\|_1\leq\sqrt d\|h\|$, division by $\|h\|$ gives an error tending to zero. This proves Fr\'echet differentiability and continuity of its derivative. Apply the same argument to the finitely many derivative maps to obtain higher differentiability.

For the bilinear rule, write the exact increment
\[
 B(F+\Delta F,G+\Delta G)-B(F,G)
 =B(\Delta F,G)+B(F,\Delta G)+B(\Delta F,\Delta G).
\]
The last term is bounded by $\|B\|\|\Delta F\|\|\Delta G\|=O(\|h\|^2)$. Subtract the two derivative terms in the first terms and use their differentiability remainders. Iteration gives the multi-index formula; its binomial coefficients satisfy Pascal's identity because a new differentiation hits either factor. Conjugate-linear first factors, such as complex inner products, are real-bilinear and the same rule applies over the real parameter field.

For inversion put $Q=M(p)^{-1}$, $Q_h=M(p+h)^{-1}$, and $E=M(p+h)-M(p)$. BA02 gives $Q_h-Q=-Q_hEQ$. Adding $QEQ$ leaves $(Q-Q_h)EQ$, whose norm is at most $\|Q_h\|\|Q\|^2\|E\|^2$. Subtracting the linear part of $E$ therefore proves the derivative formula. Local boundedness of $Q_h$ follows, when not already known, by factoring $M(p+h)=M(p)(I+QE)$ and applying the geometric inverse for $\|QE\|<1/2$. The first formula and the bilinear rule inductively prove higher regularity; differentiating $MQ=I$ and isolating the term $M\partial^\alpha Q$ gives the displayed recurrence.

For the remainder, apply the Banach-valued fundamental theorem of calculus ML-FTC to $t\mapsto F(p+th)$ and then to $t\mapsto DF(p+th)h$. The error equals
$\int_0^1\int_0^t D^2F(p+rh)[h,h]\,dr\,dt$.
The integrand norm is at most $L\|h\|^2$, and the triangle has area $1/2$. All derivatives here are continuous on a compact interval, hence Bochner integrable. This proves the bound without an unqualified exchange of derivatives and integrals.
\end{proof}

\node{CA02}{Uniform differentiation of the free inverse-square-root weights}{CA01, FD02}
\begin{lemma}\label{lem:weightderivatives}
For $b\geq1$, let $G_b(x)=(|x|^2+b)^{-1/2}$. For every coordinate multi-index $\alpha$ of order $r\leq4$, $\partial^\alpha G_b(x)$ is a finite sum
\[
 \sum_{(\beta,j)}c_{\beta,j}x^\beta(|x|^2+b)^{-1/2-j},
 \qquad |\beta|=2j-r,\quad 0\leq j\leq r.
\]
There is $C_r$ depending only on $r$ with
$|\partial^\alpha G_b(x)|\leq C_rG_b(x)$ for all $b\geq1,x$. These formulas differentiate the corresponding multiplication operators in operator norm, through order three, both into the physical space and into any fixed locally comparable weighted energy space.
\end{lemma}
\begin{proof}
At order zero take $c_{0,0}=1$. Differentiation of one term in direction $i$ gives
\[
 c\beta_i x^{\beta-e_i}w^{-1/2-j}
   -c(1+2j)x^{\beta+e_i}w^{-1/2-(j+1)},\qquad w=|x|^2+b.
\]
An absent monomial with $\beta_i=0$ contributes zero. Both new terms have the asserted relation between degree and order. The sum of absolute coefficients is at most $(3r+1)$ times the previous sum at the step from order $r$ to $r+1$, since $\beta_i\leq r$ and $j\leq r$. Starting at one therefore gives an explicit finite bound, for example $C_0=1$ and $C_{r+1}=(3r+1)C_r$. Since $|x|\leq\sqrt w$,
\[
 |x^\beta w^{-1/2-j}|/G_b(x)
 \leq w^{|\beta|/2-j}=w^{-r/2}\leq1.
\]
This proves the bound, including order four.

For norm differentiability, scalar Taylor's formula along a coordinate gives the pointwise difference quotient remainder bounded by $C_{r+2}|t|/2$ when differentiating an order-$r$ formula, $r\leq2$, and an integral difference of the order-four formula gives continuity of the order-three formula. More generally the first-order difference of a formula of order at most three is bounded by $C|t|$ using the next derivative. Multiplication by a bounded scalar function has operator norm at most its essential supremum, because integration of $|mf|^2\leq\|m\|_\infty^2|f|^2$ gives that bound. Thus these uniform remainder estimates give operator-norm derivatives and continuous derivatives. CA01 converts coordinate derivatives into the Fr\'echet derivatives.

For the energy-space assertion fix $p_0$ and let $|p-p_0|\leq1$, $x=K+p$, $x_0=K+p_0$. The inequalities
$|x_0|^2+b\leq2(|x|^2+b)+2\leq4(|x|^2+b)$ and its reverse show weight ratios at most two. Multiplying the derivative/remainder formulas by the fixed weight $\sqrt{|x_0|^2+b}$ therefore preserves uniform bounds (with a fixed enlarged constant) on this neighborhood. To control remainders after this multiplication use the same comparison at every intermediate point of the segment. This proves differentiability into the fixed energy space, not merely pointwise differentiability of an unbounded multiplier.
\end{proof}

\node{CA03}{Finite differences with small-distance weights}{CA01, ML-FTC}
\begin{lemma}\label{lem:finitedifferences}
Let $F$ be a scalar $C^r$ function of several Euclidean vector variables. Suppose every mixed derivative required below, including order zero, has operator norm at most $M$ globally. A difference in a variable by a vector $h$ is bounded by $2M\min(|h|,1)$ when a first derivative is available. A product of $l\leq r$ differences in different variables is bounded by
$2^lM\prod_{i=1}^l\min(|h_i|,1)$.
The same estimate applies to a difference of a derivative when the total order does not exceed $r$. If a real $C^2$ function $f$ is even and $\|D^2f\|\leq M$, then
$|f(h)-f(0)|\leq M|h|^2/2$.
\end{lemma}
\begin{proof}
For $|h|\leq1$ apply the fundamental theorem along the segment: the difference is $\int_0^1 DF(x+th)[h]dt$ and has absolute value at most $M|h|$. For $|h|>1$ use two values bounded by $M$. For several differences, call a displacement small when its length is at most one. Write integral derivative formulas for precisely the small displacements. Expand each large displacement as a difference of two terms; there are $2^{\#\mathrm{large}}$ terms. Each term is bounded by $M$ times the product of the small lengths, because the required mixed derivative is bounded. This gives the asserted bound and covers zero displacements as well. To apply it to a derivative, regard that derivative as the starting function and count its order in the available total order. Finally differentiating $f(-h)=f(h)$ at $h=0$ gives $Df(0)=-Df(0)$, hence zero. CA01's second-order remainder at zero gives the last bound.
\end{proof}

\node{RS02}{Norm derivatives through order three}{RS01, AN02, CA01, CA02, BA01, LD01, EN02}
\begin{lemma}\label{lem:derivatives}
For each multi-index $\alpha$ with $\abs\alpha\leq3$, $p\mapsto R_p^{(n)}$ and $p\mapsto h_p$ are $C^3$ in physical operator norm and Hilbert norm, respectively. There are constants independent of $p,g$ such that
\[
 \partial^\alpha R_p^{(n)}=G_{n,p}B_{\alpha,n,p}G_{n,p},\qquad
 \norm{B_{\alpha,n,p}}\leq C_\alpha,
 \qquad\norm{\partial^\alpha h_p}\leq C_\alpha s.
\]
Also $h_p\in\Q_{\geq1}$ and $\norm{W_{1,p}h_p}\leq Cs$.
\end{lemma}
\begin{proof}
The multiplication formula makes all derivatives explicit. On a particle sector of order $m\geq n$, write $x=K_m+p$, $b=m+1-n-e\geq1$, and $w=\abs x^2+b$. The scalar multiplier of $G$ is $w^{-1/2}$. Its first three derivatives are
\begin{align*}
 \partial_iG&=-x_iw^{-3/2},\\
 \partial_i\partial_jG&=-\delta_{ij}w^{-3/2}+3x_ix_jw^{-5/2},\\
 \partial_i\partial_j\partial_kG
 &=3(\delta_{ij}x_k+\delta_{ik}x_j+\delta_{jk}x_i)w^{-5/2}
   -15x_ix_jx_kw^{-7/2}.
\end{align*}
Dividing these multipliers by $G$ gives uniformly bounded multipliers because $\abs x\leq\sqrt w$ and $w\geq1$. CA02 proves, by its finite coefficient recurrence through order four, the uniform bound for the next derivatives and the operator-norm Taylor remainder. Apply it here with $b=m+1-n-e\geq1$. Thus the displayed multipliers are genuine operator-norm derivatives through order three. This is differentiation of multiplication operators, not of a formal unbounded $W$.

For a fixed $p_0$ and $\abs{p-p_0}\leq1$, the weights $W_{n,p}$ and $W_{n,p_0}$ are uniformly comparable. Thus the preceding derivatives hold as maps into the fixed energy space $V_{n,p_0}$. Differentiate the continuous form sandwich there. More explicitly, write $D_{\alpha,p}=G_{n,p}^{-1}\partial^\alpha G_{n,p}$ for the bounded real multiplier supplied by CA02, with $D_{0,p}=I$. For fixed physical vectors $u,w$, CA01's product rule for bounded forms gives
\[
 \langle w,\partial^\alpha T_{n,p}u\rangle
 =\sum_{\beta\leq\alpha}\binom{\alpha}{\beta}
  \Phi(v_g)(\partial^\beta G_{n,p}w,
                  \partial^{\alpha-\beta}G_{n,p}u).
\]
All derivatives here exist in the fixed local energy norm. Since each derivative of $G$ equals $G D$, each summand has bounded representative
$D_{\beta,p}^* T_{n,p}D_{\alpha-\beta,p}$. Therefore its norm is at most
$C_\beta C_{\alpha-\beta}C_Rs$. The finite sum gives $C_\alpha s$, uniformly in $p$. BA01 identifies the form derivatives with operator-norm derivatives; it is not a differentiation of an unbounded interaction operator.

Let $Q_p=M_{n,p}^{-1}$. Differentiating $M_pQ_p=I$ gives the explicit recurrence
\[
 \partial^\alpha Q_p
 =-Q_p\sum_{0<\beta\leq\alpha}\binom{\alpha}{\beta}
      (\partial^\beta M_p)(\partial^{\alpha-\beta}Q_p).
\]
It gives uniform bounds for all derivatives required. Apply the product rule to $GQG$ and factor a $G$ at each outside edge to obtain $B_{\alpha,n,p}$.

For the source, $z_p=G_{1,p}v_g$ is an actual one-particle $L^2$ vector. Lemma~\ref{lem:shiftedsource} gives $\norm{z_p}\leq Cs$, and the derivative multiplier formulas give $\norm{\partial^\alpha z_p}\leq C_\alpha s$. Here is the norm remainder verification for the source. Fix $p_0$ and restrict $|p-p_0|\leq1/2$. CA02's one-higher-derivative remainder, multiplied by $v_g$, is pointwise bounded by
$C|h|^2|G_{1,p_0}v_g|$ for a coordinate increment $h$ with $|h|\leq1/2$. The squared majorant is integrable by AN02. Taking its $L^2$ norm gives an $O(|h|^2s)$ remainder; the analogous first-difference bound gives norm continuity of each derivative. Repeating for derivative orders zero, one and two yields the $C^3$ assertion in $L^2$. This specifies the integrable majorant and the convergence topology without evaluating $v_g$ as a Hilbert vector. Now
\[
 h_p=G_{1,p}M_{1,p}^{-1}z_p.
\]
The product rule proves the Hilbert-norm assertions, and $W_{1,p}h_p=M_{1,p}^{-1}z_p$ proves energy-domain membership and its bound.
\end{proof}

\node{RS03}{The scalar self-energy has a small Hessian}{RS02, SP01, FB02, CA01, SY01}
\begin{lemma}\label{lem:selfenergy}
The function $\Ecal$ is real, even, and $C^2$, and
\[
 \sup_p\norm{D^2\Ecal(p)}\leq C_E' g^2,
 \qquad\abs{\Ecal(q)-e}\leq C_E' g^2\abs q^2.
\]
It is also uniformly bounded in absolute value by $Cg^2$.
\end{lemma}
\begin{proof}
Use the entirely bounded expression $\Ecal(p)=-\langle z_p,Q_pz_p\rangle$, where $Q_p=M_{1,p}^{-1}$. Real-parameter differentiation commutes with complex conjugation. The full formula, including all multiplicities, is
\[
 \partial^\alpha\Ecal(p)=
 -\sum_{\beta+\gamma+\delta=\alpha}
 \frac{\alpha!}{\beta!\gamma!\delta!}
 \langle\partial^\beta z_p,
       (\partial^\gamma Q_p)\partial^\delta z_p\rangle,
 \qquad |\alpha|\leq2.
\]
Cauchy--Schwarz bounds each summand by
$C_\beta s\, C_\gamma\, C_\delta s$. Summing the finitely many terms bounds each coordinate Hessian entry by $Cs^2$; for vectors $x,y\in\R^3$, summing their nine coordinate products and using $\sum_i|x_i|\leq\sqrt3|x|$ gives the asserted Hessian operator norm. The order-zero case gives the uniform bound on $|\Ecal|$. Reality follows from self-adjointness. Reflection gives $\Ecal(-p)=\Ecal(p)$, hence $D\Ecal(0)=0$. The fundamental theorem of calculus twice along $t\mapsto tq$ gives
\[
 \Ecal(q)-\Ecal(0)=\int_0^1(1-t)D^2\Ecal(tq)[q,q]\dd t.
\]
Use $\Ecal(0)=e$. No differentiability of the physical energy $P\mapsto E_g(P)$ is asserted or needed.
\end{proof}

\section{Cutoff inverses at a fixed energy and genuine pull-through statements}\label{sec:kernel}
In this section $e=E_g(0)$ is fixed once and for all by SP01. Every cutoff resolvent uses this \emph{same} number. It is not necessary to define a cutoff ground-state energy, prove its convergence, or prove that cutoff ground vectors converge. This removes SP02 from the version-1 dependency graph.

\node{UV01}{Fixed-energy cutoff convergence in the correct norms}{FB02A, FB02, RS01, AN02, FC02, BA02, FD02, LD01, EN02}
For $\Lambda\geq1$, let $v_\Lambda=v_g\mathbf1_{\abs k\leq\Lambda}$ and define $R_{p,\Lambda}^{(n)}$ by RS01 with $v_g$ replaced by $v_\Lambda$, leaving $e$ unchanged. Define $Z_{p,\Lambda},Y_{p,\Lambda},X_{P,\Lambda}$ accordingly, and
$h_{p,\Lambda}=Y_{p,\Lambda}v_\Lambda$ in the form-source sense.

\begin{lemma}\label{lem:fixedcutoff}
Uniformly for $p$ in each fixed bounded set, the inverses converge in norm from the common energy dual to the common energy space, and
\[
 h_{p,\Lambda}\longrightarrow h_p\quad\hbox{in the energy norm}.
\]
The physical inverse norms remain uniformly bounded for all $p,\Lambda$. If $A_\Lambda:\Q_1\to\Q_{\geq2}^{\times}$ is the cutoff one-to-high form coupling, then
\[
 A_\Lambda\to A,\qquad
 A_\Lambda^*X_{P,\Lambda}A_\Lambda\to A^*X_PA
\]
in the corresponding energy operator norms. Here $V^{\times}$ denotes the continuous conjugate-linear dual, or equivalently its Riesz representatives as specified in Section~\ref{sec:dualconventions}.
\end{lemma}
\begin{proof}
Use one fixed graph space on each bounded parameter set. Its norm and all the shifted norms are uniformly equivalent. If $T_{p,\Lambda},T_p$ are the bounded representatives of the compressed shifted forms, FB02 gives
\[
 \sup_{p\in K}\norm{T_{p,\Lambda}-T_p}\leq C_K\abs g\Lambda^{-1/2}.
\]
Coercivity is uniform by RS01, including for cutoffs. The exact identity
\[
 T_{p,\Lambda}^{-1}-T_p^{-1}
 =T_{p,\Lambda}^{-1}(T_p-T_{p,\Lambda})T_p^{-1}
\]
proves inverse convergence with the same rate up to $C_K$. The source-dual error obeys
\[
 \norm{v_\Lambda-v_g}_{V_{1,p}^{\times}}^2
 \leq \int_{\abs k>\Lambda}
       \frac{g^2}{\abs k^2(1+\abs{k+p}^2)}\dd k
 \leq C_Kg^2/\Lambda
\]
for large $\Lambda$, by AN02. Adding and subtracting
$R_{p,\Lambda}^{(1)}v_g$ gives the source-solution convergence. In particular convergence holds in the number-form norm, so FC02 gives convergence of their annihilation fields in $L^2(\R^3;\F)$.

The mixed tail-form estimate gives $A_\Lambda\to A$. Expand the difference of the three factors $A_\Lambda^*X_{P,\Lambda}A_\Lambda$ and $A^*X_PA$ by replacing them one at a time; each resulting norm is bounded by one norm error times two uniformly bounded factors. This proves the last assertion. No free \emph{operator}-graph norm is claimed to converge as $\Lambda\to\infty$.
\end{proof}

\node{CV01}{A norm-by-norm limit ledger for the ultraviolet approximation}{UV01, LD01, EN02, GW03, BA02, ML-BOCHNER, GW-Domination}
\begin{lemma}\label{lem:cutoffledger}
On any fixed bounded parameter set, all the following limits hold on fixed coordinate Hilbert spaces: compressed form representatives in operator norm; their inverses in operator norm; the resulting physical inverses in operator norm; source solutions in energy norm; annihilation fields of those solutions in $L^2(X;\F)$; and one-to-high Schur corrections in norm $\mathcal B(\Q_1,\Q_1^\times)$. No free-operator-graph convergence and no pointwise convergence of annihilation slices are asserted.
\end{lemma}
\begin{proof}
Choose one reference energy norm for each of the finitely many sector thresholds. EN02 bounds all coordinate changes on a bounded momentum set, so the constants in what follows are uniform on that set. Let $T_\Lambda,T$ be the form representatives, with $\|T_\Lambda-T\|\leq\delta_\Lambda\to0$ from FB02A, and let both lower coercivity constants be at least $a>0$. BA02 gives
\[
 \|T_\Lambda^{-1}-T^{-1}\|
 \leq a^{-2}\delta_\Lambda.
\]
Multiplication on both sides by the fixed physical embeddings bounds the physical inverse error by $\|j\|^2a^{-2}\delta_\Lambda$. For sources with Riesz vectors $z_\Lambda\to z$, insert and subtract $T_\Lambda^{-1}z$ to obtain
\[
 \|T_\Lambda^{-1}z_\Lambda-T^{-1}z\|
 \leq a^{-1}\|z_\Lambda-z\|+a^{-2}\delta_\Lambda\|z\|.
\]
LD01 gives the source error explicitly as the square root of the weighted tail integral. EN02 and GW03 give
$\|\mathcal A(u_\Lambda-u)\|_{L^2}\leq\|u_\Lambda-u\|_{V_N}\leq C\|u_\Lambda-u\|_V$.
Thus the field convergence has an explicit norm bound and is not obtained by evaluating slices.

For the Schur term write its energy representative as $B_\Lambda^*D_\Lambda^{-1}B_\Lambda$, with $B_\Lambda$ the one-to-high form coupling. The exact telescoping identity is
\begin{align*}
 B_\Lambda^*D_\Lambda^{-1}B_\Lambda-B^*D^{-1}B
 ={}&(B_\Lambda^*-B^*)D_\Lambda^{-1}B_\Lambda\\
 &+B^*(D_\Lambda^{-1}-D^{-1})B_\Lambda
   +B^*D^{-1}(B_\Lambda-B).
\end{align*}
The norm of the adjoint error equals that of the map error. Each product is bounded by the product of its three operator norms; one factor tends to zero and the others are uniformly bounded by FB02A and coercivity. This proves the last assertion. All limits retain the same reference energy $e$. They do not use cutoff ground-state energies, convergence of spectral projections, or convergence of eigenvectors.
\end{proof}

\node{UV02}{The cutoff operator domain equals the free operator domain}{FD02, FA00, FC01, FC03, FB02, BA02, GW06, GW04}
\begin{lemma}\label{lem:cutoffdomain}
For each finite $\Lambda$, each fiber momentum, and each compressed sector used here, the cutoff associated operator has the free multiplication-operator domain. On that domain it is the free operator plus the cutoff field and the indicated scalar shift.
\end{lemma}
\begin{proof}
Let $F=(\Pf-P)^2+N$ on the chosen sector and $V$ be the compressed cutoff field. The projection does not enlarge its norm. FC03 gives
\[
 \norm{Vu}\leq 2\norm{v_\Lambda}_2\norm{(N+1)^{1/2}u}.
\]
For $c\geq1$, the free multiplier yields
\[
 \norm{(N+1)^{1/2}(F+c)^{-1}}\leq c^{-1/2},
\]
because $(n+1)/(x+n+c)^2\leq1/c$ for $x\geq0,n\geq0$. Choose $c$ larger than the negative form lower bound and so large that
$2\norm{v_\Lambda}_2/\sqrt c<1/2$. Then
\[
 F+V+c=\bigl(I+V(F+c)^{-1}\bigr)(F+c)
\]
on $\Dom F$, and the first factor has a bounded inverse by its geometric series. Hence $F+V+c$ is bijective from $\Dom F$ to the physical sector.

For $u\in\Dom F$, the form representation and the smeared adjoint relation give $u\in\Dom H_\Lambda$ and $H_\Lambda u=Fu+Vu$. Conversely, for $u\in\Dom H_\Lambda$, solve
$(F+V+c)w=(H_\Lambda+c)u$ in $\Dom F$. The associated operator $H_\Lambda+c$ is injective by the strict lower form bound. Consequently $w=u$, proving domain equality. The same reasoning applies to all sector compressions and real scalar shifts. The chosen $c$ may grow with $\Lambda$; it is used only at a fixed cutoff and never in a limiting estimate.
\end{proof}

\node{PT01}{Annihilation pull-through proved through an energy-valued residual}{FC01, FC02, FC03, FC04, FC07, FC08, FC09, UV01, UV02, RS01, BA01, AN01, MS02, MS03, CV01, GW04, EN02}
Let $n\in\{1,2\}$ and $F\in\Pi^{\geq n}\F\cap V_N$ be a physical source. All scalar tests below are complex; the factor conjugated in the testing integral is displayed explicitly.
\begin{lemma}\label{lem:pullthrough}
For every fixed $p$, the following identity holds in $L^2(K;\Pi^{\geq n-1}\F)$ for each bounded measurable $K$:
\begin{equation}
 (\mathcal A R_p^{(n)}F)(k)
 =R_{p+k}^{(n-1)}\bigl(\mathcal AF(k)-v_g(k)R_p^{(n)}F\bigr).
 \label{eq:regularsourcePT}
\end{equation}
It consequently holds almost everywhere on $\R^3$. The corresponding statement with $v_\Lambda$ and all inverses at cutoff, with the same fixed energy $e$, also holds.
\end{lemma}
\begin{proof}
\textbf{1. Fix the cutoff and identify a legitimate energy field.}
Set $u=R_{p,\Lambda}^{(n)}F$. UV02 identifies its operator domain with the free multiplication domain. FC04 therefore gives
\[
 \int\|\mathcal Au(k)\|_{\Q_{-(p+k)}}^2dk
 =\sum_m\int m(|K_m+p|^2+m)|u_m|^2<\infty.
\]
Here the sign is chosen because $R_p^{(n)}$ uses $H(-p)$. The constant shifts in the inverse are irrelevant to this membership. Restrict to a bounded measurable $K$, and put $V=\Q_{\geq n-1}$ with the fixed graph norm at zero momentum. Equivalence of the shifted norms, uniformly for $k\in K$, implies $U(k)=\mathcal Au(k)\in L^2(K;V)$. To see strong measurability in this stronger space, apply MS02 sectorwise to the graph-weighted coordinates in the displayed finite integral and then MS03. The fixed-weight and shifted-weight products are boundedly related on $K$, so the same conclusion holds for the chosen fixed $V$; this is stronger than merely assuming each slice belongs to $V$.

\textbf{2. Write the residual in one fixed Hilbert space.}
Let $j:V\to\Pi^{\geq n-1}\F$ be the inclusion. Let $T(k)$ be the bounded representative of the form of
$L_{n-1,p+k,\Lambda}=H_{g,\Lambda}(-p-k)-e+1-(n-1)$ on this sector. Its dependence on $k$ is operator-norm continuous: subtracting its values leaves only
$2(k-k')\cdot P_{\rm f}+(|p+k|^2-|p+k'|^2)$ as a bounded form on the fixed graph space. Cauchy--Schwarz bounds each momentum-linear form by the product of graph norms, and the scalar coefficients are continuous polynomials. The cutoff interaction is fixed in $k$. Thus $\sup_K\|T(k)\|<\infty$. Define
\[
 r(k)=T(k)U(k)-j^*(\mathcal AF(k)-v_\Lambda(k)u)\in V.
\]
FC08 shows $T(k)U(k)\in L^2(K;V)$. The other two terms lie there because $\mathcal AF\in L^2(\R^3;\F)$, $j^*$ is bounded, $v_\Lambda\in L^2(K)$, and $u$ is fixed. Hence $r\in L^2(K;V)$, with no exceptional-set convention still to be chosen.

\textbf{3. Test the core identity and take the free graph limit.}
Choose $u_j\in\mathscr C\cap\Pi^{\geq n}\F$ converging to $u$ in the free operator graph norm, by FC01. FC07's identity \eqref{eq:projectedcorePT}, paired with $w\in\mathscr C\cap\Pi^{\geq n-1}\F$ and integrated against any $\chi\in L^2(K)$ (extended by zero), gives
\begin{align*}
 \int_K\overline\chi\,
    L_{n-1,p+k,\Lambda}(w,\mathcal Au_j(k))\,dk
 ={}&\langle a^*(\chi)w,L_{n,p,\Lambda}u_j\rangle\\
 &-\int_K\overline\chi\,v_\Lambda(k)\langle w,u_j\rangle\,dk.
\end{align*}
The source pairing equality is FC03's adjoint relation: the Bochner integral $\int\overline\chi\,\mathcal A(Lu_j)$ is $a(\chi)Lu_j$. Here $Lu_j\in\mathscr C$, so this use is legitimate. In contrast, no statement about convergence of these annihilation fields is made.

On the left, FC04 applied to $u_j-u$ gives convergence of the fields in $L^2(K;V)$; multiply its norm by $\|\chi\|_2\|w\|_V\sup_K\|T(k)\|$ to bound the tested difference. On the first term of the right, the cutoff field is free-operator bounded, as proved in UV02, so $L_{n,p,\Lambda}u_j\to L_{n,p,\Lambda}u=F$ in the physical Hilbert norm. Since $a^*(\chi)w\in\F$, its pairing converges. This limiting pairing equals $\int_K\overline\chi\langle w,\mathcal AF(k)\rangle dk$ by FC03, now legitimately using $F\in V_N$. The final term converges by the bound
$\|\chi\|_2\|v_\Lambda\|_{L^2(K)}\|w\|\|u_j-u\|$.
All integrals are absolutely integrable by these same inequalities. We have proved
$\int_K\overline\chi\langle w,r(k)\rangle_Vdk=0$ for every such $\chi,w$.

\textbf{4. Extract a single almost-everywhere vector equation.}
FC01 makes the test core dense in $V$. Apply FC09 to $r$, obtaining $r(k)=0$ outside a single null set in $K$. At each remaining point this is the weak form equation for all energy tests simultaneously:
\begin{equation}
 L_{n-1,p+k,\Lambda}(w,\mathcal Au(k))
 =\langle w,\mathcal AF(k)-v_\Lambda(k)u\rangle
 \quad(w\in V).
 \label{eq:weakPT}
\end{equation}
RS01 gives a coercive inverse of this form. Uniqueness follows directly: the difference of two solutions can be tested against itself, and coercivity forces it to be zero. Therefore its solution is
$R_{p+k,\Lambda}^{(n-1)}(\mathcal AF(k)-v_\Lambda(k)u)$.
This proves the cutoff statement as an $L^2(K;\F)$ identity. The equation in energy space was used only at finite cutoff.

\textbf{5. Remove the cutoff only in the claimed topology.}
UV01 gives $u_\Lambda=R_{p,\Lambda}^{(n)}F\to u=R_p^{(n)}F$ in energy norm. Since the energy norm dominates the number-form norm, FC02 gives $\mathcal Au_\Lambda\to\mathcal Au$ in global $L^2$ field norm. On a fixed bounded $K$, UV01 also gives
$\delta_\Lambda=\sup_{k\in K}\|R_{p+k,\Lambda}^{(n-1)}-R_{p+k}^{(n-1)}\|\to0$.
Choose $\Lambda$ larger than a radius containing $K$, so $v_\Lambda=v_g$ there. The difference of the right sides is bounded in $L^2(K)$ by
\[
 \delta_\Lambda\|\mathcal AF\|_{L^2(K)}
 +\|v_g\|_{L^2(K)}
   \bigl(\delta_\Lambda\|u\|+2\|u_\Lambda-u\|\bigr).
\]
This tends to zero; AN01 supplies local square integrability of $v_g$, and RS01 supplies the inverse norm bound two. For clarity, the finite-cutoff inverse family is continuous on each bounded $K$: its form representative on the fixed space $V$ depends continuously on $k$ because the kinetic difference is a polynomial form in $k$, while the interaction is independent of $k$. Coercivity gives a locally uniform inverse bound, and BA02's two-inverse identity gives continuity of its inverse. Composing with the fixed embeddings gives continuity of the physical inverse family. FC08 now gives strong measurability of its action on the fields. UV01 gives uniform-on-$K$ convergence of these inverse families, so their limits have the same measurability. Passing to the $L^2$ limit proves \eqref{eq:regularsourcePT}. Taking a countable union of the exceptional sets for balls of integer radius gives the asserted global almost-everywhere statement for this fixed $p$. It does not produce a common null set for all $p$, and no such conclusion is used later.
\end{proof}


\node{PT02}{Pull-through for the singular one-particle form source}{PT01, UV01, AN01, FC02, FC08, LD01, CV01}
\begin{lemma}\label{lem:sourcePT}
For every fixed $p$, as an equality in $L^2(K;\F)$ on each bounded $K$,
\begin{equation}
 (\mathcal A h_p)(k)=v_g(k)Z_{p+k}(\Omega-h_p).
 \label{eq:sourcePT}
\end{equation}
\end{lemma}
\begin{proof}
At cutoff use PT01 with $n=1$, $F=v_\Lambda$ in the one-particle sector, and $u=h_{p,\Lambda}$. FC02 gives
$\mathcal A v_\Lambda(k)=v_\Lambda(k)\Omega$. Thus
\[
 \mathcal A h_{p,\Lambda}(k)
   =v_\Lambda(k)Z_{p+k,\Lambda}(\Omega-h_{p,\Lambda}).
\]
Now use UV01 and precisely the local $L^2$ convergence estimate in the preceding proof. There is no application of $\mathcal A$ to the singular non-$L^2$ source $v_g$. The source itself enters only through its cutoff approximation and its continuous form functional.
\end{proof}
\paragraph{Specializations consumed below.}
For $z\in\Pi^{\geq2}\F\cap V_N$ and fixed $P$,
\begin{equation}
 (\mathcal A X_Pz)(k)
 =Y_{k-P}\bigl(\mathcal A z(k)-v_g(k)X_Pz\bigr).
 \label{eq:XfieldPT}
\end{equation}
For an input not entirely in the high sector, read $z$ as $\Pi^{\geq2}z$; because $Y$ vanishes on the vacuum, $Y\mathcal A\Pi^{\geq2}z=Y\mathcal A z$ whenever $z\in\Pi^{\geq1}\F\cap V_N$. These are the only projection conventions needed in the following kernel proof.

\section{The threshold kernel without pointwise creation or delta distributions}\label{sec:scalarkernel}
\node{KA01}{The finite-cutoff one-to-high map and its physical adjoint}{GW01, GW04, FC02, FC06, ML-PROD, ML-BOCHNER}
\begin{lemma}\label{lem:couplingadjoint}
For $v\in L^2(X)$, let $A_v:L^2(X)\to\Pi^{\geq2}\F$ be the bounded map with only two-particle component
\[
 (A_vf)_2(k,l)=\frac{v(k)f(l)+v(l)f(k)}{\sqrt2}.
\]
Then $A_vf=a^*(v)\iota_1f=a^*(f)\iota_1v$ and $\|A_v\|\leq\sqrt2\|v\|$. Its adjoint is
\[
 (A_v^*z)(k)=\sqrt2\int\overline{v(l)}z_2(k,l)\,dl.
\]
This formula is valid for every $z\in\Pi^{\geq2}\F$. If in addition $z\in V_N$, it equals $\langle\iota_1v,\mathcal Az(k)\rangle$ a.e. The latter notation is not used for an arbitrary vector outside $V_N$.
\end{lemma}
\begin{proof}
The displayed symmetric expression follows from \texttt{ordinaryCreationSector\_ae} at input particle number one, together with GW01's $j_1$. Its two summands each have product-space norm $\|v\|\|f\|$. The triangle inequality divided by $\sqrt2$ gives the norm bound, and shows it is a bounded map on the entire one-particle Hilbert space. Symmetry of the expression gives the equality with exchanged $f,v$.

For a high vector $z$, Fubini gives $z_2(k,\cdot)\in L^2$ for almost every $k$. Cauchy--Schwarz in $l$ and then integration in $k$ show the proposed adjoint is in $L^2$ with norm at most $\sqrt2\|v\|\|z_2\|$. Pair $A_vf$ with $z$. The two terms have equal integrals after exchanging $k,l$, since $z_2$ is symmetric. Each absolute integral is bounded by $\|v\|\|f\|\|z_2\|$ using Cauchy--Schwarz on $X^2$, so this exchange is legitimate. The resulting identity is $\langle A_vf,z\rangle=\langle f,A_v^*z\rangle$, establishing the adjoint. If $z\in V_N$, GW03 identifies the one-particle coordinate of its annihilation field as $\sqrt2 z_2(k,\cdot)$ and gives the final formula. Only that coordinate contributes to its inner product with $\iota_1v$.
\end{proof}

\node{PT03}{A scalar, integrated proof of the threshold kernel}{PT01, PT02, FC03, UV01, SP03, RS01, AN02, ML-PROD, ML-BOCHNER, KA01, CV01, TR01, LD01}
For $\abs P\leq1/4$ define
\begin{align}
 \mathcal J_P(q,r)
 &=\ip{\Omega-h_q}{Z_{P+q+r}(\Omega-h_r)}+\ip{h_q}{X_Ph_r},\\
 C_P(q,r)&=\frac1{1+e-\abs P^2}-\mathcal J_P(q,r).
 \label{eq:Ckernel}
\end{align}
Let $\tau_Pf(q)=f(q+P)$ and let $\widehat O_P$ be the translated threshold form $\tau_PO_P^{(0)}\tau_P^{-1}$.

\begin{lemma}\label{lem:kernelformula}
For $f,h\in\Q_1$ the translated sesquilinear threshold form is
\begin{align}
 \widehat O_P(f,h)
 ={}&\int\bigl(\abs q^2+\Ecal(q)-e\bigr)\overline{f(q)}h(q)\dd q\notag\\
 &+\iint\overline{f(q)}\frac g{\abs{q+P}}C_P(q,r)
                 \frac g{\abs{r+P}}h(r)\dd q\dd r.
 \label{eq:kernelformula}
\end{align}
The double integral is absolutely convergent. In particular the diagonal case gives the quadratic-form identity used later.
\end{lemma}
\begin{proof}
\textbf{Step 1: do all particle algebra at a finite cutoff.}
Suppress the cutoff subscript and use $v=v_\Lambda$, with the reference energy still $e$. All occurrences of $v$ in this step are genuine one-particle Hilbert vectors. For $h\in L^2(\R^3)$ let
\[
 A h=a^*(v)\iota_1h=a^*(h)\iota_1v,
 \qquad u=X_PA h.
\]
The equality is the explicit symmetric two-particle formula
\[
 (Ah)_2(k,l)=\frac{v(k)h(l)+h(k)v(l)}{\sqrt2}.
\]
It yields
$\mathcal A(Ah)(k)=h(k)v+v(k)\iota_1h$.
Using \eqref{eq:XfieldPT},
\begin{equation}
 \mathcal A u(k)=h(k)h_{k-P}
                  +v(k)Y_{k-P}(\iota_1h-u).
 \label{eq:kernelstep1}
\end{equation}
The source $Ah$ has only two particles, so it belongs to $V_N$; no further domain condition on $h$ is needed at cutoff.

\textbf{Step 2: identify the remaining scalar pairing for every parameter.}
Fix $q\in\R^3$. The vacuum component of \eqref{eq:sourcePT} says, a.e. in $l$,
\[
 (h_q)_1(l)=v(l)\ip\Omega{Z_{q+l}(\Omega-h_q)}.
\]
Taking its conjugate inside the $L^2$ pairing and using self-adjointness of $Z$ gives
\begin{equation}
 \ip{h_q}{\iota_1h}
 =\int h(l)v(l)\ip{\Omega-h_q}{Z_{q+l}\Omega}\dd l.
 \label{eq:kernelstep2}
\end{equation}
For the high term, self-adjointness of $X$ and FC03 give
\[
 \ip{h_q}{X_PA h}
 =\ip{X_Ph_q}{a^*(h)v}
 =\int h(l)\ip{\mathcal A X_Ph_q(l)}v\dd l.
\]
Here PT01 applies to $\Pi^{\geq2}h_q\in V_N$. Insert \eqref{eq:XfieldPT} and then \eqref{eq:sourcePT}. With $r=l-P$, self-adjointness of $Y_r,Z_{q+l},X_P$ gives
\begin{equation}
 \ip{\mathcal A X_Ph_q(l)}v
 =v(l)\left(\ip{\Omega-h_q}{Z_{q+l}h_{l-P}}
                         -\ip{h_q}{X_Ph_{l-P}}\right)
 \quad\hbox{a.e. in }l.
 \label{eq:kernelstep3}
\end{equation}
The vacuum component removed by $\Pi^{\geq2}$ is killed by $Y_r$, as stated after \eqref{eq:XfieldPT}. Subtracting the two scalar integrals proves, for \emph{every} fixed $q$,
\begin{equation}
 \ip{h_q}{\iota_1h-u}
       =\int h(l)v(l)\mathcal J_P(q,l-P)\dd l.
 \label{eq:kernelstep4}
\end{equation}
This is an identity of ordinary scalars for all $q$. We do not substitute a varying parameter into a merely a.e.-in-two-variables field formula.

\textbf{Step 3: pair the first annihilation identity and integrate.}
For $f\in L^2$,
\[
 D_P(f,h)=\ip{Af}{X_PAh}
        =\int\overline{f(k)}\ip v{\mathcal A u(k)}\dd k.
\]
Insert \eqref{eq:kernelstep1} and \eqref{eq:kernelstep4}, now legitimately with $q=k-P$. Since $Y$ is self-adjoint, the result is
\begin{align}
 D_P(f,h)={}&\int\overline{f(k)}h(k)\ip v{Y_{k-P}v}\dd k\notag\\
 &+\iint\overline{f(k)}v(k)\mathcal J_P(k-P,l-P)v(l)h(l)\dd k\dd l.
 \label{eq:integratedDkernel}
\end{align}
The diagonal term is an ordinary single integral, not a distributional delta. At cutoff, $\int\abs{vf}\leq\norm v_2\norm f_2$, and $\mathcal J$ is bounded by the uniform $h,Z,X$ estimates. Thus every exchange and pairing in these steps is absolutely integrable. All numerical factors have been accounted for by the displayed symmetric two-particle formula.

\textbf{Step 4: remove the cutoff in form norm and in the scalar integrals.}
UV01 gives convergence of $D_{P,\Lambda}$ in
$\mathcal B(\Q_1,\Q_1^{\times})$, convergence $h_{p,\Lambda}\to h_p$, and inverse convergence at every parameter. It follows that the scalar kernels converge pointwise (indeed uniformly on compact parameter sets). They obey one bound independent of $\Lambda,q,r$. The diagonal scalar $\ip{v_\Lambda}{Y_{q,\Lambda}v_\Lambda}$ is uniformly $O(g^2)$ and converges to $-\Ecal(q)$; this is a source-dual pairing, using UV01, not convergence of $v_\Lambda$ in unweighted $L^2$.

For $f,h\in\Q_1$, the absolute-integrability bound needed in the singular limit is
\[
 \int\abs{v_g(k)f(k)}\dd k
 \leq\left(\int\frac{g^2}{\abs k^2(1+\abs k^2)}\dd k\right)^{1/2}
          \norm f_{\Q_1}<\infty.
\]
The product of this bound for $f$ and $h$ dominates the double integrals; the diagonal terms are dominated by $C g^2\abs{fh}\in L^1$. Dominated convergence therefore proves \eqref{eq:integratedDkernel} in the singular form sense. After the substitutions $k=P+q,l=P+r$, combine it with SP03 at $\varepsilon=0$. The kinetic part is $\abs q^2$, the diagonal correction is $\Ecal(q)$, and the vacuum term supplies $(1+e-\abs P^2)^{-1}$. This gives \eqref{eq:kernelformula} with exactly the stated signs.
\end{proof}
\paragraph{Why this version is preferable for translation.}
There is no pointwise creation map, no product of distributions, no delta-function type, and no assumed invariance of a finite-particle core under a resolvent. The only unbounded-field operation surviving the cutoff limit is the continuous map $\mathcal A:V_N\to L^2(\R^3;\F)$; singular pairings stay in form duals. The commutator-to-resolvent step is the weak-equation argument in PT01, not a formal rearrangement of unbounded operators.

\section{Regularity and cancellation in the kernel}\label{sec:kernelbounds}
\node{KE01}{Kernel derivatives and reflection}{PT03, RS02, RS03, FB02, CA01, SY01}
\begin{lemma}\label{lem:Cregularity}
The kernel $C_P(q,r)$ is real, symmetric in $q,r$, and satisfies
\[
 C_{-P}(-q,-r)=C_P(q,r).
\]
On $\abs P\leq1/4$, all joint derivatives in $(P,q,r)$ through total order three are bounded by a constant $C_K$, uniformly in $q,r$ and small $g$.
\end{lemma}
\begin{proof}
Complex conjugation fixes $h_p,\Omega$ and commutes with the resolvents, so all displayed matrix elements are real. Self-adjointness then gives symmetry in $q,r$. Momentum reflection gives the covariance identity. For derivatives, use the following explicit rule. In a matrix element $\langle u(x),R(x)v(x)\rangle$, a real multi-index $\alpha$ gives
\[
 \sum_{\beta+\gamma+\delta=\alpha}
 \frac{\alpha!}{\beta!\gamma!\delta!}
 \langle\partial^\beta u(x),(\partial^\gamma R(x))\partial^\delta v(x)\rangle.
\]
Apply this with $x=(P,q,r)$ and the affine parameter $P+q+r$ for $Z$, or $-P$ for the compressed inverse. Each coordinate derivative of the first affine map has coefficient zero or one; that of the second has coefficient zero or minus one, and all higher derivatives vanish. Derivatives of $\Omega$ vanish. Every summand is therefore one of the derivatives bounded in RS02, with total order at most three, times its displayed integer coefficient. Its absolute value is bounded by the product of the three norms. There are finitely many multi-indices of total order at most three in nine real coordinates, so the maximum of these finite sums is one uniform constant. This specifies both all terms and their bounds, without differentiating an unbounded operator. The denominator stays bounded away from zero on an open neighborhood of the closed ball, for example $\abs P<1/3$ after the fixed preliminary smallness choice. Each derivative of the reciprocal through order three is therefore bounded there. This proves ordinary derivatives at boundary momenta as well, not merely within-set derivatives.
\end{proof}

Define
\begin{equation}
 c_P=C_P(0,0),\quad
 \psi_P(q)=C_P(q,0)-c_P,\quad
 F_P(q,r)=C_P(q,r)-c_P-\psi_P(q)-\psi_P(r).
 \label{eq:decomp}
\end{equation}
The symbol $F_P(q,r)$ here is a two-variable kernel and is distinct from the vacuum scalar $F_P(\lambda)$ in Section~\ref{sec:vacuum}; in Lean use different names, such as \module{vacuumSchurScalar} and \module{remainderKernel}.

\node{KE02}{Quantitative cancellation bounds}{KE01, CA03}
\begin{lemma}\label{lem:kernelbounds}
Uniformly for $\abs P\leq1/4$ and all $q,r$,
\begin{align}
 \abs{c_P-c_0}&\leq C\abs P^2,\label{eq:cbound}\\
 \abs{\psi_P(q)}&\leq C\omega(q)(\omega(q)+\abs P),\label{eq:psibound}\\
 \abs{\psi_P(q)-\psi_0(q)}&\leq C\abs P\omega(q),\label{eq:psidiff}\\
 \abs{F_P(q,r)}&\leq C\omega(q)\omega(r),\label{eq:Fbound}\\
 \abs{F_P(q,r)-F_0(q,r)}&\leq C\abs P\omega(q)\omega(r),\label{eq:Fdiff}\\
 \norm{D_qF_P(q,r)}+\norm{D_q^2F_P(q,r)}&\leq C\omega(r).\label{eq:Fderivative}
\end{align}
Furthermore $\psi_P(0)=F_P(q,0)=F_P(0,r)=0$.
\end{lemma}
\begin{proof}
The vanishing statements follow by substitution. The function $P\mapsto c_P$ is even by reflection and has bounded Hessian; Taylor's integral formula gives \eqref{eq:cbound}. This replaces the much stronger perturbative lower-bound calculation for $c_P$ in the source paper.

At $P=0$, $q\mapsto C_0(q,0)$ is even, hence its gradient at $0$ vanishes. The mixed $P,q$ derivative bound implies $\norm{D_qC_P(0,0)}\leq C\abs P$. For $\abs q\leq1$, Taylor's formula in $q$ now gives $\abs{\psi_P(q)}\leq C(\abs P\abs q+\abs q^2)$. For $\abs q>1$, use the uniform bound on $C_P$. This is \eqref{eq:psibound}.

The difference $\psi_P-\psi_0$ vanishes at $q=0$. Integrating the mixed $P,q$ derivative first along $tP$ and then along $uq$ proves $C\abs P\abs q$ for $\abs q\leq1$; for larger $q$, a single $P$ integration gives $C\abs P$. This proves \eqref{eq:psidiff}.

When $\abs q,\abs r\leq1$, the exact double-difference formula is
\[
 F_P(q,r)=\int_0^1\int_0^1D_qD_rC_P(tq,ur)[q,r]\dd u\dd t.
\]
This proves \eqref{eq:Fbound} in that region. If just one variable is small, integrate only in that variable and use bounded first derivatives; if neither is small, use boundedness of $C_P$. Apply the same argument to $C_P-C_0$, integrating once more in $P$, to get \eqref{eq:Fdiff}; at most three derivatives are used. Finally, for $j=1,2$,
\[
 D_q^jF_P(q,r)=D_q^jC_P(q,r)-D_q^jC_P(q,0).
\]
Integrate in $r$ when $\abs r\leq1$ and use uniform bounds otherwise, proving \eqref{eq:Fderivative}.
\end{proof}

\section{Bounded one-particle operators and the unnormalized form}\label{sec:boundedkernel}
\node{OP00}{A bounded integral operator from a product majorant}{ML-LP, ML-L2, ML-PROD, BA01}
\begin{lemma}\label{lem:dominatedkernel}
Let $a\geq0$ be a measurable $L^2$ function on $\R^3$, and let $k(q,r)$ be a measurable complex kernel with $|k(q,r)|\leq M a(q)a(r)$ off a fixed null set in either variable (values there may be reset to zero). The formula
$(Kf)(q)=\int k(q,r)f(r)dr$ defines a bounded linear map on $L^2$ with norm at most $M\|a\|_2^2$. If $k(q,r)=\overline{k(r,q)}$ almost everywhere on the product, $K$ is self-adjoint. If the kernel is real, $K\overline f=\overline{Kf}$.
\end{lemma}
\begin{proof}
For any measurable representative of $f\in L^2$, Cauchy--Schwarz gives
$I_f=\int a(r)|f(r)|dr\leq\|a\|_2\|f\|_2<\infty$.
For every $q$ with finite $a(q)$, the kernel integral is absolutely integrable and bounded by $Ma(q)I_f$. On the remaining null set set it to zero. To check measurability, restrict $q$ to a bounded box $D$. Since $a\in L^1(D)$ by Cauchy--Schwarz against $\mathbf1_D$, the product integrand has absolutely integrable bound $Ma(q)a(r)|f(r)|$ on $D\times\R^3$. Fubini supplies a measurable integral in $q$ there. Exhausting by boxes gives global measurability. Integrating the square of the bound proves
$\|Kf\|_2\leq M\|a\|_2I_f\leq M\|a\|_2^2\|f\|_2$.
If $f=f'$ outside a null set $N$, each of the difference integrals is zero because its integrable integrand is supported in $N$. Thus the definition descends to equivalence classes. Linearity of absolutely convergent integrals gives linearity of the map.

For $f,h\in L^2$, the absolute integral of the pairing integrand is at most
\[
 M\Big(\int|f(q)|a(q)dq\Big)
    \Big(\int a(r)|h(r)|dr\Big)
 \leq M\|a\|_2^2\|f\|_2\|h\|_2.
\]
Fubini therefore allows interchanging $q,r$ and conjugating the scalar integrals. The kernel symmetry gives $\langle f,Kh\rangle=\langle Kf,h\rangle$, the adjoint equality. Reality similarly commutes with complex conjugation, first pointwise outside the stated null sets and then as $L^2$ vectors.
\end{proof}


\node{OP01}{Construction of the bounded operators}{OP00, AN03, KE02, RS03, FA02}
Set
\[
 m(q)=1+\frac{\Ecal(q)-e}{\abs q^2}\quad(q\ne0),\qquad m(0)=1.
\]
Only measurability, evenness, and the bound $\norm{m-1}_\infty\leq Cg^2$ are needed. A continuous extension at zero is \emph{not} required. Define the integral operator with kernel
\[
 K_P(q,r)=w_P(q)F_P(q,r)w_P(r)
\]
and the vector
\[
 b_P(q)=g w_P(q)\psi_P(q).
\]
Let $B_P=M_m+g^2K_P$, where $M_m$ is multiplication by $m$ on $L^2(\R^3;\C)$.

\begin{lemma}\label{lem:boundedkernel}
The operator $K_P$ is bounded and self-adjoint, and $b_P\in L^2$. Uniformly in the momentum ball,
\[
 \norm{K_P}\leq C,\quad\norm{b_P}\leq C\abs g,\quad
 \norm{B_P-I}\leq Cg^2.
\]
For small $g$, $B_P\geq I/2$, $\norm{B_P^{-1}}\leq2$, and $m\geq3/4$.
\end{lemma}
\begin{proof}
Let $a_P(q)=\omega(q)w_P(q)$. Lemma~\ref{lem:singular} gives $a_P\in L^2$ uniformly. The kernel bound is
\[
 \abs{K_P(q,r)}\leq C a_P(q)a_P(r).
\]
Thus for $f\in L^2$, its integral exists almost everywhere and
\[
 \abs{(K_Pf)(q)}\leq C a_P(q)\norm{a_P}_2\norm f_2,
 \quad\norm{K_Pf}_2\leq C\norm{a_P}_2^2\norm f_2.
\]
This constructs the continuous linear map directly. Independence of $L^2$ representatives follows from the same integrable domination. For two test vectors, absolute integrability follows from the product $(\int\abs f a_P)(\int\abs h a_P)$; Fubini and the real symmetric kernel prove self-adjointness. No Hilbert--Schmidt theory is needed.

Since $\omega+\abs P\leq5/4$, \eqref{eq:psibound} gives $\abs{b_P}\leq C\abs g\,a_P$, proving its bound. Lemma~\ref{lem:selfenergy} bounds $m-1$. The final positivity and inverse estimates follow from $\norm{B_P-I}\leq1/2$, and the sharper $m\geq3/4$ from shrinking $g$ once more.
\end{proof}

\node{OP02}{Form decomposition valid also at zero momentum}{PT03, OP01, OP00, AN02, AN03, TR01, LD01}
For $f\in\Q_1$ define
\[
 u_f(q)=\abs q f(q),\qquad
 \ell_{g,P}(f)=\int\frac g{\abs{q+P}}f(q)\dd q.
\]
The latter is a continuous functional on $\Q_1$ by Lemma~\ref{lem:shiftedsource}.

\begin{lemma}\label{lem:unnormalizedform}
For every $\abs P\leq1/4$ and $f\in\Q_1$,
\begin{equation}
 \widehat O_P[f]=\ip{u_f}{B_Pu_f}
       +c_P\abs{\ell_{g,P}(f)}^2
       +2\Re\left(\overline{\ell_{g,P}(f)}\ip{b_P}{u_f}\right).
 \label{eq:unnormalizedform}
\end{equation}
For $P\ne0$, $\ell_{g,P}(f)=\ip{t_P}{u_f}$ with $t_P=g w_P\in L^2$.
\end{lemma}
\begin{proof}
Insert \eqref{eq:decomp} into \eqref{eq:kernelformula}. The $F_P$ part plus the kinetic/self-energy part equals $\ip{u_f}{B_Pu_f}$. The constant part is $c_P\abs{\ell_{g,P}(f)}^2$. The two one-variable parts are complex conjugates and sum to the last term. Every pairing is finite by AN02 and OP00. For the mixed terms, $\int |v_g(q+P)f(q)\psi_P(q)|dq<\infty$ because $\psi_P$ is bounded and AN02 bounds $\int|v_g(q+P)f(q)|dq$. Thus Fubini factors each product into the stated two scalar integrals, with no density argument required. For $P\ne0$ the last assertion follows by cancellation of $\abs q$ almost everywhere and \eqref{eq:tint}.
\end{proof}
Notice that this formula is meaningful at $P=0$ even though $t_0\notin L^2$. We retain the functional $\ell_{g,0}$ there instead of introducing a nonexistent Hilbert-space vector.

\node{OP03}{The polarized unnormalized decomposition}{OP02, OP00, PT03, AN02}
\begin{lemma}\label{lem:polarizedO}
For $f,h\in\Q_1$, put $u_f(q)=|q|f(q)$ and $L_f=\ell_{g,P}(f)$. Then
\begin{equation}
 \widehat O_P(f,h)=\langle u_f,B_Pu_h\rangle
 +c_P\overline{L_f}L_h
 +\overline{L_f}\langle b_P,u_h\rangle
 +\langle u_f,b_P\rangle L_h.
 \label{eq:polarizedO}
\end{equation}
Each of the four expressions is a continuous form on $\Q_1$. This identity is valid at $P=0$ without an $L^2$ vector representing $L_f$ on the unweighted physical space.
\end{lemma}
\begin{proof}
The diagonal multiplier of the threshold form equals $|q|^2m(q)$ almost everywhere. Its integral is therefore $\langle u_f,M_m u_h\rangle$. In the two-variable remainder term, cancel $|q|$ and $|r|$ outside their singleton zero sets to write it as $g^2\langle u_f,K_Pu_h\rangle$; OP00 guarantees the absolute integrability of this expression. The constant kernel term factors by Fubini as $c_P\overline{L_f}L_h$. The term with $\psi_P(r)$ factors as
$\overline{L_f}\int g|r+P|^{-1}\psi_P(r)h(r)dr
=\overline{L_f}\langle b_P,u_h\rangle$,
since $b_P=g\psi_P/(|r||r+P|)$ is real. The term with $\psi_P(q)$ is $\langle u_f,b_P\rangle L_h$. All factorizations are legitimate because $\int|g f/(|q+P|)|$ is finite by AN02 and $\psi_P$ is bounded, or because OP00 supplies the stronger $L^2$ product estimate. The first term is bounded by $\|B_P\|\|f\|_{\Q_1}\|h\|_{\Q_1}$, the second by $|c_P|\|\ell_{g,P}\|^2\|f\|_{\Q_1}\|h\|_{\Q_1}$, and the other two by $\|\ell_{g,P}\|\|b_P\|$ times these graph norms. This proves continuity and the complete polarized formula directly, without an implicit polarization step in the cancellation proof.
\end{proof}

\section{Seiringer's identity without normalization by \texorpdfstring{$c_0$}{c0}}\label{sec:seiringer}
The indispensable external argument is Seiringer's exact cancellation, not his entire spectral theorem. The following version proves the equality directly and does not assume $c_0>0$. It reorganizes the pull-through calculation using the ground-vector equation, avoiding Seiringer's extra two resolvent-product identities.

\node{SE01}{The ground source field and sector equations}{SP01, SP03, PT02, PT03, FB02, RS01, UV01, GW01, SY01}
Write
\[
 f_0=\Pi^1h_0\in\Q_1,\quad r_0=\Pi^{\geq2}h_0\in\Q_{\geq2},
 \quad u_0(k)=\abs k f_0(k)\in L^2.
\]
Use $A=W_0$ for the one-to-high form coupling, $X=X_0$ for the high inverse, and $D=A^*XA$ for the one-particle form from \eqref{eq:onephonon}.

\begin{lemma}\label{lem:groundPT}
The following equations hold, with the first in local Bochner $L^2$ and the others in the indicated form duals:
\begin{align}
 (\mathcal A h_0)(k)&=v_g(k)Z_k(\Omega-h_0),\label{eq:groundPT}\\
 (\abs k^2+1-e)f_0+A^*r_0&=v_g,\label{eq:groundone}\\
 (X^{-1}+I)r_0&=-Af_0,\quad\text{equivalently}\quad
 Xr_0=-r_0-XAf_0.\label{eq:groundhigh}
\end{align}
Here $X^{-1}$ is notation for the high-sector form, not an inverse of a zero-extended operator on all of $\F$.
\end{lemma}
\begin{proof}
The first assertion is PT02 at $p=0$. In particular it does not require commuting an annihilation operator through a singular ground-state eigenvalue equation. The nonvacuum weak equation
$(\Pi^{\geq1}H_g(0)\Pi^{\geq1}-e)h_0=v_g$
follows from the definition $h_0=Y_0v_g$. Project it onto the one and high sectors. The interaction has no diagonal particle-number part and the source has no high component, giving \eqref{eq:groundone} and \eqref{eq:groundhigh}. Applying the high weak inverse gives the second equality. The same two sector equations hold at every cutoff with the fixed reference energy $e$ and $h_{0,\Lambda}=Y_{0,\Lambda}v_\Lambda$. At cutoff no assertion that $e$ is a cutoff eigenvalue is needed.
\end{proof}

Define a bounded real function
\[
 \lambda_{\mathrm s}(k)=\ip{\Omega-h_k}{Z_k(\Omega-h_0)}+\ip{h_k}{Xh_0}.
\]
By \eqref{eq:Ckernel},
\begin{equation}
 c_0=(1+e)^{-1}-\lambda_{\mathrm s}(0),\qquad
 \psi_0(k)=\lambda_{\mathrm s}(0)-\lambda_{\mathrm s}(k).\label{eq:lambda}
\end{equation}

\node{SE02}{The one-particle source identity with its cutoff-dual passage}{SE01, PT01, PT02, PT03, UV01, FC03, AN02, KA01, LD01, CV01}
\begin{lemma}\label{lem:keyidentity}
In $\Q_1^{\times}$,
\begin{equation}
 (|k|^2-e-D)f_0=v_g(k)(1-\lambda_{\rm s}(k)).
 \label{eq:keyidentity}
\end{equation}
The right side denotes the continuous conjugate-linear functional obtained by integration against the displayed measurable function, not a claim that the function is in unweighted $L^2$.
\end{lemma}
\begin{proof}
\textbf{1. Work entirely at a finite cutoff, at the fixed number $e$.}
In the next calculation write $v=v_\Lambda$, $h_p=h_{p,\Lambda}$, $X=X_{0,\Lambda}$, $Z=Z_\Lambda$, and $f_0=\Pi^1h_{0,\Lambda}$, $r_0=\Pi^{\geq2}h_{0,\Lambda}$. Both sector equations in SE01 hold for these vectors because they are the source solution at energy $e$. They do not use a cutoff ground state. The cutoff coupling $A:L^2\to\Pi^2\F$ is the bounded map $Af=a^*(v)\iota_1f$; its adjoint on the high sector is the one-particle function
\[
 (A^*z)(k)=\sqrt2\int\overline{v(l)}z_2(k,l)\,dl
         =\langle v,\mathcal Az(k)\rangle
 \quad\text{a.e.}
\]
The first formula is the two-sector adjointness calculation FC06. The second uses the normalization $\sqrt2$ in FC02. The right side ignores all components of $\mathcal Az$ except the one-particle component, as the physical vector $v$ has just that component.

\textbf{2. Use vector field equalities with fixed source, not a diagonal substitution.}
PT02 with parameter zero is the vector identity
$\mathcal Ah_0(k)=v(k)Z_k(\Omega-h_0)$ outside one null set. Its vacuum component is $f_0(k)\Omega$; its nonvacuum component is $\mathcal Ar_0(k)$. Since $h_k$ has zero vacuum component and $X$ kills the one-particle component,
\begin{align*}
 v(k)\lambda_{\rm s}(k)
 &=v(k)\langle\Omega-h_k,Z_k(\Omega-h_0)\rangle
        +v(k)\langle h_k,Xr_0\rangle\\
 &=f_0(k)-\langle h_k,\mathcal Ar_0(k)\rangle
        +v(k)\langle h_k,Xr_0\rangle.
\end{align*}
Pairing the vector identity at a remaining $k$ with the continuously defined vector $h_k$ is legitimate: the identity holds as a vector, not just for each fixed test separately.

Apply PT01 with the fixed physical source $r_0\in V_N\cap\Pi^{\geq2}\F$ and parameter zero. Its vector identity is
$\mathcal A Xr_0(k)=Y_k(\mathcal Ar_0(k)-v(k)Xr_0)$.
Pair with the physical cutoff vector $v$ and use self-adjointness of $Y_k$ and $Y_kv=h_k$. Subtraction yields
\[
 v(k)\lambda_{\rm s}(k)=f_0(k)-\langle v,\mathcal A Xr_0(k)\rangle
                       =f_0(k)-(A^*Xr_0)(k).
\]
Each field equality has a fixed parameter and fixed source. Taking the union of their two null sets is enough; there is no substitution $p=k$ into an identity quantified as ``for each $p$, almost every $k$''.

\textbf{3. Substitute the two sector source equations.}
The high equation gives $Xr_0=-r_0-XAf_0$. Apply the bounded cutoff map $A^*$ to obtain
$-A^*Xr_0=A^*r_0+D f_0$, where $D=A^*XA$.
The one-particle equation gives $A^*r_0=v-(|k|^2+1-e)f_0$. Therefore
\[
 v\lambda_{\rm s}=f_0+A^*r_0+D f_0
                 =v-(|k|^2-e)f_0+D f_0.
\]
Rearranging proves the cutoff identity. Every vector involved is in the relevant domain: $h_{0,\Lambda}$ is in the cutoff operator domain and thus the free operator domain by UV02, so its one-particle component has $|k|^2f_0\in L^2$. Alternatively all these equalities can be read in the form dual, which is the topology used next.

\textbf{4. Pass to the singular equation in the one-particle form dual.}
Restore cutoff subscripts. UV01 gives $f_{0,\Lambda}\to f_0$ in $\Q_1$ and $D_\Lambda\to D$ as maps $\Q_1\to\Q_1^{\times}$. The kinetic multiplication is bounded between these two spaces, because
\[
 \left|\int |k|^2\overline{w(k)}f(k)dk\right|
 \leq\||k|w\|_2\||k|f\|_2.
\]
Thus the left sides converge in the form-dual norm, by adding and subtracting $D_\Lambda f_0$ in the last term.

The scalar functions $\lambda_{{\rm s},\Lambda}(k)$ converge for every fixed $k$, using UV01 on each resolvent and source vector in their definition. RS01 bounds all of them and their limit by a constant independent of $\Lambda,k$. Hence the squared weighted difference of the right sides is bounded by
\[
 \frac{C g^2}{|k|^2(1+|k|^2)},
\]
an integrable majorant by AN02. The difference converges pointwise almost everywhere to zero, since $v_\Lambda(k)\to v_g(k)$ and the scalar factors converge. Dominated convergence gives convergence in the weighted $L^2$ norm with weight $(1+|k|^2)^{-1}$. Cauchy--Schwarz against a test $w\in\Q_1$ bounds its dual norm by this weighted norm. Thus the right sides converge in $\Q_1^{\times}$ as well, proving \eqref{eq:keyidentity}. This proves all domain and convergence assertions in the statement, with no unweighted limit of $v_\Lambda$.
\end{proof}

\node{SE03}{Exact saturation forced by infrared nonintegrability}{SE02, OP03, OP01, SP01, AN01}
\begin{theorem}[Unnormalized Seiringer identity]\label{thm:saturation}\label{lem:saturation}
For all sufficiently small nonzero $g$,
\begin{equation}
 N_0:=\langle b_0,B_0^{-1}b_0\rangle=c_0.
 \label{eq:saturation}
\end{equation}
No positive lower bound on $c_0$, and no division by $\sqrt{c_0}$, is used.
\end{theorem}
\begin{proof}
Write $\tau=1+e>0$, $b=b_0$, $B=B_0$, $c=c_0$, and $u=|k|f_0\in L^2$. The vacuum Schur root gives
$\ell_g(f_0)=\ell_g(h_0)=-e$; here the functional only sees the one-particle sector. The threshold form at zero is
$O_0=|k|^2-e-D+\tau^{-1}\ell_g^*\ell_g$.
Consequently SE02, paired with any $f\in\Q_1$, implies
\begin{align*}
 O_0(f,f_0)
 &=\int\overline{f(k)}v_g(k)(1-\lambda_{\rm s}(k))dk
       -e\tau^{-1}\overline{\ell_g(f)}\\
 &=\int\overline{f(k)}v_g(k)(\tau^{-1}-\lambda_{\rm s}(k))dk\\
 &=c\,\overline{\ell_g(f)}+\int\overline{f(k)}v_g(k)\psi_0(k)dk.
\end{align*}
The second equality uses $1-e/\tau=1/\tau$. All integrals are finite: $v_g$ times a bounded scalar function is in the one-particle form dual by AN02, and the second integral can also be written using $b$ after the next substitution.

Fix $m\geq1$ and the annulus $A_m=\{1/m\leq|k|\leq m\}$. For any $z\in L^2$ supported there, put $f_z(k)=z(k)/|k|$ on the annulus and zero elsewhere. Then $\|f_z\|_2\leq m\|z\|_2$ and $|k|f_z=z$, so $f_z\in\Q_1$. Put $t_0(k)=g/|k|^2$ away from zero and zero at zero; it is square integrable on each $A_m$ but not near zero globally. We have
$\overline{\ell_g(f_z)}=\langle z,t_0\rangle$ on this annulus and
$\int\overline{f_z}v_g\psi_0=\langle z,b\rangle$.
The polarized formula OP03, with $\ell_g(f_0)=-e$, gives
\[
 O_0(f_z,f_0)=\langle z,Bu\rangle
       -ec\langle z,t_0\rangle
       +\langle z,t_0\rangle\langle b,u\rangle
       -e\langle z,b\rangle.
\]
Equating this with the preceding source expression and collecting terms proves
\[
 \langle z,\tau b-Bu\rangle
  =\langle z,t_0\rangle\bigl(\langle b,u\rangle-\tau c\bigr)
 \qquad(z\text{ supported in }A_m).
\]
The difference of the two vector functions in this equality is in $L^2(A_m)$. Choose $z$ to be that difference restricted to $A_m$; its squared norm is zero. Thus they agree almost everywhere on $A_m$. The union of these annuli is $\R^3\setminus\{0\}$, and the singleton has zero Lebesgue measure. Taking a countable union of null sets proves
\begin{equation}
 \tau b-Bu=t_0\bigl(\langle b,u\rangle-\tau c\bigr)
 \quad\text{almost everywhere on }\R^3.
 \label{eq:infraredcancel}
\end{equation}

The left side is in $L^2$: $b,u\in L^2$ and $B$ is bounded. If the scalar coefficient on the right were nonzero, division by that constant would make $t_0$ an $L^2$ function. AN01 rules this out, because $|t_0|^2=g^2|k|^{-4}$ and $g\ne0$. Therefore $\langle b,u\rangle=\tau c$ and $Bu=\tau b$. Invertibility of $B$ gives $u=\tau B^{-1}b$. Pair with $b$ and use real $\tau$ to get
$\tau c=\tau\langle b,B^{-1}b\rangle$. Since $\tau>0$, divide by $\tau$ to obtain \eqref{eq:saturation}. All divisions are by a nonzero radius on an annulus, a hypothetically nonzero scalar in a contradiction, or $\tau\geq1/2$; none is by $c$ or by the coupling in the final algebra.
\end{proof}


\section{The weighted integral equation on ordinary \texorpdfstring{$L^\infty$}{L infinity}}\label{sec:weighted}
We use the existing space $L^\infty(\R^3;\R)$, not a new pointwise weighted-supremum Banach space. This removes a completeness construction and makes null-set conventions explicit. The one-particle Hilbert space remains complex; the auxiliary functions here are real.

\node{WE01}{A bounded $L^\infty$ kernel operator}{KE02, AN03, OP01, ML-LP, ML-PROD}
For a measurable function $\eta$ with $\abs\eta\leq C\omega$ a.e. put
\[
 (\mathcal L_P\eta)(q)=\int F_P(q,r)\mu_P(r)\eta(r)\dd r.
\]
On $L^\infty$ define, for $q\ne0$,
\[
 (\mathscr L_P\zeta)(q)=\frac1{\omega(q)}
       \int F_P(q,r)\mu_P(r)\omega(r)\zeta(r)\dd r,
 \qquad (\mathscr L_P\zeta)(0)=0,
\]
and $\chi_P(q)=\psi_P(q)/\omega(q)$ for $q\ne0$, with $\chi_P(0)=0$.

\begin{lemma}\label{lem:Linftyoperator}
These formulas define bounded real-linear maps $\mathscr L_P:L^\infty\to L^\infty$, with norm at most a constant $C_W$ independent of $P,g$. Also $\norm{\chi_P}_\infty\leq C_W$ and multiplication by $m^{-1}$ has norm at most $4/3$.
\end{lemma}
\begin{proof}
For $q\ne0$, the kernel absolute value is bounded by
\[
 \frac{\abs{F_P(q,r)}}{\omega(q)}\mu_P(r)\omega(r)
 \leq C\omega(r)^2\mu_P(r).
\]
This is integrable in $r$ uniformly by Lemma~\ref{lem:singular}. Therefore the integral is independent of the $L^\infty$ representative, is measurable in $q$ by the measurable integral theorem, and has absolute value at most $C_W\norm\zeta_\infty$ for almost every $q$. Choose a representative bounded outside a null set, assign zero on that set, and use the same domination if an everywhere-bounded representative is needed for the Bochner integral. Different choices give the same output because integration is against a density with respect to Lebesgue measure. The source bound follows from \eqref{eq:psibound}; the multiplication bound follows from $m\geq3/4$. All these claims descend to the existing $L^\infty$ quotient.
\end{proof}

\node{WE02}{Existence, uniqueness, and the representative equation}{WE01, OP01, BA02, ML-COMPLETE}
\begin{lemma}\label{lem:weightedequation}
If $\norm{g^2M_{m^{-1}}\mathscr L_P}\leq1/2$, then the equation
\begin{equation}
 m\zeta_P+g^2\mathscr L_P\zeta_P=\chi_P\quad\text{in }L^\infty\label{eq:Linftyequation}
\end{equation}
has a unique solution with $\norm{\zeta_P}_\infty\leq C$. The function $\eta_P=\omega\zeta_P$ satisfies, a.e.,
\begin{equation}
 m\eta_P+g^2\mathcal L_P\eta_P=\psi_P,\qquad\abs{\eta_P(q)}\leq C\omega(q).
 \label{eq:etaequation}
\end{equation}
\end{lemma}
\begin{proof}
Multiply \eqref{eq:Linftyequation} by $m^{-1}$. The inverse of $I+g^2M_{m^{-1}}\mathscr L_P$ is the norm-convergent series in $-g^2M_{m^{-1}}\mathscr L_P$, with norm at most $2$. This proves existence, uniqueness, and the uniform estimate. Multiply the representative equation by $\omega$ to obtain \eqref{eq:etaequation}. For a convenient representative valid pointwise, replace $\eta_P$ by
\[
 m(q)^{-1}\bigl(\psi_P(q)-g^2(\mathcal L_P\eta_P)(q)\bigr).
\]
Call this representative $\eta'_P$. The almost-everywhere equation proves $\eta'_P=\eta_P$ almost everywhere. For every fixed $q$, the two integrals $\mathcal L_P\eta'_P(q)$ and $\mathcal L_P\eta_P(q)$ coincide: their integrands differ only on that same null set in $r$, and the previous absolute-integrability bound applies. Substituting this equality into the definition of $\eta'_P$ proves its equation at every $q$. At $q=0$, both $\psi_P(0)$ and $F_P(0,r)$ vanish, while $m(0)=1$, so $\eta'_P(0)=0$. The same pointwise domination bounds $|\eta'_P(q)|$ by a fixed multiple of $\omega(q)$. All later inequalities may still be read almost everywhere, so this representative is a convenience, not an extra regularity assumption.
\end{proof}

\node{WE03}{Quadratic vanishing by symmetric second differences}{WE02, KE01, KE02, AN00, AN03, CA01}
\begin{lemma}\label{lem:etaquadratic}
The zero-momentum solution is even a.e. and
\begin{equation}
 \abs{\eta_0(q)}\leq C\omega(q)^2\quad\almosteverywhere.\label{eq:etaquadratic}
\end{equation}
\end{lemma}
\begin{proof}
At $P=0$, $m,\omega,\chi_0$ are even and $F_0(-q,-r)=F_0(q,r)$. The change of variables $r\mapsto-r$ shows that reflection commutes with $\mathscr L_0$. Uniqueness in \eqref{eq:Linftyequation} therefore makes $\zeta_0$, hence $\eta_0$, even a.e.

Set $J(q)=\int F_0(q,r)\mu_0(r)\eta_0(r)\dd r$. The bound $\abs{\eta_0}\leq C\omega$ and AN03 make this an absolutely convergent integral for every $q$. The same change of variables proves that $J$ is exactly even, and $F_0(0,r)=0$ gives $J(0)=0$. For each fixed $r$, apply the fundamental theorem twice to $t\mapsto F_0(tq,r)$. The Hessian bound in \eqref{eq:Fderivative} gives
\[
 \abs{F_0(q,r)+F_0(-q,r)-2F_0(0,r)}
       \leq C\abs q^2\omega(r).
\]
For example the two remainders are the integrals over $t\in[0,1]$ of $(1-t)D_q^2F_0(\pm tq,r)[q,q]$, each bounded by $C\abs q^2\omega(r)/2$. Integrate the inequality against $\mu_0(r)\abs{\eta_0(r)}$ to obtain
\[
 2\abs{J(q)}\leq C\abs q^2\int\omega(r)^2\mu_0(r)\dd r
       \leq C\abs q^2.
\]
For $\abs q\geq1$, the ordinary kernel bound gives $\abs J\leq C$ instead. Hence globally $\abs J\leq C\omega^2$.

Finally $\abs{\psi_0}\leq C\omega^2$ by \eqref{eq:psibound}; solve \eqref{eq:etaequation} for $\eta_0$ and use $m^{-1}\leq4/3$. This proves \eqref{eq:etaquadratic}. No differentiation under an integral, differentiability of an $L^\infty$ representative, or differentiability of $m$ at zero is used.
\end{proof}

\node{WE04}{Lipschitz comparison in momentum}{WE03, WE02, AN03, KE02}
\begin{lemma}\label{lem:etacomparison}
Uniformly for $\abs P\leq1/4$,
\begin{align}
 \abs{\eta_P(q)-\eta_0(q)}&\leq C\abs P\omega(q)\quad\almosteverywhere,\label{eq:etadiff}\\
 \abs{\eta_P(q)}&\leq C\omega(q)(\omega(q)+\abs P)\quad\almosteverywhere.\label{eq:etasharp}
\end{align}
\end{lemma}
\begin{proof}
The difference of the two integral operators on the \emph{specific} vector $\eta_0$ is
\begin{align*}
 ((\mathcal L_P-\mathcal L_0)\eta_0)(q)
 ={}&\int(F_P-F_0)(q,r)\mu_P(r)\eta_0(r)\dd r\\
 &+\int F_0(q,r)(\mu_P(r)-\mu_0(r))\eta_0(r)\dd r.
\end{align*}
The first term is at most $C\abs P\omega(q)\int\omega^3\mu_P$, using \eqref{eq:Fdiff} and \eqref{eq:etaquadratic}. The second is at most $C\omega(q)\int\omega^3\abs{\mu_P-\mu_0}$. Lemma~\ref{lem:singular} bounds their sum by $C\abs P\omega(q)$.

Subtract the two equations \eqref{eq:etaequation}:
\[
 (m+g^2\mathcal L_P)(\eta_P-\eta_0)
 =\psi_P-\psi_0-g^2(\mathcal L_P-\mathcal L_0)\eta_0.
\]
Divide by $\omega$ and use the uniform $L^\infty$ inverse from Lemma~\ref{lem:weightedequation}. The right side divided by $\omega$ is bounded by $C\abs P$, by \eqref{eq:psidiff} and the just-proved bound. This gives \eqref{eq:etadiff}; add \eqref{eq:etaquadratic} to obtain \eqref{eq:etasharp}.
\end{proof}
It is important that this is not an assertion of $O(\abs P)$ operator-norm continuity of $\mathcal L_P$ on the whole weighted space. The argument only needs its action on the better-vanishing vector $\eta_0$.

\node{WE05}{Identification with the Hilbert-space inverse}{WE02, OP01, OP00, AN03}
\begin{lemma}\label{lem:identifyinverse}
The $L^2$ vector
\[
 x_P(q)=g w_P(q)\eta_P(q)
\]
is well-defined and equals $B_P^{-1}b_P$.
\end{lemma}
\begin{proof}
By the $C\omega$ bound in WE02, $\norm{x_P}_2^2\leq Cg^2\int\omega^2\mu_P<\infty$. The kernel integral is absolutely convergent by the same domination used to construct $K_P$. Multiplying \eqref{eq:etaequation} by $g w_P(q)$ gives, a.e.,
\[
 m(q)x_P(q)+g^2(K_Px_P)(q)=g w_P(q)\psi_P(q)=b_P(q).
\]
Here $(K_Px_P)(q)=g w_P(q)(\mathcal L_P\eta_P)(q)$ follows by substituting $x_P(r)=g w_P(r)\eta_P(r)$ in the already absolutely convergent single integral and using $w_P(r)^2=\mu_P(r)$. No interchange of two integrations is needed. Thus $B_Px_P=b_P$ in $L^2$; invertibility of $B_P$ gives the result.
\end{proof}

\node{WE06}{The three scalar estimates}{WE03, WE04, WE05, AN03, KE02, OP01, BA02}
For all $P$ define $N_P=\ip{b_P}{B_P^{-1}b_P}$. For $P\ne0$ define
\[
 U_P=\ip{b_P}{B_P^{-1}t_P},\qquad T_P=\ip{t_P}{B_P^{-1}t_P}.
\]
These scalars are real in this model, though the abstract positivity lemma below allows complex $U$.

\begin{proposition}\label{prop:scalars}
Uniformly for $0<\abs P\leq1/4$,
\begin{equation}
 \abs{U_P}\leq C_Ug^2,\qquad
 0\leq T_P\leq C_Tg^2/\abs P,\qquad
 \abs{N_P-N_0}\leq C_Ng^2\abs P.\label{eq:scalarbounds}
\end{equation}
\end{proposition}
\begin{proof}
Self-adjointness, reality, and Lemma~\ref{lem:identifyinverse} give the absolutely convergent integral formulas
\[
 U_P=g^2\int\eta_P(q)\mu_P(q)\dd q,
 \qquad N_P=g^2\int\psi_P(q)\eta_P(q)\mu_P(q)\dd q.
\]
The $U_P$ estimate follows from \eqref{eq:etasharp} and \eqref{eq:weightint}. For $T_P$, positivity and $\norm{B_P^{-1}}\leq2$ give
\[
 0\leq T_P\leq2\norm{t_P}_2^2\leq C_Tg^2/\abs P.
\]
For $N_P$, the bounds on $\psi$ and $\eta$ imply
\[
 \abs{\psi_P\eta_P-\psi_0\eta_0}\leq C\abs P\omega^2,
 \qquad\abs{\psi_0\eta_0}\leq C\omega^4\leq C\omega^3.
\]
Indeed expand the first difference as $(\psi_P-\psi_0)\eta_P+\psi_0(\eta_P-\eta_0)$ and use the uniform $C\omega$ bounds. Therefore
\begin{align*}
 \abs{N_P-N_0}
 &\leq g^2\int\abs{\psi_P\eta_P-\psi_0\eta_0}\mu_P
    +g^2\int\abs{\psi_0\eta_0}\abs{\mu_P-\mu_0}\\
 &\leq C_Ng^2\abs P
\end{align*}
by \eqref{eq:weightint} and \eqref{eq:weightdiff}. This is the needed Lipschitz estimate; a weaker $\abs P^{1/2}$ estimate would not compensate the singular size of $T_P$.
\end{proof}

\section{Positivity with no square roots and no positive \texorpdfstring{$c_P$}{cP} assumption}\label{sec:positivity}
The following bounded-Hilbert-space lemma replaces both the normalization by $\sqrt{c_P}$ and the Sherman--Morrison calculation. It uses only an ordinary completion of a square and Cauchy--Schwarz for a positive form.

\node{PO01}{An unnormalized two-vector positivity lemma}{FA02, BA02}
\begin{lemma}\label{lem:rankpositivity}
Let $B$ be bounded self-adjoint on a complex Hilbert space with $B\geq\beta I$, $\beta>0$. Let $b,t$ be vectors and $c\in\R$. Put
\[
 x=B^{-1}b,\quad z=B^{-1}t,\quad
 N=\ip b{B^{-1}b},\quad U=\ip b{B^{-1}t},\quad T=\ip t{B^{-1}t}.
\]
Suppose $1+U\ne0$ and
\begin{equation}
 \abs{1+U}^2+T(c-N)\geq0.\label{eq:scalarcriterion}
\end{equation}
Then the bounded Hermitian form
\[
 s[u]=\ip u{Bu}+c\abs{\ip t u}^2
          +2\Re\bigl(\overline{\ip t u}\ip b u\bigr)
\]
is nonnegative. No hypothesis on the sign of $c$ is required.
\end{lemma}
\begin{proof}
Set $a=\ip t u$ and $y=u+ax$. Expanding gives the exact identity
\begin{equation}
 s[u]=\ip y{By}+(c-N)\abs a^2.\label{eq:twovectorsquare}
\end{equation}
Also
\[
 \ip t y=(1+\ip t x)a=(1+\overline U)a.
\]
The form $(v,w)\mapsto\ip v{Bw}$ is an inner product, and hence
\[
 \abs{\ip t y}^2=\abs{\ip z{By}}^2
 \leq\ip z{Bz}\ip y{By}=T\ip y{By}.
\]
For completeness this particular inequality follows without a new Hilbert-space structure. If $T=\langle z,Bz\rangle=0$, coercivity gives $z=0$ and both sides are zero. Otherwise $T>0$. Put $r=\langle z,By\rangle$ and $\gamma=r/T$. Expanding positivity at $y-\gamma z$ gives
$0\leq\langle y,By\rangle-|r|^2/T$.
Multiplication by $T$ gives exactly the inequality above, including its complex conjugations. Thus
\[
 \abs a^2\leq\frac{T}{\abs{1+U}^2}\ip y{By}.
\]
If $c-N\geq0$, \eqref{eq:twovectorsquare} is already nonnegative. Otherwise multiply the last inequality by the negative number $c-N$ and insert it into \eqref{eq:twovectorsquare} to get
\[
 s[u]\geq
 \frac{\abs{1+U}^2+T(c-N)}{\abs{1+U}^2}\ip y{By}\geq0.
\]
This proves exactly the sufficient implication used in the main theorem; no unnecessary converse is included.
\end{proof}

\node{PO02}{Zero-momentum positivity}{OP02, SE03, OP01}
\begin{lemma}\label{lem:zeropositive}
For small nonzero $g$, $\widehat O_0\geq0$ on $\Q_1$.
\end{lemma}
\begin{proof}
Let $x_0=B_0^{-1}b_0$, $u=\abs q f$, and $a=\ell_{g,0}(f)$. By \eqref{eq:unnormalizedform} and \eqref{eq:saturation},
\[
 \widehat O_0[f]
 =\ip{u+ax_0}{B_0(u+ax_0)}+(c_0-N_0)\abs a^2
 =\ip{u+ax_0}{B_0(u+ax_0)}\geq0.
\]
The scalar $a$ is a finite form-domain pairing, and $u,x_0$ are Hilbert-space vectors. In particular no use of $t_0$ as an $L^2$ vector occurs in this argument.
\end{proof}

\node{PO03}{Uniform nonzero-momentum positivity}{OP02, PO01, PO02, WE06, SE03, KE02, AN03, TR01}
\begin{lemma}\label{lem:nonzeropositive}
There is $g_{\mathrm{pos}}>0$ such that for $0<\abs g<g_{\mathrm{pos}}$ and $0<\abs P\leq1/4$, $\widehat O_P\geq0$.
\end{lemma}
\begin{proof}
Choose a common $C\geq1$ at least as large as the constants in \eqref{eq:cbound} and \eqref{eq:scalarbounds}. By saturation,
\[
 \abs{c_P-N_P}\leq C(\abs P^2+g^2\abs P).
\]
For $Cs^2\leq1/4$, $\abs{1+U_P}\geq1-Cs^2\geq3/4$, so in particular $1+U_P\ne0$. With $r=\abs P>0$,
\begin{align}
 \abs{1+U_P}^2+T_P(c_P-N_P)
 &\geq(1-Cs^2)^2-C^2(s^2r+s^4)\notag\\
 &\geq(1-Cs^2)^2-C^2(s^2/4+s^4).
 \label{eq:finalscalar}
\end{align}
For example, $s\leq1/(4C)$ makes the first term at least $9/16$ and the subtracted quantity at most $5/256<1/4$. Thus the expression is at least $5/16>0$. Apply Lemma~\ref{lem:rankpositivity} to $B_P,b_P,t_P,c_P$ and then use \eqref{eq:unnormalizedform}, where $\ell_{g,P}(f)=\ip{t_P}{u_f}$.
\end{proof}

\begin{corollary}\label{cor:thresholdpositive}
For sufficiently small nonzero $g$, $O_P^{(0)}\geq0$ for all $\abs P\leq1/4$.
\end{corollary}
\begin{proof}
Use Lemma~\ref{lem:zeropositive} at zero momentum and Lemma~\ref{lem:nonzeropositive} otherwise. Translation is an isomorphism of the form domain, so the same assertion holds for the untranslated forms.
\end{proof}

\section{An explicit finite smallness ledger}\label{sec:constants}
\node{SM01}{One coupling cutoff with all constants chosen in order}{FB03, RS01, RS02, KE02, AN03, OP01, WE01, WE02, WE03, WE04, WE06, PO03}
This section turns every use of ``sufficiently small'' into one ordered finite selection. Constants in intermediate estimates are positive real witnesses supplied by their already proved inequalities; no constant is allowed to depend on the particular $g$ or $P$ later quantified in the theorem.

\paragraph{First range: forms, source root, and auxiliary inverses.}
Take the already fixed constants $C_E=2C_L$, $C_1=1+C_L$, $C_2=4C_L$, $C_\Phi=2\max\{1,C_L\}$, $C_R=3C_\Phi$, and $C_D=16C_\Phi^2$ from FB02--FB03, SP03, and RS01. The last constant controls the one-phonon Schur correction and must not be omitted. Put
\[
 s_{\rm pre}=\min\left\{\frac14,\frac1{8(1+C_1)},
 \frac1{\sqrt{8(1+C_E)}},\frac1{\sqrt{8(1+C_2)}},
 \frac1{2(1+C_R)},\frac1{\sqrt{2(1+C_D)}}\right\}.
\]
Every entry is positive. For $0\leq s\leq s_{\rm pre}$ we have
$s\leq1/4$, $C_1s\leq1/8$, $C_Es^2\leq1/8$, $C_2s^2\leq1/8$, $C_Rs\leq1/2$, and $C_Ds^2\leq1/2$. These give the root, high-sector coercivity, auxiliary inverses, and the Schur graph lower bound $a=b=1/2$. For the entire required one-phonon window,
\[
 1+e-|P|^2-\varepsilon
 \geq1-\tfrac18-\tfrac1{16}-\tfrac18=\tfrac{11}{16}>\tfrac12.
\]
Derivatives and kernel estimates are established uniformly on this fixed preliminary range, so their constants are now fixed before another coupling is chosen.

\paragraph{Second range: the two bounded inverses.}
Let $I\geq1$ be a single constant for all three bounds in AN03. Let $K\geq1$ be a single constant for the six bounds in KE02, including the derivative bound. Choose each as the maximum of one and the finitely many constants in those lemmas. Let $C_m\geq0$ satisfy $\|m-1\|_\infty\leq C_m s^2$, as proved by RS03. The dominated-kernel proof gives $\|K_P\|\leq KI$, so put $C_B=C_m+KI$. WE01 has operator bound $C_W=KI$, and $\|\chi_P\|_\infty\leq2K$, since $\omega+|P|\leq5/4<2$. Define
\[
 s_{\rm aux}=\min\left\{\frac{s_{\rm pre}}2,
 \frac1{2\sqrt{1+C_m}},\frac1{\sqrt{2(1+C_B)}},
 \sqrt{\frac3{8(1+KI)}}\right\}.
\]
For $s\leq s_{\rm aux}$, these choices give $m\geq3/4$, $B_P\geq I/2$, and
$\|s^2M_{m^{-1}}\mathscr L_P\|\leq(4/3)s^2KI\leq1/2$.
Consequently the \emph{normalized} Neumann inverse obeys
\[
 \|(I+g^2M_{m^{-1}}\mathscr L_P)^{-1}\|\leq2.
\]
The inverse for the original left-hand side is
$(M_m+g^2\mathscr L_P)^{-1}=(I+g^2M_{m^{-1}}\mathscr L_P)^{-1}M_{m^{-1}}$, so its norm is at most $2(4/3)=8/3$, not two. The constants below already use this factor.

\paragraph{Explicit comparison constants, not a new choice of range.}
The following constants implement the inequalities in WE02--WE06 using only $I,K$:
\begin{align*}
 C_\eta&=16K/3, & C_J&=KIC_\eta,\\
 C_0&=\tfrac43(K+C_J), & C_\Delta&=\tfrac83(K+2KIC_0),\\
 C_{\rm sharp}&=C_0+C_\Delta, & C_U&=I C_{\rm sharp},\\
 C_T&=2I, & C_N&=KI(C_\eta+C_\Delta+C_0).
\end{align*}
Indeed WE02 gives $\|\zeta_P\|_\infty\leq2(4/3)(2K)=C_\eta$. In WE03 the symmetric-difference and ordinary kernel bounds give $|J(q)|\leq C_J\omega(q)^2$; since $s\leq1$, solving for $\eta_0$ gives $C_0$. In WE04 the two terms in $(\mathcal L_P-\mathcal L_0)\eta_0$ have total bound $2KIC_0|P|\omega(q)$. Multiplying the remaining right side by the inverse norm $8/3$ gives $C_\Delta$. Adding zero-momentum and difference bounds gives $C_{\rm sharp}$. Finally the three products in WE06 give exactly the listed $C_U,C_T,C_N$. These constants are fixed on the auxiliary range and do not depend on a subsequently chosen $s$.

\paragraph{Complete guard ledger.}
The following are the precise meanings of the small-coupling qualifiers in the preceding statements. These are finite conjunctions, not additional unproved assumptions on the final theorem.
\begin{center}\small
\begin{tabular}{p{.23\textwidth}p{.65\textwidth}}
\toprule Nodes & Sufficient numerical hypotheses in addition to their displayed domains\\\midrule
FB03 & $s\leq1/4$, and $C_2s^2\leq1/8$ for its two-sector assertion.\\
SP01 & $s\leq1/4$, $C_1s\leq1/8$, $C_Es^2\leq1/8$.\\
SP03 & The preceding bounds, $C_2s^2\leq1/8$, $C_Ds^2\leq1/2$, and $0\leq\varepsilon\leq C_1s$.\\
RS01--KE02 & $s\leq s_{\rm pre}$; reference $e\in[-1/4,0]$ and then $e=E_g(0)$.\\
OP01--SE03 & $s\leq s_{\rm aux}$, and $g\ne0$ only for the infrared contradiction.\\
WE01--WE06 & $s\leq s_{\rm aux}$; $P\ne0$ only for the quantities involving $t_P$.\\
PO01 & Its stated abstract algebraic hypotheses, with no model or coupling assumption.\\
PO02--MN01 & $0<s<g_*$ as selected below.\\
MN02 & $g=0$, treated directly without any division by $g$.\\
\bottomrule
\end{tabular}\normalsize\end{center}
The preliminary constants were obtained from form estimates that hold on a fixed interval of couplings and, for inverse derivatives, a fixed interval of reference energies. Hence choosing their maximum does not depend on a later cutoff. The constants of the weighted equation are chosen only after its two inverse bounds are fixed, as shown explicitly above.

\paragraph{Final cutoff and scalar positivity.}
Put
\[
 C=\max\{1,K,C_U,C_T,C_N\},\qquad
 \boxed{g_*:=\min\{s_{\rm aux}/2,\,1/(4C)\}>0.}
 \tag{S}\label{eq:gstar}
\]
For $s<g_*$, $Cs^2\leq1/16$ and hence $|1+U_P|\geq15/16$. Also
\[
 C^2(s^2/4+s^4)\leq1/64+1/256=5/256.
\]
Thus the lower bound in \eqref{eq:finalscalar} is at least
$(15/16)^2-5/256=220/256>0$. In particular it is stronger than the coarser bound used in PO03. All previous form, inverse, derivative, weighted-equation and positivity requirements hold at once. No infimum over an uncontrolled family of positive cutoffs is being taken; every minimum displayed here is finite and has strictly positive entries.

\section{Closing the spectral argument, including endpoints and zero coupling}\label{sec:assembly}
\node{MN01}{Nonzero coupling: all subthreshold spectral points}{SP01, SP03, PO02, PO03, SM01, FA05, GW07}
\begin{theorem}\label{thm:nonzeromain}
For sufficiently small nonzero $g$ and all $\abs P\leq1/4$,
\[
 \sigR(H_g(P))\cap(-\infty,e+1)=\{e_P\},
 \qquad e_P=E_g(P),\quad e=E_g(0).
\]
\end{theorem}
\begin{proof}
Choose the coupling so that all previous estimates hold, $L_s>1/2$, and $e+1>1/2$. The constructed eigenvalue satisfies
\[
 e_P\leq\abs P^2\leq1/16<e+1,
\]
so the singleton is contained in the left side.

Let $\lambda<e+1$ with $\lambda\ne e_P$. If $\lambda<L_s$, Lemma~\ref{lem:vacuumroot} puts $\lambda$ in the resolvent set. If $\lambda\geq L_s$, let $\varepsilon=e+1-\lambda>0$. Since $e\leq0$,
\[
 \varepsilon\leq e+1-L_s=e+C_1s\leq C_1s.
\]
Corollary~\ref{cor:thresholdpositive} and Lemma~\ref{lem:nearresolvent} again put $\lambda$ in the resolvent set. The equality case $\lambda=L_s$ is included in this second case, so there is no uncovered endpoint. This proves the reverse containment.
\end{proof}
The two cases need not be described as overlapping open intervals. They are an exact exhaustive split, with the boundary assigned to the near-threshold argument. The threshold energy itself is excluded by the open interval in the theorem and is not analyzed.

\node{MN02}{Free coupling}{FD02, FA05, GW07}
\begin{lemma}\label{lem:free}
For $g=0$ and $\abs P\leq1/4$,
\[
 E_0(P)=\abs P^2,\qquad E_0(0)=0,\qquad
 \sigR(H_0(P))\cap(-\infty,1)=\{\abs P^2\}.
\]
\end{lemma}
\begin{proof}
The vacuum is an operator-domain eigenvector with eigenvalue $\abs P^2$. On the nonvacuum sector, the free multiplication operator is at least $1$ because $N\geq1$ and the kinetic term is nonnegative. For $\lambda<1$, its inverse is the multiplication operator
\[
 ((\abs{K_n-P}^2+n)-\lambda)^{-1},\qquad n\geq1,
\]
with norm at most $(1-\lambda)^{-1}$ and range in the operator domain. The latter follows because the multiplier $(\abs{K_n-P}^2+n)/((\abs{K_n-P}^2+n)-\lambda)$ is bounded for fixed $\lambda<1$. If $\lambda\ne\abs P^2$, invert the vacuum scalar as well. This constructs the full inverse. Since $\abs P^2\leq1/16<1$, the stated energy and spectrum follow. This proof needs neither the exact spectrum of the higher sectors nor an essential-spectrum computation.
\end{proof}

\node{MN03}{The stated theorem for the actual fiber operators}{MN01, MN02, SM01, FA05, GW06, GW07}
\begin{proof}[Proof of Theorem~\ref{thm:target}]
Choose the single positive coupling cutoff in Section~\ref{sec:constants}. Split into $g=0$ and $g\ne0$. Use Lemma~\ref{lem:free} in the first case and Theorem~\ref{thm:nonzeromain} in the second. All arguments used $s=\abs g$, not the sign of $g$. The operators are precisely the GW form operators fixed in GW06 on the canonical GW carrier. GW07 identifies the real bounded-inverse definition with the real slice of the full complex spectrum through the existing GW predicate. Thus the real-spectrum formulation is exactly the stated equality under the real embedding. RE01 is not needed for this canonical realization. Eigenvector existence was constructed in both cases.
\end{proof}

\clearpage
\section{Dependency architecture and source modules}\label{sec:registry}
The registry below records the explicit nodes of this revision. The accompanying JSON records their direct dependencies and a topological ordering. Library leaves have the exact statements and source locations in Section~\ref{sec:mathlib}. RE01 is an optional source-transport branch, not an assumed premise of the canonical theorem. Form-cutoff convergence and limiting estimates are deliberately separate nodes; WE01 explicitly consumes OP01's lower bound on $m$.

\small\begin{longtable}{p{.09\textwidth}p{.49\textwidth}p{.34\textwidth}}\toprule Node & Mathematical contract & Implementation action \\\midrule\endhead
\hyperlink{node:GW01}{GW01} & Canonical sectors, scalar vacuum, one-particle inclusion and projections & Local proof in \path{GW01.lean}.\\
\hyperlink{node:GW02}{GW02} & GW currying in the correct direction, including symmetric sections & Local proof in \path{GW02.lean}.\\
\hyperlink{node:GW03}{GW03} & Number graph coordinates and the assembled annihilation field & Local proof in \path{GW03.lean}.\\
\hyperlink{node:GW04}{GW04} & Restrictions of the maximal GW smeared operators, not replacement domains & Local proof in \path{GW04.lean}.\\
\hyperlink{node:GW05}{GW05} & Signed coupling and closed ultraviolet balls & Local proof in \path{GW05.lean}.\\
\hyperlink{node:EN01}{EN01} & The complete fiber energy carrier and its embedding into GW Fock space & Local proof in \path{EN01.lean}.\\
\hyperlink{node:EN02}{EN02} & Energy norm changes, restrictions, and the dual change-of-coordinates rule & Local proof in \path{EN02.lean}.\\
\hyperlink{node:MS01}{MS01} & Rectangle density with its generating-set argument & Reuse GW; no new construction.\\
\hyperlink{node:MS02}{MS02} & The scalar-product to Bochner-$L^2$ identification & Reuse GW; no new construction.\\
\hyperlink{node:MS03}{MS03} & Countable Hilbert sums and fields assembled from components & Reuse GW; no new construction.\\
\hyperlink{node:FD01}{FD01} & Canonical symmetric sectors and their projections & Reuse GW; no new construction.\\
\hyperlink{node:FD02}{FD02} & Weighted domains, dense embeddings and multiplication operators & Wrap the fixed GW adapter.\\
\hyperlink{node:BA01}{BA01} & Riesz coordinates for conjugate-dual sources and forms & Local proof in \path{BA01.lean}.\\
\hyperlink{node:BA02}{BA02} & Geometric inverses, continuity, and order on dual sources & Reuse ML-COMPLETE and ML-SERIES; retain the inverse-order and dual-source specialization.\\
\hyperlink{node:BA03}{BA03} & Mixed form bounds from diagonal bounds & Local proof in \path{BA03.lean}.\\
\hyperlink{node:DF01}{DF01} & Domain and spectrum definitions on the fixed partial-operator interface & Local proof in \path{DF01.lean}.\\
\hyperlink{node:FA00}{FA00} & Construction and uniqueness of the associated self-adjoint operator & Reuse GW; no new construction.\\
\hyperlink{node:FC01}{FC01} & A bounded-momentum algebraic core & Local proof in \path{FC01.lean}.\\
\hyperlink{node:FC02}{FC02} & The annihilation map into a Bochner $L^2$ space & Wrap the fixed GW adapter.\\
\hyperlink{node:FC06}{FC06} & Smeared adjointness and the creation norm, with all indices & Wrap the fixed GW adapter.\\
\hyperlink{node:FC03}{FC03} & Smeared operators and their adjoints & Wrap the fixed GW adapter.\\
\hyperlink{node:FC04}{FC04} & Energy control for cutoff annihilation fields & Local proof in \path{FC04.lean}.\\
\hyperlink{node:FC07}{FC07} & Field commutations and the compressed free shift & Local proof in \path{FC07.lean}.\\
\hyperlink{node:FC08}{FC08} & Measurable bounded operator families acting on fields & Local proof in \path{FC08.lean}.\\
\hyperlink{node:FC09}{FC09} & A common weak equation from an $L^2$ residual & Local proof in \path{FC09.lean}.\\
\hyperlink{node:FC05}{FC05} & The weak-equation extraction interface used in pull-through & Optional interface outside MN03\textquotesingle s declared closure; no initial implementation ticket.\\
\hyperlink{node:FA01}{FA01} & Energy-space inversion & Local proof in \path{FA01.lean}.\\
\hyperlink{node:FA02}{FA02} & Coercivity gives an energy-space inverse & Reuse GW-Coercive on the fixed energy Hilbert space; only a signature/inverse-identity adapter.\\
\hyperlink{node:FA03}{FA03} & Two estimates imply graph-norm coercivity & Local proof in \path{FA03.lean}.\\
\hyperlink{node:FA04}{FA04} & Schur elimination on energy spaces & Local proof in \path{FA04.lean}.\\
\hyperlink{node:FA05}{FA05} & Kernel obstruction, uniqueness of inverses, and the real spectrum & Local proof in \path{FA05.lean}.\\
\hyperlink{node:GW07}{GW07} & Interpreting the spectrum by GW's actual inverse predicate & Local proof in \path{GW07.lean}.\\
\hyperlink{node:AN00}{AN00} & Translations, scalar dilations and positive finite ball volume & Reuse GW-Model volume clauses; check the affine measure adapters.\\
\hyperlink{node:AN01}{AN01} & Power integrals and infrared nonintegrability & Local proof in \path{AN01.lean}.\\
\hyperlink{node:AN02}{AN02} & Uniform shifted source norm & Local proof in \path{AN02.lean}.\\
\hyperlink{node:LD01}{LD01} & Weighted source functionals and their exact convergence norm & Local proof in \path{LD01.lean}.\\
\hyperlink{node:TR01}{TR01} & Momentum translation as a specific form-domain isomorphism & Local proof in \path{TR01.lean}.\\
\hyperlink{node:AN03}{AN03} & Moving singularities & Local proof in \path{AN03.lean}.\\
\hyperlink{node:FB01}{FB01} & Finite-cutoff Lieb--Yamazaki inequalities, before any singular form exists & Local proof in \path{FB01.lean}.\\
\hyperlink{node:FB02A}{FB02A} & Norm construction of the singular interaction and its limiting estimates & Local proof in \path{FB02A.lean}.\\
\hyperlink{node:GW06}{GW06} & The precise coercive-form record defining the fiber Hamiltonian & Local proof in \path{GW06.lean}.\\
\hyperlink{node:FB02}{FB02} & Closed fiber forms, associated operators and selection rules & Local proof in \path{FB02.lean}.\\
\hyperlink{node:SY01}{SY01} & Conjugation and reflection on forms, sources and weak inverses & Local proof in \path{SY01.lean}.\\
\hyperlink{node:RE01}{RE01} & Unitary transport of the entire form construction & Optional transport, not required by MN03.\\
\hyperlink{node:FB03}{FB03} & Uniform fiber and sector bounds & Local proof in \path{FB03.lean}.\\
\hyperlink{node:SP00}{SP00} & Monotonicity of a scalar source resolvent & Local proof in \path{SP00.lean}.\\
\hyperlink{node:SP01}{SP01} & Vacuum Schur function & Local proof in \path{SP01.lean}.\\
\hyperlink{node:SP03}{SP03} & One-phonon elimination and monotonicity & Local proof in \path{SP03.lean}.\\
\hyperlink{node:RS01}{RS01} & Uniform auxiliary resolvents & Local proof in \path{RS01.lean}.\\
\hyperlink{node:CA01}{CA01} & Norm differentiation, product rules and integrated remainders & Reuse ML-DIFF and ML-FTC; prove the stated quantitative recurrence/remainders.\\
\hyperlink{node:CA02}{CA02} & Uniform differentiation of the free inverse-square-root weights & Local proof in \path{CA02.lean}.\\
\hyperlink{node:CA03}{CA03} & Finite differences with small-distance weights & Local proof in \path{CA03.lean}.\\
\hyperlink{node:RS02}{RS02} & Norm derivatives through order three & Local proof in \path{RS02.lean}.\\
\hyperlink{node:RS03}{RS03} & The scalar self-energy has a small Hessian & Local proof in \path{RS03.lean}.\\
\hyperlink{node:UV01}{UV01} & Fixed-energy cutoff convergence in the correct norms & Local proof in \path{UV01.lean}.\\
\hyperlink{node:CV01}{CV01} & A norm-by-norm limit ledger for the ultraviolet approximation & Local proof in \path{CV01.lean}.\\
\hyperlink{node:UV02}{UV02} & The cutoff operator domain equals the free operator domain & Local proof in \path{UV02.lean}.\\
\hyperlink{node:PT01}{PT01} & Annihilation pull-through proved through an energy-valued residual & Local proof in \path{PT01.lean}.\\
\hyperlink{node:PT02}{PT02} & Pull-through for the singular one-particle form source & Local proof in \path{PT02.lean}.\\
\hyperlink{node:KA01}{KA01} & The finite-cutoff one-to-high map and its physical adjoint & Local proof in \path{KA01.lean}.\\
\hyperlink{node:PT03}{PT03} & A scalar, integrated proof of the threshold kernel & Local proof in \path{PT03.lean}.\\
\hyperlink{node:KE01}{KE01} & Kernel derivatives and reflection & Local proof in \path{KE01.lean}.\\
\hyperlink{node:KE02}{KE02} & Quantitative cancellation bounds & Local proof in \path{KE02.lean}.\\
\hyperlink{node:OP00}{OP00} & A bounded integral operator from a product majorant & Local proof in \path{OP00.lean}.\\
\hyperlink{node:OP01}{OP01} & Construction of the bounded operators & Local proof in \path{OP01.lean}.\\
\hyperlink{node:OP02}{OP02} & Form decomposition valid also at zero momentum & Local proof in \path{OP02.lean}.\\
\hyperlink{node:OP03}{OP03} & The polarized unnormalized decomposition & Local proof in \path{OP03.lean}.\\
\hyperlink{node:SE01}{SE01} & The ground source field and sector equations & Local proof in \path{SE01.lean}.\\
\hyperlink{node:SE02}{SE02} & The one-particle source identity with its cutoff-dual passage & Local proof in \path{SE02.lean}.\\
\hyperlink{node:SE03}{SE03} & Exact saturation forced by infrared nonintegrability & Local proof in \path{SE03.lean}.\\
\hyperlink{node:WE01}{WE01} & A bounded $L^\infty$ kernel operator & Local proof in \path{WE01.lean}.\\
\hyperlink{node:WE02}{WE02} & Existence, uniqueness, and the representative equation & Local proof in \path{WE02.lean}.\\
\hyperlink{node:WE03}{WE03} & Quadratic vanishing by symmetric second differences & Local proof in \path{WE03.lean}.\\
\hyperlink{node:WE04}{WE04} & Lipschitz comparison in momentum & Local proof in \path{WE04.lean}.\\
\hyperlink{node:WE05}{WE05} & Identification with the Hilbert-space inverse & Local proof in \path{WE05.lean}.\\
\hyperlink{node:WE06}{WE06} & The three scalar estimates & Local proof in \path{WE06.lean}.\\
\hyperlink{node:PO01}{PO01} & An unnormalized two-vector positivity lemma & Local proof in \path{PO01.lean}.\\
\hyperlink{node:PO02}{PO02} & Zero-momentum positivity & Local proof in \path{PO02.lean}.\\
\hyperlink{node:PO03}{PO03} & Uniform nonzero-momentum positivity & Local proof in \path{PO03.lean}.\\
\hyperlink{node:SM01}{SM01} & One coupling cutoff with all constants chosen in order & Local proof in \path{SM01.lean}.\\
\hyperlink{node:MN01}{MN01} & Nonzero coupling: all subthreshold spectral points & Local proof in \path{MN01.lean}.\\
\hyperlink{node:MN02}{MN02} & Free coupling & Local proof in \path{MN02.lean}.\\
\hyperlink{node:MN03}{MN03} & The stated theorem for the actual fiber operators & Local proof in \path{MN03.lean}.\\
\bottomrule\end{longtable}\normalsize
\paragraph{A cycle-free file plan.}
Retain the complete original \path{GW/} and \path{GW.lean} unchanged. Each genuinely new node has its own file \path{Froehlich/Nodes/<ID>.lean}; it imports the files of its earlier local prerequisites and the exact GW modules required by its source leaves. The companion \path{metadata/dependencies.json} gives these paths and a topological order; it is regenerated from the node headers, not independent mathematical authority. Keeping one node per file initially avoids introducing a module cycle by grouping an early adapter and a late Hamiltonian definition into the same file. Source-reuse-only nodes require no new construction: their mathematical explanation can be omitted from the implementation, but the actual specialization must still compile before a consumer treats it as proved. FA02 is now a direct GW reuse; BA02, AN00, and CA01 combine library reuse with their stated local specializations. The declared target closure has 81 of the 83 nodes; FC05 and RE01 are outside it and receive no initial proof ticket. The wrappers FD02, FC02, FC03 and FC06 use EN01--EN02 and GW01--GW04; they do not introduce another carrier or another field operator. RE01 is retained as a proved optional transport theorem and is not on the canonical target's dependency closure.

\paragraph{Reading order and the two branches.}
The ground root SP01 precedes every use of $e=E_g(0)$. SP03 states an implication from threshold positivity to a physical resolvent; it does not assume the positivity theorem has already been proved. PT03 may therefore use SP03's definition without a cycle. SE01--SE02 can run alongside the kernel/operator branch, but SE03 explicitly waits for OP01 and OP03. Only after those interfaces exist can the cancellation and weighted-comparison branches proceed independently toward PO02--PO03. SM01 chooses a single cutoff after the finite family of uniform bounds has been established. MN01--MN03 perform the exact spectral partition and zero-coupling branch.

\paragraph{Suggested local source grouping.}
Use separate local modules for measure/field assembly (MS), weighted domains (FD), bounded forms and associated operators (BA/FA), sector fields (FC), singular integrals (AN), fiber forms (FB), spectral reductions (SP), parameter calculus (CA/RS), cutoff limits (UV), integrated pull-through and the kernel (PT/KE/OP), cancellation (SE), weighted comparison (WE), positivity (PO), and final assembly (SM/MN). These are organizational names, not Lean declarations or external packages. Retained GW source is an implementation of selected foundational nodes, not another assumed analytic layer.



\subsection{Interface gates and the two-worker schedule}\label{sec:execution}
This subsection specializes the companion \path{Froehlich-Codex-instructions-v5.3.txt}. It is an execution contract, not an additional mathematical hypothesis. The manuscript is authoritative for mathematics; the instruction file is authoritative for job management. A source lookup, a proved paragraph here, or membership in the planning DAG does not make a new Lean declaration verified.

\paragraph{Reuse overrides expository reproof.}
For a stated general result, use a faithful existing GW or pinned-mathlib declaration with its hypotheses discharged on the prescribed carriers. Prove only a necessary specialization or compatibility identity. This does not authorize a new carrier, a different operator domain, additional premises, or weakened constants. For an external-library port, include the licensed proof and its necessary source dependencies locally, record an immutable upstream reference and attribution, and rebuild against the unchanged pin. An external theorem header or literature citation never closes a dependency.

\paragraph{Eight interface checks, staged when their objects exist.}
Each check is a small actual application, not just a name lookup. Run it before its first dependent branch, rather than blocking the whole project until every later object is constructed.
\begin{enumerate}
\item GW01--GW02: scalar vacuum, one-particle normalization, and the direction and range of the existing product/currying isometry.
\item EN01--EN02: complete energy coordinates, dense physical embedding, and the specific norm-changing maps; no ambient-norm domain subtype substitutes for the energy norm.
\item GW03--GW04: the original smeared maps and domains, and the new momentum-indexed annihilation field with its almost-everywhere and number-norm identities.
\item BA01: conjugation of the linear source into a conjugate-dual load and the real-part/self-adjoint convention for its Riesz representative.
\item BA02/FA02: the bounded-map completeness instance, noncommutative Neumann inverse, exact GW coercive inverse, both inverse identities, and the norm bounds.
\item FA00/GW06: the actual \path{formOperator}, its domain equivalence and the chosen fiber form data; no full-space Hamiltonian theorem substitutes for this fiber construction.
\item DF01/GW07: the retained bounded two-sided inverse predicate and the real-spectrum complement, with the actual shifted partial operator and domain.
\item CA01: real-parameter differentiation on bounded complex-linear maps, the units bridge for \path{Ring.inverse}, and the multiplication order in the derivative.
\end{enumerate}
Inspect all binders and meaning-carrying definitions for these interfaces once, and recheck only after a relevant change. Generic coercive/inverse tests can run early; the actual fiber form and spectrum tests wait for their defining nodes. Record \contract{SOURCE-MATCHED}, \contract{APPLICATION-CHECKED}, and \contract{VERIFIED} as distinct evidence, not synonyms.

\paragraph{Three roles; one compiler queue.}
The coordinator approves the common definitions and signatures, assigns disjoint owned files, integrates candidates, and alone edits the authoritative ledger and shared imports. Workers A and B develop separate ready nodes or perform bounded read-only searches/reviews. They do not spawn further agents. For the two-core, 8\,GB laptop, only the coordinator launches Lean/Lake checks, one active expensive job globally. A submitted file and its imported dependency revision are frozen until the job exits; another worker may edit only files outside that build's import closure. Do not mistake one Lake process for a guaranteed limit on its internal Lean jobs: inspect the pinned client's supported concurrency controls and use conservative settings. No obsolete revision runs beside its replacement.

\begin{longtable}{p{.22\textwidth}p{.33\textwidth}p{.36\textwidth}}
\toprule Stage & Worker A & Worker B / readiness rule\\\midrule\endhead
Six-node pilot & GW01, then EN01 and GW02 & BA01, then BA02 and FA02. Accept all named source leaves; respect the actual DAG, not the row order alone.\\
Shared foundations & GW03--GW05, EN02 and ready FC nodes & Ready AN, LD/TR, FA, CA nodes; OP00 and PO01 when their prerequisites are accepted.\\
After auxiliary inverses & UV01/CV01, UV02, then PT01--PT03 & RS02--RS03 after their CA prerequisites. KA01 is accepted before its PT03 consumer.\\
After PT03 & SE01--SE02 & KE01--KE02, then OP01--OP03, with OP00 accepted.\\
After operator join & SE03, then PO02 & WE01--WE06 in dependency order. SE03 cannot bypass OP03.\\
Final integration & \multicolumn{2}{p{.70\textwidth}}{Coordinator assembles PO03, SM01, MN01--MN03; MN02 is the independent free-coupling case when ready.}\\
\bottomrule
\end{longtable}
These are candidate work pools, not licenses to ignore a missing direct prerequisite. If no second independent proof task exists, use the second worker for a focused review or leave it idle. On repeated identical failures, inspect a short reproducer and profile imports, statement elaboration, and proof elaboration before another full attempt. A disposable diagnostic proof hole must never enter production imports.

\paragraph{Verification without a second bookkeeping project.}
Keep one compact row per source node, not a manual registry of inherited GW helpers. Local acceptance checks a changed module, its public statement and transitive axioms. Integration checks affected consumers at genuine joins or changes to shared definitions. Final acceptance builds every retained production module and checks the fully elaborated target with its actual defining objects and permitted axioms. No previous evidence survives a changed dependency fingerprint. Optional nodes may remain documented but unimplemented; report 81 target-required source nodes separately from the complete 83-node catalogue. A checkpoint is not a stopping condition, and an audit is not an additional main theorem.

\section{Exact library leaves and source references}\label{sec:mathlib}
\paragraph{Reference revision and scope.}
This leaf list originated in the version 3 documentation audit, whose source link resolved to commit
\texttt{af6dcfbf23867d9df7ee4835e28567fe509a2677}. The actual retained GW source pins mathlib to \texttt{9bac2d0692944a775113f7c130fd5e5b4608a854}. The companion reference index records occurrences of these APIs in the uploaded GW source separately from documentation-only evidence; neither is described as a fresh build at that pin. The companion \path{metadata/library-contracts.json} records the selected rechecked modules and declarations separately; the inherited catalogue below remains the reference for the other leaf groups. These leaves are theorems about ordinary Hilbert spaces, $L^p$ spaces and integrals, not assertions that the physics operators or pull-through formulas already exist in mathlib. The local proofs specify the exact mathematical specialization when the library statement is more general.

\subsection{Verified GW leaves: exact source, not paper-level assumptions}
The following source groups are existing proved declarations in the included release. Their names are leaves in the dependency graph. An adapter node states the specialization and any change of coordinates; the indexed source gives the full generic theorem and its retained proof. The complete source file, not a truncated theorem header, fixes all implicit variables and instances.
\paragraph{\texttt{GW-Sectors}. Existing scalar-vacuum symmetric sectors and Hilbert sum}\label{leaf:GW-Sectors}
Use the exact declarations under this group in Section~\ref{sec:gwverify}. Source files: \path{GW/Conventions.lean}, \path{GW/Fock/Sectors.lean}.
\paragraph{\texttt{GW-Coordinates}. Measure-preserving coordinate splits and sector inclusions}\label{leaf:GW-Coordinates}
Use the exact declarations under this group in Section~\ref{sec:gwverify}. Source files: \path{GW/Fock/Coordinates.lean}.
\paragraph{\texttt{GW-Product}. Bochner products, amplification, scalar contraction and density}\label{leaf:GW-Product}
Use the exact declarations under this group in Section~\ref{sec:gwverify}. Source files: \path{GW/Fock/FockAmplification.lean}.
\paragraph{\texttt{GW-Generalized}. Fixed-sector projection and the existing ordinary field operators}\label{leaf:GW-Generalized}
Use the exact declarations under this group in Section~\ref{sec:gwverify}. Source files: \path{GW/Fock/FockGeneralized.lean}.
\paragraph{\texttt{GW-Number}. Maximal number multipliers and their square-root domains}\label{leaf:GW-Number}
Use the exact declarations under this group in Section~\ref{sec:gwverify}. Source files: \path{GW/Fock/Number.lean}.
\paragraph{\texttt{GW-Bochner}. Existing sector-field reconstruction, with a common exceptional set}\label{leaf:GW-Bochner}
Use the exact declarations under this group in Section~\ref{sec:gwverify}. Source files: \path{GW/Fock/Bochner.lean}.
\paragraph{\texttt{GW-Fields}. Existing root inputs and the field on the number-form domain}\label{leaf:GW-Fields}
Use the exact declarations under this group in Section~\ref{sec:gwverify}. Source files: \path{GW/Fock/FockFieldDefinition.lean}.
\paragraph{\texttt{GW-Bounds}. Existing smeared creation and annihilation estimates}\label{leaf:GW-Bounds}
Use the exact declarations under this group in Section~\ref{sec:gwverify}. Source files: \path{GW/Fock/FockGeneralizedBounds.lean}.
\paragraph{\texttt{GW-Adjoints}. Adjoint pairing of maximal sector fields}\label{leaf:GW-Adjoints}
Use the exact declarations under this group in Section~\ref{sec:gwverify}. Source files: \path{GW/Fock/FockAdjoints.lean}.
\paragraph{\texttt{GW-CCR}. Existing smeared norm identity and canonical commutations}\label{leaf:GW-CCR}
Use the exact declarations under this group in Section~\ref{sec:gwverify}. Source files: \path{GW/Fock/FockCcr.lean}.
\paragraph{\texttt{GW-Weighted}. Weighted coordinate spaces and their physical embeddings}\label{leaf:GW-Weighted}
Use the exact declarations under this group in Section~\ref{sec:gwverify}. Source files: \path{GW/Spaces/WeightedL2.lean}.
\paragraph{\texttt{GW-Multiplier}. Bounded and maximal multiplication operators}\label{leaf:GW-Multiplier}
Use the exact declarations under this group in Section~\ref{sec:gwverify}. Source files: \path{GW/Analysis/L2Quotient.lean}, \path{GW/Analysis/RealMultiplier.lean}.
\paragraph{\texttt{GW-Model}. Kernel normalization, cutoff membership, and Euclidean radial integrals}\label{leaf:GW-Model}
Use the exact declarations under this group in Section~\ref{sec:gwverify}. Source files: \path{GW/Model/Coefficient.lean}, \path{GW/Model/ModelFrohlichEstimates.lean}, \path{GW/Model/ModelPowerIntegrals.lean}.
\paragraph{\texttt{GW-Form}. Existing coercive-form record, associated operator and exact domain}\label{leaf:GW-Form}
Use the exact declarations under this group in Section~\ref{sec:gwverify}. Source files: \path{GW/Operators/FormRepresentation.lean}, \path{GW/Operators/Polarization.lean}.
\paragraph{\texttt{GW-Coercive}. Bounded coercive inversion on complete complex Hilbert spaces}\label{leaf:GW-Coercive}
Use \path{GW/Operators/CoerciveBounded.lean}: \path{GW.Operators.coercive_bijective}, \path{GW.Operators.coerciveInverse}, \path{GW.Operators.coerciveInverse_left}, \path{GW.Operators.coerciveInverse_right}, and \path{GW.Operators.coerciveInverse_norm_le}. Inputs are a bounded complex-linear $T$, $T^*=T$, $c>0$, and $c\|u\|^2\leq\Re\langle u,Tu\rangle$ for all $u$, on a complete complex Hilbert space. Output is the actual bounded two-sided inverse with norm at most $c^{-1}$. FA02 supplies only the specialization; all five declarations and their proof bodies are retained unchanged.
\paragraph{\texttt{GW-Resolvent}. Existing two-sided bounded-resolvent predicate for partial operators}\label{leaf:GW-Resolvent}
Use the exact declarations under this group in Section~\ref{sec:gwverify}. Source files: \path{GW/Operators/ResolventAlgebra.lean}.
\paragraph{\texttt{GW-Mixed}. Bounded sesquilinear forms and their norm limits}\label{leaf:GW-Mixed}
Use the exact declarations under this group in Section~\ref{sec:gwverify}. Source files: \path{GW/Analysis/MixedForms.lean}.
\paragraph{\texttt{GW-Domination}. An already proved vector-valued L2 dominated-convergence theorem}\label{leaf:GW-Domination}
Use the exact declarations under this group in Section~\ref{sec:gwverify}. Source files: \path{GW/Analysis/L2Domination.lean}.
\subsection{The finite leaf list}
\hypertarget{leaf:ML-LP}{}\paragraph{ML-LP: equivalence classes, representatives, and zero norm.}
Use \path{Mathlib/MeasureTheory/Function/LpSpace/Basic}. The checked declarations include
\path{MeasureTheory.MemLp.toLp},
\path{MeasureTheory.MemLp.toLp_congr},
\path{MeasureTheory.Lp.ext},
\path{MeasureTheory.Lp.ext_iff},
\path{MeasureTheory.Lp.stronglyMeasurable},
\path{MeasureTheory.Lp.memLp}, and
\path{MeasureTheory.Lp.eq_zero_iff_ae_eq_zero}.
Their role is exactly: a strongly measurable square-integrable representative defines an $L^2$ class; almost-everywhere equal representatives define the same class; class equality may be proved almost everywhere; a zero class has zero representative almost everywhere. The $L^\infty$ use is the same quotient construction at exponent $\infty$, with essential-supremum norm. None permits pointwise evaluation of an arbitrary class at a prescribed momentum.
\url{https://leanprover-community.github.io/mathlib4_docs/Mathlib/MeasureTheory/Function/LpSpace/Basic.html}.

\hypertarget{leaf:ML-L2}{}\paragraph{ML-L2: Hilbert structure and the integral inner product.}
Use \path{Mathlib/MeasureTheory/Function/L2Space}, specifically
\path{MeasureTheory.L2.innerProductSpace},
\path{MeasureTheory.L2.inner_def},
\path{MeasureTheory.L2.integrable_inner}, and
\path{MeasureTheory.memLp_two_iff_integrable_sq_norm}.
The specialization is exponent two and complete real or complex Hilbert-valued functions. It identifies the $L^2$ inner product with the integral of pointwise inner products and makes the squared norm the integral of the squared pointwise norm. Integrability of each product also follows from Hilbert Cauchy--Schwarz plus the scalar $L^2$ inequality. In all applications here both factors have their $L^2$ membership proved before the inner-product identity is used.
\url{https://leanprover-community.github.io/mathlib4_docs/Mathlib/MeasureTheory/Function/L2Space.html}.

\hypertarget{leaf:ML-SIMPLE}{}\paragraph{ML-SIMPLE: density of simple functions.}
Use \path{Mathlib/MeasureTheory/Function/SimpleFuncDenseLp}, specifically
\path{MeasureTheory.Lp.simpleFunc.dense},
\path{MeasureTheory.Lp.simpleFunc.denseRange},
\path{MeasureTheory.MemLp.exists_simpleFunc_eLpNorm_sub_lt}, and
\path{MeasureTheory.Lp.induction}.
The needed specialization is $p=2<\infty$ (or $p=1$ for an integrable approximation), with complete normed target. It supplies finite simple functions with arbitrarily small $L^p$ error. It does \emph{not} by itself identify product-space sections: MS01--MS02 prove that additional compatibility. In FC09 the underlying measure is the restriction to a bounded Euclidean set, so all indicator tests have finite $L^2$ norm.
\url{https://leanprover-community.github.io/mathlib4_docs/Mathlib/MeasureTheory/Function/SimpleFuncDenseLp.html}.

\hypertarget{leaf:ML-COMPLETE}{}\paragraph{ML-COMPLETE: the Banach property of $L^p$.}
Use \path{Mathlib/MeasureTheory/Function/LpSpace/Complete}, in particular
\path{MeasureTheory.Lp.instCompleteSpace}.
Its range is $1\leq p\leq\infty$ with complete normed target. Thus both Bochner $L^2$ and real $L^\infty$ used here are complete. For bounded maps, reuse \path{ContinuousLinearMap.instCompleteSpace} from \path{Mathlib/Topology/Algebra/Module/Spaces/ContinuousLinearMap.lean}. In the normed setting the source is a sequential topological vector space and the target is complete; its bounded-convergence uniformity agrees with the operator-norm uniformity. Check the actual complex-linear and real-linear specializations, rather than assuming that completeness of a physical Hilbert space makes a domain subtype complete. The Cauchy arguments below explain the construction but do not require a new completeness instance. No completeness of an unstated weighted-supremum space is assumed.
\url{https://leanprover-community.github.io/mathlib4_docs/Mathlib/MeasureTheory/Function/LpSpace/Complete.html}.

\hypertarget{leaf:ML-HSUM}{}\paragraph{ML-HSUM: square-summable Hilbert direct sums.}
Use \path{Mathlib/Analysis/InnerProductSpace/l2Space}, specifically
\path{lp.instInnerProductSpace} and \path{lp.inner_eq_tsum}.
For a family of Hilbert spaces and exponent two these give the Hilbert inner product of square-summable vectors as the sum of component inner products. Coordinate inclusions, projections and finite-tail norm estimates are verified directly in FD01 and MS03. The component spaces here are the closed symmetric $L^2$ subspaces already constructed in FD01, rather than an abstract Fock-space existence assumption.
\url{https://leanprover-community.github.io/mathlib4_docs/Mathlib/Analysis/InnerProductSpace/l2Space.html}.

\hypertarget{leaf:ML-PROD}{}\paragraph{ML-PROD: Tonelli, Fubini and measurable parameter integrals.}
Use \path{Mathlib/MeasureTheory/Measure/Prod} and
\path{Mathlib/MeasureTheory/Integral/Prod}. The checked declarations include
\path{MeasureTheory.lintegral_prod},
\path{MeasureTheory.integrable_prod_iff},
\path{MeasureTheory.integral_prod},
\path{MeasureTheory.integral_integral_swap},
\path{MeasureTheory.Integrable.prod_left_ae},
\path{MeasureTheory.Integrable.prod_right_ae}, and
\path{MeasureTheory.StronglyMeasurable.integral_prod_right}.
Lebesgue measure and counting measure are sigma-finite, as are their finite products and measurable restrictions. For nonnegative measurable integrands use Tonelli; for complex or Banach-valued integrands the document first supplies an integrable norm majorant and then uses Fubini. Countable sums of nonnegative sector norms are Tonelli on counting measure. The local parameter-integrability proof in OP00 avoids requiring a globally integrable majorant in the output variable when only local measurability is needed.
\url{https://leanprover-community.github.io/mathlib4_docs/Mathlib/MeasureTheory/Measure/Prod.html};
\url{https://leanprover-community.github.io/mathlib4_docs/Mathlib/MeasureTheory/Integral/Prod.html}.

\hypertarget{leaf:ML-MEAS}{}\paragraph{ML-MEAS: simple approximation and strong measurability.}
Use \path{Mathlib/MeasureTheory/Function/StronglyMeasurable/Basic}, with
\path{MeasureTheory.StronglyMeasurable.approx},
\path{MeasureTheory.StronglyMeasurable.tendsto_approx}, and
\path{stronglyMeasurable_of_tendsto}.
Strongly measurable functions admit pointwise simple approximants; limits and continuous compositions preserve the relevant measurability. FC08 proves the extra operator-family composition needed here by applying bounded continuous operators to those approximants. The symmetric-subspace version in MS02 is obtained by applying the finite symmetry projection to the approximants, so the subspace-valued field does not require a new measurable-selection theorem.
\url{https://leanprover-community.github.io/mathlib4_docs/Mathlib/MeasureTheory/Function/StronglyMeasurable/Basic.html}.

\hypertarget{leaf:ML-BOCHNER}{}\paragraph{ML-BOCHNER: norm bounds, continuous maps and dominated convergence.}
Use \path{Mathlib/MeasureTheory/Integral/Bochner/Basic},
\path{Mathlib/MeasureTheory/Integral/Bochner/ContinuousLinearMap}, and
\path{Mathlib/MeasureTheory/Integral/DominatedConvergence}. Checked declarations include
\path{MeasureTheory.integral_congr_ae},
\path{MeasureTheory.norm_integral_le_lintegral_norm},
\path{ContinuousLinearMap.integral_comp_comm}, and
\path{MeasureTheory.tendsto_integral_of_dominated_convergence}.
The latter is used with almost-everywhere convergence, measurable integrands and an explicitly integrable norm majorant. FC01 uses the original weighted squared norm as majorant; MS02 uses sums of squared errors; PT03 uses $C|v_gf||v_gh|$ and $Cg^2|fh|$; SE02 uses $Cg^2/(|k|^2(1+|k|^2))$. In each case the local lemma verifies the hypothesis, so the reference is not a request to discover a dominating function.
\url{https://leanprover-community.github.io/mathlib4_docs/Mathlib/MeasureTheory/Integral/DominatedConvergence.html};
\url{https://leanprover-community.github.io/mathlib4_docs/Mathlib/MeasureTheory/Integral/Bochner/ContinuousLinearMap.html}.

\hypertarget{leaf:ML-RIESZ}{}\paragraph{ML-RIESZ: the ordinary complex Riesz correspondence.}
Use \path{Mathlib/Analysis/InnerProductSpace/Dual}, specifically
\path{InnerProductSpace.toDual},
\path{InnerProductSpace.toDual_apply_apply}, and
\path{InnerProductSpace.toDual_symm_apply}.
For a complete complex Hilbert space, a continuous linear functional is represented by $w\mapsto\langle z,w\rangle$, and the correspondence in $z$ is conjugate-linear. BA01 explicitly conjugates a form load to obtain the linear isometric identification of conjugate-dual loads with Riesz vectors. All sesquilinear forms, adjoints and inverse-source pairings in the manuscript use that fixed orientation.
\url{https://leanprover-community.github.io/mathlib4_docs/Mathlib/Analysis/InnerProductSpace/Dual.html}.

\hypertarget{leaf:ML-HILBERT}{}\paragraph{ML-HILBERT: orthogonal-complement characterization of density.}
Use \path{Mathlib/Analysis/InnerProductSpace/Projection/Submodule}, in particular
\path{Submodule.orthogonal_orthogonal_eq_closure},
\path{Submodule.orthogonal_eq_bot_iff}, and
\path{Submodule.topologicalClosure_eq_top_iff}.
For a subspace of a complete Hilbert space, having zero orthogonal complement is equivalent to dense range by \path{Submodule.topologicalClosure_eq_top_iff}; \path{orthogonal_orthogonal_eq_closure} gives the corresponding double-complement description. In contrast, \path{Submodule.orthogonal_eq_bot_iff} requires an orthogonal-projection instance and concludes that the subspace itself is top. Do not apply that stronger interface to an arbitrary not-yet-closed range. The explanatory FA00 argument checks the zero-complement hypothesis for its self-adjoint injective $R$; the implementation reuses GW-Form. Injectivity alone does not imply dense range without the self-adjoint argument.
\url{https://leanprover-community.github.io/mathlib4_docs/Mathlib/Analysis/InnerProductSpace/Projection/Submodule.html}.

\hypertarget{leaf:ML-SERIES}{}\paragraph{ML-SERIES: geometric series in the appropriate algebra.}
Use \path{Mathlib/Analysis/SpecificLimits/Normed}. For bounded operators use \path{summable_geometric_of_norm_lt_one}, \path{Units.oneSub}, \path{geom_series_mul_neg}, \path{mul_neg_geom_series}, \path{hasSum_geom_series_inverse}, and \path{tsum_geometric_le_of_norm_lt_one} in a normed ring with \path{HasSummableGeomSeries}; completeness supplies this instance. Multiplication of endomorphisms is composition, so the two telescoping identities provide the two inverse identities without commutativity. The norm estimate retains $\|I\|$; use $\|I\|\leq1$ for the bounded-operator specialization, including a zero-dimensional space. For the scalar dyadic sums in AN01 one may instead use \path{hasSum_geometric_of_norm_lt_one} and \path{tsum_geometric_of_norm_lt_one}, whose division-ring assumptions are satisfied there, but not for the general operator algebra. BA02 retains the order/duality and inverse-identity adapters, not a new proof of geometric summability.
\url{https://raw.githubusercontent.com/leanprover-community/mathlib4/9bac2d0692944a775113f7c130fd5e5b4608a854/Mathlib/Analysis/SpecificLimits/Normed.lean}.

\hypertarget{leaf:ML-VOLUME}{}\paragraph{ML-VOLUME: interval and product-rectangle volumes.}
Use \path{Mathlib/MeasureTheory/Measure/Lebesgue/Basic}, with
\path{Real.volume_Ico}, \path{Real.volume_singleton},
\path{Real.volume_Icc_pi}, \path{Real.volume_pi_Ioo}, and
\path{Real.volume_pi_Ico}.
The product-rectangle measure is the product of side lengths; degenerate intervals have zero measure. AN00 gives a mathematical derivation of affine substitutions from these identities, but its implementation must reuse the indexed GW ball-volume facts and first check exact pinned affine measure interfaces. The generating-set proof is retained only for any affine specialization not supplied by an applicable existing theorem. It is therefore unnecessary to import spherical-coordinate integration or an exact convolution formula for the Fr\"ohlich singularity.
\url{https://leanprover-community.github.io/mathlib4_docs/Mathlib/MeasureTheory/Measure/Lebesgue/Basic.html}.

\hypertarget{leaf:ML-FTC}{}\paragraph{ML-FTC: Banach-valued fundamental theorem on a segment.}
Use \path{Mathlib/MeasureTheory/Integral/IntervalIntegral/FundThmCalculus}, specifically
\path{intervalIntegral.integral_eq_sub_of_hasDerivAt}.
For a complete real normed target, a derivative defined on the compact segment and integrable there has integral equal to the endpoint difference. All applications here have a continuous derivative on a neighborhood of that segment; continuity on a compact interval supplies integrability. Complex Hilbert spaces and operator spaces are viewed as real normed spaces for the parameter derivative. CA01 supplies the quantitative product/inverse recurrence using the ML-DIFF bounded calculus interfaces, and CA02 supplies the uniform multiplication-operator remainder bounds. Neither asserts an unbounded-family differentiability theorem.
\url{https://leanprover-community.github.io/mathlib4_docs/Mathlib/MeasureTheory/Integral/IntervalIntegral/FundThmCalculus.html}.

\hypertarget{leaf:ML-DIFF}{}\paragraph{ML-DIFF: bounded bilinear and noncommutative inverse differentiation.}
Use \path{Mathlib/Analysis/Calculus/FDeriv/Bilinear.lean}, in particular \path{IsBoundedBilinearMap.hasFDerivAt} and \path{ContinuousLinearMap.hasFDerivAt_of_bilinear}, for composition/product differentiation over $\R$. Use \path{hasFDerivAt_ringInverse} and \path{fderiv_inverse} from \path{Mathlib/Analysis/Calculus/FDeriv/Mul.lean} at an actual unit of a normed algebra with summable geometric series. In CA01 the algebra is the bounded complex-linear endomorphism space with its real scalar structure. Its units are supplied by the proved two-sided inverse; the derivative sends $A$ to $-M^{-1}AM^{-1}$. The implementation must identify \path{Ring.inverse} with the chosen inverse, discharge the real-parameter instances, and then use the chain rule. The higher-order recurrence and uniform remainder estimates in CA01--CA03 are not supplied merely by looking up the first-derivative theorem.
\url{https://raw.githubusercontent.com/leanprover-community/mathlib4/9bac2d0692944a775113f7c130fd5e5b4608a854/Mathlib/Analysis/Calculus/FDeriv/Bilinear.lean};
\url{https://raw.githubusercontent.com/leanprover-community/mathlib4/9bac2d0692944a775113f7c130fd5e5b4608a854/Mathlib/Analysis/Calculus/FDeriv/Mul.lean}.

\hypertarget{leaf:ML-IVT}{}\paragraph{ML-IVT: the one real scalar zero.}
Use \path{Mathlib/Topology/Order/IntermediateValue}, specifically
\path{intermediate_value_Icc'} (or \path{intermediate_value_Icc} with the endpoint order reversed). For a continuous real function on $[a,b]$, every value between its endpoint values is attained. SP01 checks continuity, $-1\leq |P|^2$, and the signs at these two endpoints. Strict monotonicity is proved separately in SP00; it is not supplied by IVT.
\url{https://leanprover-community.github.io/mathlib4_docs/Mathlib/Topology/Order/IntermediateValue.html}.

\hypertarget{leaf:ML-REAL}{}\paragraph{ML-REAL: elementary scalar and Hilbert algebra.}
This designation includes finite sums/products, real order and completeness, Euclidean norm inequalities, and the defining algebra of bounded linear maps. The Hilbert identities used are in \path{Mathlib/Analysis/InnerProductSpace/Basic}, including
\path{inner_conj_symm}, \path{inner_add_left}, \path{inner_add_right},
\path{inner_smul_left}, \path{inner_smul_right},
\path{inner_self_eq_norm_sq_to_K}, and the Cauchy--Schwarz inequality.
The latter may equivalently be obtained by minimizing $\|x-ay\|^2$, as is done explicitly for the positive-form version in PO01. Real quadratic Young inequalities follow from $(\sqrt\delta a-b/\sqrt\delta)^2\geq0$. No spectral, field-domain, compactness, or representation theorem is hidden in this designation.
\url{https://leanprover-community.github.io/mathlib4_docs/Mathlib/Analysis/InnerProductSpace/Basic.html}.

\subsection{Pinned-source checks and application status}\label{sec:pinnedcontracts}
For version 5, the named Riesz, density, operator-completeness, geometric-series, bilinear-derivative and ring-inverse interfaces above were inspected in source at the actual pin \path{9bac2d0692944a775113f7c130fd5e5b4608a854}. The relevant density and Riesz sources are respectively
\path{Mathlib/Analysis/InnerProductSpace/Projection/Submodule.lean} and \path{Mathlib/Analysis/InnerProductSpace/Dual.lean}; their immutable links and the other inspected modules are recorded in \path{metadata/library-contracts.json}. The rest of the inherited reference list keeps its previous evidence level; these selected checks are not a blanket re-certification of every leaf. Every new fiber specialization remains \contract{APPLICATION PENDING}. Section~\ref{sec:execution} prescribes when to compile those applications. The GW index records names, source ranges and complete-file hashes; the unchanged source archive, not a truncated header, supplies the implicit contexts and proof bodies.

\subsection{Translation boundary and assumptions audit}
The source-level GW declarations and pinned mathlib commit are now mapped in Section~\ref{sec:gwverify}; actual compilation and declaration-type application on that pin remain outstanding. This is a compatibility check, not permission to replace any local result by an axiom. The final theorem must contain no hypotheses asserting the kernel identity, saturation identity, existence of the Hamiltonian or a spectral gap: those are conclusions of earlier nodes. The canonical source-reuse and adapter decisions are fixed in the registry. RE01 is an optional mathematical transport result, not a required branch of the canonical implementation.

The implementation must preserve four distinctions: physical vectors versus form-dual loads; energy coercivity versus mere positivity in the physical norm; almost-everywhere fields versus pointwise operators; and full resolvent exclusion versus a zero-kernel statement. The source and destination spaces in the proofs fix these choices. Norm limits and all integral identities are established before the corresponding class or operator is formed. A completed Lean implementation should be checked for local proof placeholders and extra analytic axioms, and should build with no external GW package. Classical choice and the usual logical axioms of the mathlib foundation are not additional physical assumptions.


\section{Source concordance and omitted branches}\label{sec:concordance}
The source of the nonzero-momentum method is Mitrouskas~\cite{mitrouskas}; the exact infrared cancellation comes from Seiringer~\cite{seiringer}. The paper-level foundation for the available formalization is Griesemer--W\"unsch~\cite{gw}. The proof above is a specialized and reorganized argument, not a claim that every displayed lemma occurs verbatim in one of those papers.

\begin{longtable}{>{\raggedright\arraybackslash}p{.25\textwidth}>{\raggedright\arraybackslash}p{.36\textwidth}>{\raggedright\arraybackslash}p{.31\textwidth}}
\toprule
Blueprint portion & Source argument & Deliberate modification\\
\midrule\endhead
MS01--MS03, FD01--FD02, FC01--FC09, UV01--UV02, PT01--PT02 & Fock sector definitions; the pull-through steps used in Seiringer and Mitrouskas & Spell out weak domains, Bochner fields, the single residual null set and fixed-energy cutoff limits. These expanded lemmas are not asserted to occur verbatim in the sources.\\
FB01--FB03 & Mitrouskas, Section 2 and Appendix A & Uniform form estimates replace every appeal to the global energy minimum.\\
SP01 and SP03 & Vacuum/one-phonon reductions in Mitrouskas, Section 2 & Full energy-space inverse construction replaces localization based on essential spectrum.\\
RS01--RS03 & Mitrouskas, Lemma 4.1 and its proof & Retain weighted bounds and derivatives; omit perturbative remainder estimates used only in Section 4.2.\\
PT03, KE01--KE02 & Mitrouskas, Lemma 3.1 and Lemma 4.3 & Derive the exact kernel by scalar smeared identities, with no pointwise creation or delta distributions. Reflection suffices for the needed cancellations.\\
AN03, WE01--WE06 & Mitrouskas, Lemma 4.4 and Section 4.4 & Give direct three-dimensional estimates; use ordinary $L^\infty$ and symmetric differences instead of differentiation under the integral.\\
SE01--SE03 & Seiringer, Section 2, especially Lemmas 2.1 and 2.3; Mitrouskas, Lemma 3.2 & Derive the key identity from the ground sector equations; prove $N_0=c_0$ without normalizing by $c_0$.\\
PO01--PO03 & Mitrouskas, scalar criterion following Proposition 3.3 & Completion of a square proves the sufficient criterion without $c_P>0$ or Sherman--Morrison.\\
\bottomrule
\end{longtable}

The following source branches have no downstream use in this proof and should not be translated for this milestone: the extension to all momenta by Dahlb\ae k; the essential-spectrum formula; positivity improvement and uniqueness by that route; the global minimum theorem for the dispersion; the estimate $E_g(P)\leq E_g(0)+\abs P^2$; the positive lower bound and perturbative expansion for $c_P$ in Mitrouskas, Section 4.2; arbitrary dimensions and arbitrary form factors in Seiringer. A few algebraic constructions incidentally imply one-dimensionality of a low eigenspace, but no separate simplicity theorem is used.

\appendix
\section{Implementation-sensitive proofs and acceptance contracts}\label{app:domain}
This appendix records the small analytical steps most likely to be lost when translating the proof. It is part of the mathematical specification. The repository-specific identification of GW declarations is now given in Section~\ref{sec:gwverify}. Translation of each new local mathematical construction must reuse these audited declarations when applicable.

\subsection{Object and domain inventory}
\begin{longtable}{>{\raggedright\arraybackslash}p{.22\textwidth}>{\raggedright\arraybackslash}p{.34\textwidth}>{\raggedright\arraybackslash}p{.36\textwidth}}
\toprule
Symbol & Mathematical type & Operations permitted here\\\midrule\endhead
$v_g(k)$ & Measurable scalar function; locally $L^2$ & Weighted source integrals; no unweighted global $L^2$ claim.\\
$v_\Lambda$ & One-particle $L^2$ vector with bounded support & Ordinary $a(v_\Lambda)$ and $a^*(v_\Lambda)$.\\
$\mathfrak v_g$ & Continuous conjugate-linear functional on $\Q_1$ & Weak inverses and Riesz representatives.\\
$\mathcal A u$ & Bochner $L^2$ field for $u\in V_N$ & Equality as an $L^2$ class or a.e.; local restriction; scalar pairing.\\
$a^*(f)$ & Number-form-domain to physical Fock space for $f\in L^2$ & The explicit sector formula and its smeared adjoint. No $a_k^*$ is defined.\\
$X_P,Y_p,Z_p$ & Physical bounded inverses on their own sectors; also energy-dual-to-energy inverse maps & State which version acts on each argument. Zero extension to other sectors is not an inverse there.\\
$h_p$ & Energy-domain vector solving a singular-source equation & Number-field map, Hilbert pairings, graph-norm limits, parameter derivatives.\\
$D_P f$ & Element of $\Q_1^{\times}$ & Evaluation on graph-domain tests; not assumed to be $L^2$.\\
$\ell_{g,P}$ & Continuous scalar functional on $\Q_1$ & Vacuum and threshold rank-one \emph{forms}; finite at $P=0$.\\
$t_P$ & $L^2$ vector only for $P\ne0$ (or the zero function for $g=0$) & Bounded rank-one forms only in that branch.\\
$t_0=g/\abs q^2$ & Locally $L^2$ away from zero, not globally $L^2$ if $g\ne0$ & Annulus identities and the infrared contradiction, never $B^{-1}t_0$.\\
$\eta_P,\zeta_P$ & Measurable representatives of real $L^\infty$-based constructions & a.e. inequalities and integrals against absolutely continuous densities.\\
$H_g(P)$ & The actual unbounded associated operator & Its domain characterization and bounded two-sided resolvent inverse, not bounded-algebra spectrum applied to $H_g(P)$.\\
\bottomrule
\end{longtable}

\subsection{Polarization and off-diagonal bounds}
Suppose a Hermitian form $a$ on a complex Hilbert space satisfies
$\abs{a[u]}\leq c\norm u^2$ on a dense linear subspace. The polarization identity bounds its mixed values by $2c\norm u\norm w$, which is sufficient everywhere here. To see the product bound without losing a factor depending on the vectors, first apply polarization to unit vectors, where each of $u+i^j w$ has squared norm summing to $8$, and then use sesquilinearity to rescale. The form therefore extends uniquely to the completion. Its Hermitian symmetry is preserved by continuity. A sharper constant can be obtained from its self-adjoint representative but is unnecessary.

Apply this argument to each cutoff difference in FB02 before taking the limit. The sector-selection rule also passes to the limit: a form entry which vanishes for every cutoff has zero limit. A diagonal particle-number entry is zero, the vacuum couples only to the one-particle sector, and the one-to-high coupling has only a two-particle test component even though the high inverse mixes particle numbers.

The high-momentum commutator in FB01 must first use
\[
 w_{j,L,\Lambda}(k)=k_j\abs k^{-2}v_g(k)
             \mathbf1_{\{L<\abs k\leq\Lambda\}}.
\]
At finite $\Lambda$, both $w_{j,L,\Lambda}$ and $k_iw_{j,L,\Lambda}$ are square-integrable, so the core commutator is an equality of genuine smeared operators. The estimate is uniform in $\Lambda$. Only after obtaining the mixed graph bound is the limit taken. Writing the same formal commutator with an uncut $w_j$ and treating $a(k_jw_j)$ as an ordinary smeared operator would not be justified.

\subsection{Positive inverses, monotonicity, and graph coercivity}
For a bounded positive invertible $T$ on an energy Hilbert space, put $v=T^{-1}z$. Then
\[
 \ip z{T^{-1}z}=\ip{Tv}v=\ip v{Tv}\geq0.
\]
Self-adjointness of $T^{-1}$ follows by taking adjoints of $TT^{-1}=T^{-1}T=I$. For any $u$ the exact square-completion identity is
\[
 2\Re\ip z u-\ip u{Tu}
 =\ip z{T^{-1}z}-\ip{u-T^{-1}z}{T(u-T^{-1}z)}.
\]
Its supremum over $u$ is attained at $T^{-1}z$. This proves the inverse-order comparison in SP03 directly and for arbitrary form-dual loads after Riesz conversion. It does not rely on a spectral measure.

For the interpolation in FA03, positivity in the physical norm and a lower kinetic bound are separate inputs. The convex combination gives an explicit positive coefficient on both parts of the graph norm. For the shifted one-sector weight $1+\abs{k-P}^2$, comparison with $1+\abs k^2$ is uniform on $\abs P\leq1/4$. Consequently the Schur representative is onto the \emph{energy space}, which is stronger than merely having no physical kernel. FA01 then checks both inverse equations and the actual operator domain of the inverse image.

\subsection{Weighted multiplication derivatives: a finite proof recipe}
Here is a fully finite way to justify the norm differentiability in RS02. On a sector of order $m\geq n$, let
\[
 x=K_m+p,\qquad b=m+1-n-e\geq1,\qquad w=\abs x^2+b.
\]
The multiplier $G=w^{-1/2}$ and its first three derivatives are given in RS02. Its fourth derivative is
\begin{align*}
 \partial_i\partial_j\partial_k\partial_lG
 ={}&3(\delta_{ij}\delta_{kl}+\delta_{ik}\delta_{jl}+\delta_{il}\delta_{jk})w^{-5/2}\\
 &-15(\delta_{ij}x_kx_l+\delta_{ik}x_jx_l+\delta_{il}x_jx_k
       +\delta_{jk}x_ix_l+\delta_{jl}x_ix_k+\delta_{kl}x_ix_j)w^{-7/2}\\
 &+105x_ix_jx_kx_lw^{-9/2}.
\end{align*}
Since $\abs x\leq\sqrt w$ and $w\geq1$, division by $G$ gives coordinate-derivative bounds $1,4,24,204$ for orders $1,2,3,4$, respectively. Multilinear operator norms differ by at most a fixed dimension-dependent factor, obtained by expanding along the three coordinates and using $\abs{h_i}\leq\norm h$.

For $\abs{p-p_0}\leq1$, $w_{p_0}\leq4w_p$, so multiplication by $W_{n,p_0}$ increases each of these bounds by at most $2$. Apply the scalar fundamental theorem along $p+th$ to the derivatives through order three; the next derivative gives a uniform Taylor remainder in the weighted essential-supremum norm. The operator norm of a multiplication map is at most the essential supremum of its multiplier. Therefore these scalar remainders prove the asserted operator-norm derivatives into the fixed energy space. This is not a mere pointwise-differentiability assertion.

For $z_p=G_{1,p}v_g$, the same remainder estimate is multiplied by $G_{1,p_0}v_g\in L^2$, whose norm is $O(\abs g)$ by AN02. Thus $p\mapsto z_p$ has the stated Hilbert-norm derivatives without differentiating an improper integral. The form sandwich $T_p$ is a composition of bounded maps on the fixed energy space, so ordinary product rules apply. For a multi-index $\alpha$, inverse differentiation is the finite recursion
\[
 D^\alpha Q_p=-Q_p\sum_{0<\beta\leq\alpha}
       \binom{\alpha}{\beta}(D^\beta M_p)(D^{\alpha-\beta}Q_p),
 \qquad Q_p=M_p^{-1}.
\]
The zeroth inverse norm is at most $2$; every other factor has its previously fixed bound. Induction on $\abs\alpha\leq3$ produces constants independent of $p,g$. When extracting outer $G$ factors from derivatives of $GQG$, move a multiplier ratio through an adjacent $G$, \emph{not} through the noncommuting $Q$. The displayed bounded expression for $\Ecal(p)$ then makes its derivatives finite products with two $O(\abs g)$ source factors, proving the Hessian estimate.

\subsection{Reality and reflection without representation theory}
Let $J$ be componentwise complex conjugation and $\mathcal R$ the simultaneous change $k_i\mapsto-k_i$. The explicit sector formulas, followed by the mixed-form limit, give
\[
 h_{g,P}(Ju,Jw)=\overline{h_{g,P}(u,w)},\qquad
 h_{g,-P}(\mathcal R u,\mathcal R w)=h_{g,P}(u,w).
\]
Both maps preserve the relevant form domains. To prove a corresponding identity for an inverse, apply the map to its weak equation and use uniqueness; no separately developed antiunitary or group-representation library is needed. In particular
\[
 Jh_p=h_p,\quad J\Omega=\Omega,\quad
 \mathcal R h_p=h_{-p},\quad
 \mathcal R Z_p\mathcal R=Z_{-p},\quad
 \mathcal R X_P\mathcal R=X_{-P}.
\]
The scalar source pairing is invariant under $\mathcal R$ because $v_g$ is real and even. These equalities prove the reality, symmetry and reflection properties used in RS03 and KE01. They do not assert any ordering between physical fiber energies.

The multiplication function
\[
 m(q)=1+\frac{\Ecal(q)-e}{\abs q^2},\quad q\ne0,
\]
can be assigned any convenient bounded positive value, for example $m(0)=1$, at the origin. Reflection alone does not imply that this quotient has a direction-independent limit at $0$. No such extension is required: $m$ is an even measurable bounded multiplier, and its inverse is bounded. The WE03 argument explicitly avoids differentiating $m$.

\subsection{Kernel maps on quotient spaces}
If a measurable kernel $K(q,r)$ satisfies $\abs{K(q,r)}\leq C a(q)a(r)$ with $a\in L^2$, then for $f\in L^2$,
\[
 \int\abs{K(q,r)f(r)}\dd r
 \leq C a(q)\norm a_2\norm f_2
\]
for a.e. $q$. This proves both existence and the $L^2$ output bound. Joint measurability plus the parameter-integral measurability theorem gives a measurable output. Replacing $f$ on a null set leaves the integrals unchanged for each $q$ for which the domination holds. Two test vectors give the integrable product majorant
$C\abs{f(q)}a(q)a(r)\abs{h(r)}$ on the product measure. Fubini therefore proves the adjoint identity from the symmetric kernel. These are the exact steps needed for OP01; no compact-operator or Hilbert--Schmidt class is required.

For WE01, write
\[
 H_P(q,r)=\frac{F_P(q,r)}{\omega(q)}\mu_P(r)\omega(r)
 \quad(q\ne0),\qquad H_P(0,r)=0.
\]
It is measurable, and $\abs{H_P(q,r)}\leq C\omega(r)^2\mu_P(r)$, an $L^1$ majorant independent of $q$. Each $L^\infty$ class has an a.e. bounded representative; modifying it to zero on the exceptional set makes a bounded representative. The integral formula defines a bounded measurable output. The input null-set change cannot affect it, because the density $\omega^2\mu_P$ is Lebesgue-integrable and hence gives zero mass to null sets. The output's essential supremum is at most the $L^1$ majorant norm times the input essential supremum. Linearity, quotient independence and the operator norm bound now follow in that order.

Reflection on $L^\infty$ is a well-defined isometry since $q\mapsto-q$ preserves Lebesgue measure. At $P=0$ it commutes with the multiplier and integral operators, by change of variables and $F_0(-q,-r)=F_0(q,r)$. Uniqueness of the Neumann-series solution thus gives a.e. evenness. All uses of that evenness occur under integrals against the absolutely continuous densities just described.

\subsection{No differentiation under the weighted integral is required}
The strengthened WE03 proof uses a symmetric second difference. Set
\[
 J(q)=\int F_0(q,r)\mu_0(r)\eta_0(r)\dd r.
\]
Then $J$ is even and $J(0)=0$. The Hessian bound on $F_0$ yields
\[
 \abs{F_0(q,r)+F_0(-q,r)-2F_0(0,r)}
       \leq C\abs q^2\omega(r).
\]
Integrate this bound against $\mu_0\abs{\eta_0}$ and use the original
$\abs{\eta_0}\leq C\omega$. It follows that
$2\abs{J(q)}\leq C\abs q^2\int\omega^2\mu_0$.
For $\abs q\geq1$, the ordinary kernel bound gives $\abs J\leq C$. Thus $J=O(\omega^2)$, exactly what the weighted equation needs. The derivative of an $L^\infty$ representative, differentiation under a singular integral, and differentiability of $\eta_0$ are all unnecessary.

\subsection{Explicit annulus extraction in the saturation argument}
Let $A_j=\{k:1/j<\abs k<j\}$, $j\geq2$. For any $z\in L^2(A_j)$, extend it by zero and set
$f_z(k)=\mathbf1_{A_j}(k)z(k)/\abs k$. Then
\[
 \norm{f_z}_2\leq j\norm z_2,\qquad
 \norm{\abs k f_z}_2=\norm z_2,
\]
so $f_z\in\Q_1$ and is a permissible test in the weak identity SE02. Substitute this test into the polarized OP03 identity and the SE02 equation. All resulting divided terms are $L^2$ on $A_j$: the right source $g/\abs k^2$ is bounded there, and the left consists of restrictions of the global $L^2$ vectors $b_0$, $u_0$, and $B_0u_0$. Nondegeneracy of the $L^2(A_j)$ inner product, tested against \emph{all} $z$, gives the divided equality a.e. on $A_j$. A countable union of exceptional sets and exhaustion of $\R^3\setminus\{0\}$ gives the global a.e. identity. Singletons have zero Lebesgue measure.

If a nonzero scalar multiplied by $g/\abs k^2$ equalled a global $L^2$ vector, scalar division would put that function in $L^2$, contradicting AN01. The argument does not cancel a formal infinite vector; it compares a legitimate local equality with a proved global nonintegrability statement.

\subsection{Dyadic estimates without exact geometric constants}
For the lower annulus volume used in AN01, one may avoid any open/closed sphere convention by using
$B(0,r)\setminus B(0,r/2)$ and the identity
$\operatorname{vol}(B(0,r))=r^3\operatorname{vol}(B(0,1))$.
The smaller ball is a subset of the larger, so the annulus measure is their difference regardless of how its outer boundary is excluded. It equals $(7/8)v_3r^3$. On it $\abs k^{-4}\geq r^{-4}$, so the integral is at least $(7/8)v_3/r$. Arbitrarily small radii already contradict finiteness, without summing disjoint shells or computing an exact polar integral. AN00 explains finite positive unit-ball volume by cube inclusions and the affine formulas by rectangle volumes; the implementation reuses GW-Model for the volume clauses and applicable pinned measure interfaces for substitutions, as its implementation contract requires. AN02--AN03 use precisely those formulas, not a rotational-coordinate integration theorem.

At $P\ne0$, the two singularities $0$ and $-P$ are distinct. After scaling by $\rho=\abs P$, split into balls of radius $1/4$ around $0$ and $-P/\rho$, their remaining bounded complement, and the exterior of the radius-$2$ ball. This partition gives constants independent of the direction of $P$. No choice of a rotation sending $P$ to a coordinate vector is required.

\subsection{A finite checklist for every translated node}
For each node the local source file should establish the statement with all parameters and domains shown here, identify every load as physical or form-dual, and discharge measurability/integrability before invoking an integral theorem. Constants must be quantified before the final coupling is chosen. Quotient-space equalities are a.e. unless a continuous scalar kernel or explicitly constructed representative has been selected.

The source integration uses the canonical GW realization and the explicitly proved adapters; it must preserve their physical embeddings and existing partial-operator domains. The Hamiltonian construction, weak-domain and inverse equivalences, and the cases $g=0$, $P=0$, $\abs P=1/4$ and $\lambda=L_s$ are proved in FA00--FA05 and MN01--MN03, rather than left as additional mathematical assumptions. After translation, the corresponding declarations and the assumptions of the final Lean theorem must be checked by Lean. A successful compilation of this TeX file certifies none of those Lean facts. GW01--GW07 and EN01--EN02 now fix the source-level mathematical interface: any translation mismatch must be reported against its stated domain or convention, not patched by silently reimplementing the foundation.

\subsection{What remains unverified in this revision}
The analytical arguments above are expanded mathematical proofs, not machine-checked proofs. The uploaded GW source has now been examined, and exact reuse declarations, import closure, and convention corrections are documented in Section~\ref{sec:gwverify}. This is a static source audit, not a compilation. A mathematical review should in particular check cutoff commutation PT01, integrated kernel PT03, and weak-to-annulus saturation SE02--SE03; these domain-sensitive links are new and do not follow from GW. Selected pinned mathlib source interfaces were rechecked for version 5 as recorded in Section~\ref{sec:pinnedcontracts}; neither a local mathlib build nor the new application checks were performed here. The remaining index retains its inherited evidence labels. This is an evidence boundary, not an additional mathematical hypothesis.

\section{Mathematical outputs to test independently}
Before connecting the entire graph, the following checks help expose a convention or domain mismatch.

\paragraph{Free reference check.}
At $g=0$, $h_p=0$, $e=0$, and
\[
 C_P(q,r)=\frac1{1-\abs P^2}-\frac1{1+\abs{P+q+r}^2}.
\]
Although this kernel is nonzero, its contribution to the threshold form is multiplied by $g^2$. Hence $B_P=I$, $b_P=0$, and $\widehat O_P[f]=\norm{\abs q f}^2$. This tests the signs and the $+1$ shift in $Z_p$.

\paragraph{Ground scalar check.}
The vacuum equation for $\Psi_0=\Omega-h_0$ is $e=-\ell_g(h_0)$, so $\ell_g(f_0)=-e$. Reversing this sign breaks \eqref{eq:infraredcancel}. The zero-momentum one-particle denominator before the high elimination is $\abs k^2+1-e$, whereas the threshold denominator is $\abs k^2-e$; the difference of one is essential in SE02.

\paragraph{Singular scaling check.}
For $P\ne0$, $\norm{t_P}^2$ is proportional to $g^2/\abs P$, not $g^2/\abs P^2$. The compensating comparison estimate is $\abs{N_P-N_0}\leq Cg^2\abs P$. No assertion of norm convergence of $t_P$ at the origin should appear in the proof graph.

\paragraph{Spectral endpoint check.}
The final resolvent proof uses $\varepsilon>0$. It makes no claim at $\varepsilon=0$, and therefore does not identify the threshold as essential spectrum or even assert here that it belongs to the spectrum. The target theorem does not need either statement.


\clearpage
\section{Verified GW source concordance and mandatory reuse}\label{sec:gwverify}
\noindent\textbf{Audited uploaded source:} \path{inputs/GW-Froehlich-GitHub-release.zip}; SHA-256 \path{e180dfcab487864a8c8c03182daf2ba06ba40c1f48bd2a12ea6240a54e62e132}. GitHub project \url{https://github.com/PolaronFormalization/GW-Froehlich-formalization}. The source archive identifies Lean \path{leanprover/lean4:v4.36.0-rc2} and mathlib commit \path{9bac2d0692944a775113f7c130fd5e5b4608a854}. The public Git commit matching the uploaded release is not established. There are 115 local Lean modules under \path{GW/}; \path{GW.lean} imports their library root. The source imports only GW and mathlib modules. No Lean or Lake executable was available in this audit, so the release is \emph{not certified to build here}. The archive README likewise warns the source-only release was not independently rebuilt. No instance of \path{sorry}, \path{admit}, or a source-level \path{axiom} was found by text scan, but this is \emph{not} a compiled axiom audit.

\subsection{Non-negotiable type and normalization contract}
\begin{enumerate}
\item Use $\texttt{GW.Space}=\texttt{EuclideanSpace}\ \R\ (\texttt{Fin}\ 3)$, \path{GW.OnePhonon}, \path{GW.Fock.Sector} and \path{GW.Fock.Carrier = lp Sector 2} as the \emph{canonical} position/momentum, one-boson, symmetric-sector and bosonic-Fock types. \path{GW.Fock.Physical = Lp Carrier 2 volume} describes the \emph{full electron+field} Hilbert space, not a fixed-momentum fiber. Build the new fiber operator directly on \path{GW.Fock.Carrier}.
\item The prescribed vacuum is \path{GW.Fock.vacuum}. The sector zero is definitionally $\C$. GW has the exact ordinary sector creation and annihilation formulas, their $\sqrt{n+1}$ factors, adjoints, and CCR. Use \path{GW.Fock.ordinaryField} without inserting an extra $\sqrt{2}$ or Fourier factor.
\item To use signed $g$, put $\alpha=g^2$ and $\varepsilon_g=1$ for $g\ge0$, $-1$ for $g<0$. Then, including the origin under the defined zero-divisor convention, $g/\abs k=\varepsilon_g\,\texttt{GW.Model.v}(g^2,k)$ because $\sqrt{g^2}=\abs g$. Apply the identical sign adapter to finite UV cutoffs. The GW interaction convention is \emph{not} obtained by setting $\alpha=g$.
\item GW's \path{GW.Fock.fiberPointwiseAnnihilator} is indexed by \emph{electron position} $x$, and \path{GW.Fock.ordinary_annihilation_number_pullthrough} is the \emph{smeared number-operator} commutation formula. Neither is the momentum-indexed $a_k$ required by Seiringer's resolvent pull-through. The new measurable $L^2(k;\F)$ field must be constructed from the already-proved product-isometry and sector formulas, preserving its almost-everywhere nature.
\item \path{GW.Fock.freeMomentumSquare} is the full-space \emph{electron} momentum operator and is not phonon field momentum $P_f$. Define $P_f$ as the sector multiplier $\sum_j k_j$ on \path{GW.Fock.Carrier}; use the existing GW number operator and form representation for $H_g(P)$.
\end{enumerate}

\subsection{Required no-reimplementation protocol}
Copy \path{GW/} and \path{GW.lean} into the new Lean repository, preserving module names and source attribution. Retain the audited Lean and mathlib revisions initially and verify \path{lake build GW} before new work; mathlib remains the sole external Lake dependency. Use \path{import GW} only for a retained-source baseline or final coverage check. New node and scratch modules must use the specific compiled GW modules and accepted node interfaces they need. Retaining all GW sources does not require loading the umbrella in every attempt. For every node whose prerequisites include a GW source group, use the exact declarations indexed under that group; \emph{do not} translate duplicate elementary proof text from the earlier sections. This rule also applies to the existing general mathlib interfaces in BA02 and CA01. FA02 is a direct reuse of GW-Coercive, and AN00 reuses the GW volume clauses. Only the required mathematical specialization is new. The required adapters are GW01--GW07, EN01--EN02, LD01 and TR01, whose definitions and proofs are supplied here. Do not design an alternative adapter or a second carrier during translation; report any source-statement mismatch against the specified contract. All actual new work must go into separate fiber/threshold modules. The full-space Gross-transform results and \path{GW.Main.construction} do not establish the fiber construction. Check legal attribution and permission before redistribution: the uploaded README does not grant an open-source license.

\subsection{Exact declaration index and immutable source evidence}
The table records 178 distinct declarations directly indexed as source leaves: the original 173 plus five coercive-inverse declarations. Lines are one-based in the unchanged supplied release. The companion \path{metadata/gw-source-index.json} records source locations, complete-file hashes, source groups, and consuming nodes. It does not replace generic contexts by guessed arguments; read the full indexed module in the unchanged archive for section variables, instances, statements and proof bodies.
\small\begin{longtable}{p{.58\textwidth}p{.35\textwidth}}\toprule Declaration & Source location \\\midrule\endhead
\multicolumn{2}{l}{\textbf{GW-Sectors}}\\*
\path{GW.Space} & \path{GW/Conventions.lean}: 15\\
\path{GW.OnePhonon} & \path{GW/Conventions.lean}: 22\\
\path{GW.IsSelfAdjointOperator} & \path{GW/Conventions.lean}: 25\\
\path{GW.Fock.Coordinates} & \path{GW/Fock/Sectors.lean}: 14\\
\path{GW.Fock.sectorMeasure} & \path{GW/Fock/Sectors.lean}: 16\\
\path{GW.Fock.AmbientSector} & \path{GW/Fock/Sectors.lean}: 19\\
\path{GW.Fock.symmetricSector} & \path{GW/Fock/Sectors.lean}: 67\\
\path{GW.Fock.Sector} & \path{GW/Fock/Sectors.lean}: 88\\
\path{GW.Fock.Carrier} & \path{GW/Fock/Sectors.lean}: 111\\
\path{GW.Fock.vacuum} & \path{GW/Fock/Sectors.lean}: 117\\
\path{GW.Fock.vacuum_zero} & \path{GW/Fock/Sectors.lean}: 119\\
\path{GW.Fock.vacuum_succ} & \path{GW/Fock/Sectors.lean}: 121\\
\path{GW.Fock.vacuum_norm} & \path{GW/Fock/Sectors.lean}: 122\\
\path{GW.Fock.carrier_inner} & \path{GW/Fock/Sectors.lean}: 128\\
\path{GW.Fock.carrier_norm_sq} & \path{GW/Fock/Sectors.lean}: 130\\
\path{GW.Fock.cutoff} & \path{GW/Fock/Sectors.lean}: 138\\
\path{GW.Fock.cutoff_apply} & \path{GW/Fock/Sectors.lean}: 147\\
\path{GW.Fock.cutoff_norm_le} & \path{GW/Fock/Sectors.lean}: 152\\
\path{GW.Fock.cutoff_tendsto} & \path{GW/Fock/Sectors.lean}: 168\\
\path{GW.Fock.cutoff_pairing} & \path{GW/Fock/Sectors.lean}: 179\\
\path{GW.Fock.P} & \path{GW/Fock/Sectors.lean}: 173\\
\multicolumn{2}{l}{\textbf{GW-Coordinates}}\\*
\path{GW.Fock.sectorMeasure_zero} & \path{GW/Fock/Coordinates.lean}: 10\\
\path{GW.Fock.vacuumAdapter} & \path{GW/Fock/Coordinates.lean}: 35\\
\path{GW.Fock.sectorInclusion} & \path{GW/Fock/Coordinates.lean}: 50\\
\path{GW.Fock.splitCoordinates} & \path{GW/Fock/Coordinates.lean}: 102\\
\path{GW.Fock.splitCoordinates_preserving} & \path{GW/Fock/Coordinates.lean}: 108\\
\path{GW.Fock.splitCoordinates_inverse_preserving} & \path{GW/Fock/Coordinates.lean}: 113\\
\path{GW.Fock.splitUnitary} & \path{GW/Fock/Coordinates.lean}: 140\\
\path{GW.Fock.splitUnitary_ae} & \path{GW/Fock/Coordinates.lean}: 144\\
\multicolumn{2}{l}{\textbf{GW-Product}}\\*
\path{GW.Fock.amplify} & \path{GW/Fock/FockAmplification.lean}: 17\\
\path{GW.Fock.amplify_ae} & \path{GW/Fock/FockAmplification.lean}: 18\\
\path{GW.Fock.amplify_norm_le} & \path{GW/Fock/FockAmplification.lean}: 20\\
\path{GW.Fock.contractionValue} & \path{GW/Fock/FockAmplification.lean}: 89\\
\path{GW.Fock.contraction_integrable} & \path{GW/Fock/FockAmplification.lean}: 81\\
\path{GW.Fock.contractionValue_bound} & \path{GW/Fock/FockAmplification.lean}: 93\\
\path{GW.Fock.contraction} & \path{GW/Fock/FockAmplification.lean}: 126\\
\path{GW.Fock.contraction_pairing} & \path{GW/Fock/FockAmplification.lean}: 161\\
\path{GW.Fock.tensorEmbedding} & \path{GW/Fock/FockAmplification.lean}: 142\\
\path{GW.Fock.tensorEmbedding_ae} & \path{GW/Fock/FockAmplification.lean}: 151\\
\path{GW.Fock.tensorEmbedding_adjoint} & \path{GW/Fock/FockAmplification.lean}: 169\\
\path{GW.Fock.rectangle_dense_criterion} & \path{GW/Fock/FockAmplification.lean}: 514\\
\path{GW.Fock.productUnitary} & \path{GW/Fock/FockAmplification.lean}: 550\\
\path{GW.Fock.productUnitary_rectangle} & \path{GW/Fock/FockAmplification.lean}: 553\\
\path{GW.Fock.productUnitary_natural} & \path{GW/Fock/FockAmplification.lean}: 575\\
\path{GW.Fock.productUnitary_tensor} & \path{GW/Fock/FockAmplification.lean}: 671\\
\path{GW.Fock.product_map_ext} & \path{GW/Fock/FockAmplification.lean}: 559\\
\path{GW.Fock.product_slice_memLp} & \path{GW/Fock/FockAmplification.lean}: 696\\
\multicolumn{2}{l}{\textbf{GW-Generalized}}\\*
\path{GW.Fock.kernelPermutation} & \path{GW/Fock/FockGeneralized.lean}: 143\\
\path{GW.Fock.kernelSymmetrizer_of_fixed} & \path{GW/Fock/FockGeneralized.lean}: 275\\
\path{GW.Fock.kernelFixedInclusion} & \path{GW/Fock/FockGeneralized.lean}: 323\\
\path{GW.Fock.kernelFixedProjection} & \path{GW/Fock/FockGeneralized.lean}: 327\\
\path{GW.Fock.ordinarySectorKernel} & \path{GW/Fock/FockGeneralized.lean}: 458\\
\path{GW.Fock.ordinaryKernel} & \path{GW/Fock/FockGeneralized.lean}: 713\\
\path{GW.Fock.ordinaryCreation} & \path{GW/Fock/FockGeneralized.lean}: 715\\
\path{GW.Fock.ordinaryAnnihilation} & \path{GW/Fock/FockGeneralized.lean}: 718\\
\path{GW.Fock.sectorCreation} & \path{GW/Fock/FockGeneralized.lean}: 637\\
\path{GW.Fock.sectorAnnihilation} & \path{GW/Fock/FockGeneralized.lean}: 641\\
\path{GW.Fock.sectorCreation_pairing} & \path{GW/Fock/FockGeneralized.lean}: 663\\
\path{GW.Fock.kernelOneBoson} & \path{GW/Fock/FockGeneralized.lean}: 826\\
\path{GW.Fock.kernelOneBoson_coe} & \path{GW/Fock/FockGeneralized.lean}: 829\\
\path{GW.Fock.oneBosonCoordinate_preserving} & \path{GW/Fock/FockGeneralized.lean}: 770\\
\path{GW.Fock.kernelFixedVacuum_coe} & \path{GW/Fock/FockGeneralized.lean}: 832\\
\path{GW.Fock.kernelCreation_vacuum} & \path{GW/Fock/FockGeneralized.lean}: 834\\
\path{GW.Fock.kernelAnnihilation_vacuum} & \path{GW/Fock/FockGeneralized.lean}: 842\\
\path{GW.Fock.ordinaryCreationSector_ae} & \path{GW/Fock/FockGeneralized.lean}: 886\\
\path{GW.Fock.ordinaryCreation_domain} & \path{GW/Fock/FockGeneralized.lean}: 926\\
\path{GW.Fock.ordinaryAnnihilation_domain} & \path{GW/Fock/FockGeneralized.lean}: 931\\
\path{GW.Fock.ordinaryCreation_zero} & \path{GW/Fock/FockGeneralized.lean}: 936\\
\path{GW.Fock.ordinaryCreation_succ} & \path{GW/Fock/FockGeneralized.lean}: 939\\
\path{GW.Fock.ordinaryAnnihilation_apply} & \path{GW/Fock/FockGeneralized.lean}: 943\\
\multicolumn{2}{l}{\textbf{GW-Number}}\\*
\path{GW.Fock.numberPower} & \path{GW/Fock/Number.lean}: 76\\
\path{GW.Fock.numberPower_domain} & \path{GW/Fock/Number.lean}: 141\\
\path{GW.Fock.numberPower_domain_eq} & \path{GW/Fock/Number.lean}: 147\\
\path{GW.Fock.numberPower_apply} & \path{GW/Fock/Number.lean}: 171\\
\path{GW.Fock.numberPower_norm_sq} & \path{GW/Fock/Number.lean}: 187\\
\path{GW.Fock.shiftedNumberPower} & \path{GW/Fock/Number.lean}: 70\\
\path{GW.Fock.shiftedNumberPower_apply} & \path{GW/Fock/Number.lean}: 177\\
\path{GW.Fock.shiftedNumberInput} & \path{GW/Fock/Number.lean}: 182\\
\path{GW.Fock.numberOperator} & \path{GW/Fock/Number.lean}: 245\\
\multicolumn{2}{l}{\textbf{GW-Bochner}}\\*
\path{GW.Fock.BochnerSectorSum} & \path{GW/Fock/Bochner.lean}: 14\\
\path{GW.Fock.BochnerFock} & \path{GW/Fock/Bochner.lean}: 15\\
\path{GW.Fock.bochnerReconstruct} & \path{GW/Fock/Bochner.lean}: 56\\
\path{GW.Fock.bochnerReconstruct_ae} & \path{GW/Fock/Bochner.lean}: 58\\
\path{GW.Fock.bochnerReconstruct_norm_sq} & \path{GW/Fock/Bochner.lean}: 61\\
\path{GW.Fock.bochnerCoordinate_reconstruct} & \path{GW/Fock/Bochner.lean}: 65\\
\path{GW.Fock.bochnerUnitary} & \path{GW/Fock/Bochner.lean}: 143\\
\path{GW.Fock.bochnerUnitary_single} & \path{GW/Fock/Bochner.lean}: 193\\
\multicolumn{2}{l}{\textbf{GW-Fields}}\\*
\path{GW.Fock.ordinaryCreationRootInput} & \path{GW/Fock/FockFieldDefinition.lean}: 85\\
\path{GW.Fock.ordinaryAnnihilationRootInput} & \path{GW/Fock/FockFieldDefinition.lean}: 90\\
\path{GW.Fock.ordinaryField} & \path{GW/Fock/FockFieldDefinition.lean}: 94\\
\path{GW.Fock.ordinaryField_domain} & \path{GW/Fock/FockFieldDefinition.lean}: 101\\
\path{GW.Fock.ordinaryField_action} & \path{GW/Fock/FockFieldDefinition.lean}: 104\\
\path{GW.Fock.ordinaryField_dense} & \path{GW/Fock/FockFieldDefinition.lean}: 112\\
\path{GW.Fock.ordinaryField_symmetric} & \path{GW/Fock/FockFieldDefinition.lean}: 116\\
\multicolumn{2}{l}{\textbf{GW-Bounds}}\\*
\path{GW.Fock.rootNumberCoefficient} & \path{GW/Fock/FockGeneralizedBounds.lean}: 14\\
\path{GW.Fock.creation_mem_shiftedRoot} & \path{GW/Fock/FockGeneralizedBounds.lean}: 39\\
\path{GW.Fock.annihilation_mem_root} & \path{GW/Fock/FockGeneralizedBounds.lean}: 49\\
\path{GW.Fock.creation_root_norm} & \path{GW/Fock/FockGeneralizedBounds.lean}: 60\\
\path{GW.Fock.annihilation_root_norm} & \path{GW/Fock/FockGeneralizedBounds.lean}: 80\\
\multicolumn{2}{l}{\textbf{GW-Adjoints}}\\*
\path{GW.Fock.creation_annihilation_pairing} & \path{GW/Fock/FockAdjoints.lean}: 183\\
\path{GW.Fock.ordinaryCreation_adjoint_eq} & \path{GW/Fock/FockAdjoints.lean}: 330\\
\multicolumn{2}{l}{\textbf{GW-CCR}}\\*
\path{GW.Fock.ordinaryCreation_domain_eq_annihilation} & \path{GW/Fock/FockCcr.lean}: 751\\
\path{GW.Fock.ordinaryCreation_norm_identity} & \path{GW/Fock/FockCcr.lean}: 755\\
\path{GW.Fock.ordinaryCCR_mixed_graphs} & \path{GW/Fock/FockCcr.lean}: 768\\
\path{GW.Fock.ordinaryCCR_annihilation} & \path{GW/Fock/FockCcr.lean}: 781\\
\path{GW.Fock.ordinaryCCR_creation} & \path{GW/Fock/FockCcr.lean}: 785\\
\multicolumn{2}{l}{\textbf{GW-Weighted}}\\*
\path{GW.Spaces.WeightedL2} & \path{GW/Spaces/WeightedL2.lean}: 18\\
\path{GW.Spaces.WeightedL2.coordinates} & \path{GW/Spaces/WeightedL2.lean}: 34\\
\path{GW.Spaces.WeightedL2.embed} & \path{GW/Spaces/WeightedL2.lean}: 45\\
\path{GW.Spaces.WeightedL2.embed_ae} & \path{GW/Spaces/WeightedL2.lean}: 51\\
\path{GW.Spaces.WeightedL2.embed_norm_le} & \path{GW/Spaces/WeightedL2.lean}: 56\\
\path{GW.Spaces.WeightedL2.weighted_embed_ae} & \path{GW/Spaces/WeightedL2.lean}: 67\\
\path{GW.Spaces.WeightedL2.embed_mem_domain} & \path{GW/Spaces/WeightedL2.lean}: 76\\
\path{GW.Spaces.WeightedL2.embed_injective} & \path{GW/Spaces/WeightedL2.lean}: 81\\
\path{GW.Spaces.WeightedL2.embed_range} & \path{GW/Spaces/WeightedL2.lean}: 90\\
\path{GW.Spaces.WeightedL2.weighted_norm_embed} & \path{GW/Spaces/WeightedL2.lean}: 109\\
\path{GW.Spaces.WeightedL2.embed_denseRange} & \path{GW/Spaces/WeightedL2.lean}: 190\\
\multicolumn{2}{l}{\textbf{GW-Multiplier}}\\*
\path{GW.Analysis.boundedMultiplier} & \path{GW/Analysis/L2Quotient.lean}: 154\\
\path{GW.Analysis.boundedMultiplier_ae} & \path{GW/Analysis/L2Quotient.lean}: 163\\
\path{GW.Analysis.multiplierDomain} & \path{GW/Analysis/L2Quotient.lean}: 33\\
\path{GW.Analysis.multiplierValue} & \path{GW/Analysis/L2Quotient.lean}: 63\\
\path{GW.Analysis.realMultiplier} & \path{GW/Analysis/RealMultiplier.lean}: 14\\
\path{GW.Analysis.realMultiplier_domain} & \path{GW/Analysis/RealMultiplier.lean}: 17\\
\path{GW.Analysis.realMultiplier_ae} & \path{GW/Analysis/RealMultiplier.lean}: 21\\
\path{GW.Analysis.realMultiplier_selfadjoint} & \path{GW/Analysis/RealMultiplier.lean}: 243\\
\path{GW.Analysis.realMultiplier_closed} & \path{GW/Analysis/RealMultiplier.lean}: 184\\
\path{GW.Analysis.realCutoff_tendsto} & \path{GW/Analysis/RealMultiplier.lean}: 159\\
\multicolumn{2}{l}{\textbf{GW-Model}}\\*
\path{GW.Model.v} & \path{GW/Model/Coefficient.lean}: 14\\
\path{GW.Model.cutoff} & \path{GW/Model/Coefficient.lean}: 40\\
\path{GW.Model.cutoff_memLp} & \path{GW/Model/ModelFrohlichEstimates.lean}: 56\\
\path{GW.Model.cutoffL2} & \path{GW/Model/ModelFrohlichEstimates.lean}: 65\\
\path{GW.Model.cutoffL2_ae} & \path{GW/Model/ModelFrohlichEstimates.lean}: 68\\
\path{GW.Model.cutoffL2_norm_sq} & \path{GW/Model/ModelFrohlichEstimates.lean}: 72\\
\path{GW.Model.unitBallVolume} & \path{GW/Model/ModelPowerIntegrals.lean}: 12\\
\path{GW.Model.unitBallVolume_pos} & \path{GW/Model/ModelPowerIntegrals.lean}: 14\\
\path{GW.Model.spaceBallVolume} & \path{GW/Model/ModelPowerIntegrals.lean}: 21\\
\path{GW.Model.spaceSphereVolume} & \path{GW/Model/ModelPowerIntegrals.lean}: 29\\
\path{GW.Model.spaceClosedBallVolume} & \path{GW/Model/ModelPowerIntegrals.lean}: 33\\
\path{GW.Model.exteriorPowerIntegral} & \path{GW/Model/ModelPowerIntegrals.lean}: 278\\
\path{GW.Model.interiorPowerIntegral} & \path{GW/Model/ModelPowerIntegrals.lean}: 279\\
\path{GW.Model.exteriorPowerIntegral_finite_iff} & \path{GW/Model/ModelPowerIntegrals.lean}: 344\\
\path{GW.Model.interiorPowerIntegral_finite_iff} & \path{GW/Model/ModelPowerIntegrals.lean}: 357\\
\path{GW.Model.exteriorPowerIntegral_four_bound} & \path{GW/Model/ModelPowerIntegrals.lean}: 383\\
\path{GW.Model.interiorPowerIntegral_two_bound} & \path{GW/Model/ModelPowerIntegrals.lean}: 406\\
\multicolumn{2}{l}{\textbf{GW-Form}}\\*
\path{GW.Operators.AlgebraicForm} & \path{GW/Operators/Polarization.lean}: 12\\
\path{GW.Operators.CoerciveFormData} & \path{GW/Operators/FormRepresentation.lean}: 75\\
\path{GW.Operators.energyForm} & \path{GW/Operators/FormRepresentation.lean}: 93\\
\path{GW.Operators.formInverse} & \path{GW/Operators/FormRepresentation.lean}: 147\\
\path{GW.Operators.formSolution} & \path{GW/Operators/FormRepresentation.lean}: 177\\
\path{GW.Operators.formSolution_equation} & \path{GW/Operators/FormRepresentation.lean}: 182\\
\path{GW.Operators.formSolution_unique} & \path{GW/Operators/FormRepresentation.lean}: 187\\
\path{GW.Operators.formOperator} & \path{GW/Operators/FormRepresentation.lean}: 253\\
\path{GW.Operators.formOperator_selfadjoint} & \path{GW/Operators/FormRepresentation.lean}: 258\\
\path{GW.Operators.operatorFormCoordinate} & \path{GW/Operators/FormRepresentation.lean}: 269\\
\path{GW.Operators.operatorFormCoordinate_embedding} & \path{GW/Operators/FormRepresentation.lean}: 272\\
\path{GW.Operators.formOperator_form_identity} & \path{GW/Operators/FormRepresentation.lean}: 281\\
\path{GW.Operators.formOperator_domain_of_equation} & \path{GW/Operators/FormRepresentation.lean}: 290\\
\path{GW.Operators.formOperator_domain_iff} & \path{GW/Operators/FormRepresentation.lean}: 308\\
\path{GW.Operators.formOperator_unique} & \path{GW/Operators/FormRepresentation.lean}: 328\\
\path{GW.Operators.formOperator_inverse_left} & \path{GW/Operators/FormRepresentation.lean}: 340\\
\path{GW.Operators.formOperator_inverse_right} & \path{GW/Operators/FormRepresentation.lean}: 345\\
\multicolumn{2}{l}{\textbf{GW-Coercive}}\\*
\path{GW.Operators.coercive_bijective} & \path{GW/Operators/CoerciveBounded.lean}: 33\\
\path{GW.Operators.coerciveInverse} & \path{GW/Operators/CoerciveBounded.lean}: 47\\
\path{GW.Operators.coerciveInverse_left} & \path{GW/Operators/CoerciveBounded.lean}: 56\\
\path{GW.Operators.coerciveInverse_right} & \path{GW/Operators/CoerciveBounded.lean}: 62\\
\path{GW.Operators.coerciveInverse_norm_le} & \path{GW/Operators/CoerciveBounded.lean}: 68\\
\multicolumn{2}{l}{\textbf{GW-Resolvent}}\\*
\path{GW.Operators.IsResolvent} & \path{GW/Operators/ResolventAlgebra.lean}: 12\\
\path{GW.Operators.isResolvent_unique} & \path{GW/Operators/ResolventAlgebra.lean}: 35\\
\path{GW.Operators.resolvent_isResolvent} & \path{GW/Operators/ResolventAlgebra.lean}: 27\\
\multicolumn{2}{l}{\textbf{GW-Mixed}}\\*
\path{GW.Analysis.BoundedForm} & \path{GW/Analysis/MixedForms.lean}: 16\\
\path{GW.Analysis.form_bound} & \path{GW/Analysis/MixedForms.lean}: 19\\
\path{GW.Analysis.form_norm_le_iff} & \path{GW/Analysis/MixedForms.lean}: 25\\
\path{GW.Analysis.mixed_form_estimate} & \path{GW/Analysis/MixedForms.lean}: 36\\
\path{GW.Analysis.form_dense_determination} & \path{GW/Analysis/MixedForms.lean}: 49\\
\path{GW.Analysis.mixed_form_tendsto} & \path{GW/Analysis/MixedForms.lean}: 60\\
\multicolumn{2}{l}{\textbf{GW-Domination}}\\*
\path{GW.Analysis.l2_limit_measurable} & \path{GW/Analysis/L2Domination.lean}: 17\\
\path{GW.Analysis.dominated_l2_membership} & \path{GW/Analysis/L2Domination.lean}: 28\\
\path{GW.Analysis.dominated_l2_convergence} & \path{GW/Analysis/L2Domination.lean}: 48\\
\bottomrule\end{longtable}\normalsize

\subsection{Exactly what remains genuinely new}
These are \emph{not} further GW copies: the fiber field momentum multiplier $P_f$, the free operator $(P-P_f)^2+N$ and its momentum-support core, fiber-uniform Lieb--Yamazaki and singular form limits, the integrated $L^2(k;\F)$ annihilation field and energy-valued cutoff pull-through, the vacuum and one-phonon Schur complements, the resolvent identities used by Seiringer, the exact infrared saturation identity, and Mitrouskas's kernel/weighted integral estimates. The full proofs in this revision and their explicit adapter nodes are the mathematical specification for these modules; the exact source declarations above fix their implementation foundation.

\begin{thebibliography}{99}
\bibitem{mitrouskas}
D. Mitrouskas,
\emph{Absence of Excited Bound States for All Fibers of the Fr\"ohlich Polaron at Weak Coupling},
arXiv:2611.09586, version 1 (2026).
\url{https://arxiv.org/abs/2611.09586}.
Equation and section references in this document refer to this version.

\bibitem{seiringer}
R. Seiringer,
\emph{Absence of excited eigenvalues for Fr\"ohlich type polaron models at weak coupling},
Journal of Spectral Theory \textbf{13} (2023), 1045--1055.
DOI: 10.4171/JST/469.
Preprint arXiv:2210.17123; the proof was checked against its October 2022 version.
\url{https://arxiv.org/abs/2210.17123};
\url{https://ems.press/journals/jst/articles/13394944}.

\bibitem{gw}
M. Griesemer and A. W\"unsch,
\emph{Self-adjointness and domain of the Fr\"ohlich Hamiltonian},
Journal of Mathematical Physics \textbf{57} (2016), 021902.
Preprint arXiv:1508.02533, version 2. \url{https://arxiv.org/abs/1508.02533}. The paper's Lemma 2.1 and Theorem 2.2 motivate the selected form-limit/representation foundation, not a direct inference about fibers. The user-supplied Lean release is source-audited in Section~\ref{sec:gwverify}; no fiber theorem is imported from the full-space GW formalization.

\bibitem{mathlib-dual}
mathlib documentation,
\emph{Mathlib.Analysis.InnerProductSpace.Dual}.
\url{https://leanprover-community.github.io/mathlib4_docs/Mathlib/Analysis/InnerProductSpace/Dual.html}.

\bibitem{mathlib-inverse}
mathlib documentation,
\emph{Mathlib.Analysis.Calculus.FDeriv.Mul}.
\url{https://leanprover-community.github.io/mathlib4_docs/Mathlib/Analysis/Calculus/FDeriv/Mul.html}.

\bibitem{mathlib-l2}
mathlib documentation,
\emph{Mathlib.MeasureTheory.Function.L2Space}.
\url{https://leanprover-community.github.io/mathlib4_docs/Mathlib/MeasureTheory/Function/L2Space.html}.

\bibitem{mathlib-fubini}
mathlib documentation,
\emph{Mathlib.MeasureTheory.Integral.Prod}.
\url{https://leanprover-community.github.io/mathlib4_docs/Mathlib/MeasureTheory/Integral/Prod.html}.

\bibitem{mathlib-laxmilgram}
mathlib documentation,
\emph{Mathlib.Analysis.InnerProductSpace.LaxMilgram}.
\url{https://leanprover-community.github.io/mathlib4_docs/Mathlib/Analysis/InnerProductSpace/LaxMilgram.html}.

\bibitem{mathlib-spectrum}
mathlib documentation,
\emph{Mathlib.Analysis.CStarAlgebra.Spectrum}.
\url{https://leanprover-community.github.io/mathlib4_docs/Mathlib/Analysis/CStarAlgebra/Spectrum.html}.
\end{thebibliography}
\end{document}
